% Produire le plan détaillé d’une dissertation — Français Tle Tech — Cours signé OnBuch+
% Programme officiel MINESEC. Compiler avec : tectonic N19-produire-plan-detaille-dissertation.tex
\newcommand{\DOCMATIERE}{Français}
\newcommand{\DOCNIVEAU}{Tle Tech}
\newcommand{\DOCMODULE}{Français – classe de Terminale technique}
\newcommand{\DOCLECON}{N19}
\newcommand{\DOCTITRE}{Produire le plan détaillé d’une dissertation}
% OnBuch+ — préambule « cours signés » ENRICHI (compilé avec tectonic / XeLaTeX)
% Système visuel « Pop » d'OnBuch+ : fond crème (jamais blanc), bordures encre
% épaisses, ombres « dures » décalées, titres Archivo Black en capitales.
% Version MATHÉMATIQUES : graphes (pgfplots/tikz), boîtes pédagogiques
% (définition, propriété, méthode, attention, le savais-tu, exercice résolu,
% exercice, TP, bilan), page de garde et signature OnBuch+ ancrée en bas.
%
% Macros attendues dans main.tex AVANT \input{preamble} :
%   \DOCMATIERE  \DOCNIVEAU  \DOCMODULE  \DOCLECON (numéro)  \DOCTITRE
\documentclass[11pt]{article}

\usepackage{fontspec}
\usepackage[french]{babel}
\usepackage[a4paper,margin=2cm,top=3.4cm,bottom=2.8cm,headheight=1.4cm,footskip=1.3cm]{geometry}
\usepackage{amsmath,amssymb,amsfonts}
\usepackage{mathtools} % \xmapsto, \coloneqq... souvent utilisés par les modèles
\usepackage{cancel} % simplifications barrées (bilans de forces)
% \abs{z} (module, valeur absolue) et \norm{u} (norme), utilisés par les modèles
\providecommand{\abs}{}\let\abs\relax\DeclarePairedDelimiter{\abs}{\lvert}{\rvert}
\providecommand{\norm}{}\let\norm\relax\DeclarePairedDelimiter{\norm}{\lVert}{\rVert}
\usepackage[table]{xcolor} % [table] : fournit aussi \rowcolors (alternance de lignes)
\usepackage{pagecolor}
\usepackage{tikz}
\usetikzlibrary{arrows.meta,positioning,calc,shapes.geometric,shapes.misc,shapes.arrows,decorations.pathmorphing,decorations.pathreplacing,decorations.markings,patterns,shadows,mindmap,trees,fit,backgrounds,matrix,intersections,angles,quotes}
\usepackage{pgfplots}
\usepackage[version=4]{mhchem} % équations nucléaires \ce{^{235}_{92}U}
\usepackage[european,siunitx]{circuitikz} % schémas électriques
\pgfplotsset{compat=1.18}
\usepgfplotslibrary{fillbetween,groupplots} % aires entre courbes (calcul intégral)
\usepackage{enumitem}
\usepackage{fancyhdr}
% [most] charge toutes les bibliothèques tcolorbox (listings, minted, raster,
% documentation...) et gonfle inutilement la mémoire TeX sur un document long
% (~300 boîtes) : on ne charge que ce qui est réellement utilisé.
\usepackage[skins,breakable]{tcolorbox}
\usepackage{graphicx}
\usepackage{colortbl}
\usepackage{booktabs}
\usepackage{tabularx}
\usepackage{multirow}
\usepackage{diagbox} % case d'en-tête diagonale des tableaux à double entrée
% Colonnes extensibles centrée / gauche / droite (C, L, R), souvent utilisées par les modèles
\newcolumntype{C}{>{\centering\arraybackslash}X}
\newcolumntype{L}{>{\raggedright\arraybackslash}X}
\newcolumntype{R}{>{\raggedleft\arraybackslash}X}
\usepackage{array}
\usepackage{multicol}
\usepackage{siunitx}
% parse-numbers=false : accepte aussi \SI{2,5 \times 10^{-3}}{...}
\sisetup{locale=FR,output-decimal-marker={,},group-minimum-digits=4,parse-numbers=false}
\DeclareSIUnit\ohm{\ensuremath{\symup{\Omega}}} % U+2126 absent de la police
\DeclareSIUnit\division{div} % réglages d'oscilloscope (ms/div, V/div)
\DeclareSIUnit\tr{tr} % tours (vitesse de rotation, ex. \SI{1500}{\tr\per\minute})
\DeclareSIUnit\molar{M} % concentration molaire (ex. \SI{3}{\molar}, \SI{1}{\micro\molar})
\DeclareSIUnit\an{an} % année (le modèle confond parfois avec \year, un compteur TeX) — ex. \SI{5}{\kilo\meter\per\an}
\DeclareSIUnit\FCFA{FCFA} % franc CFA (ex. \SI{500}{\FCFA\per\kilo\gram})
\DeclareSIUnit\degreeBrix{\textdegree Brix} % teneur en sucre d'un sirop/jus (réfractométrie)
\DeclareSIUnit\month{mois} % le modèle utilise parfois \month, un compteur TeX, comme unité de durée
\DeclareSIUnit\microliter{\micro\liter} % le modèle écrit parfois \microliter en un seul token
\DeclareSIUnit\million{millions} % dénombrement (ex. \SI{15}{\million\per\mL} spermatozoïdes)
\usepackage{qrcode}
% \degree : commande fréquemment utilisée par les modèles pour un angle ou une
% température en dehors de \SI{}{} (non fournie par les paquets chargés).
\providecommand{\degree}{^{\circ}}
% \qed (fin de démonstration), utilisé par les modèles sans amsthm
\providecommand{\qed}{\hfill$\square$}
% \barre : double barre des valeurs interdites dans un tableau de variations (tkz-tab)
\providecommand{\barre}{\ensuremath{\|}}
% environnement proof (amsthm non chargé) utilisé par les modèles
\ifdefined\proof\else\newenvironment{proof}[1][Démonstration]{\par\noindent\textit{#1.}\ }{\qed\par}\fi
\usepackage{lastpage}
\usepackage{hyperref}
\hypersetup{hidelinks}

% ============================================================
% Palette « Pop » OnBuch+ — mêmes hex que la classe P (plus_ui.dart)
% ============================================================
\definecolor{poporange}{HTML}{F59321}
\definecolor{popdark}{HTML}{E07A0C}
\definecolor{popcream}{HTML}{FFF6E8}
\definecolor{popink}{HTML}{1C1714}
\definecolor{poppurple}{HTML}{7A5AE0}
\definecolor{popgreen}{HTML}{2E9E6A}
\definecolor{poppink}{HTML}{E0457B}
\definecolor{popblue}{HTML}{2F7DE1}
\definecolor{popgold}{HTML}{C98A12}
\definecolor{popmuted}{HTML}{8A7F76}
\definecolor{popline}{HTML}{EADFD2}
\definecolor{popgray}{HTML}{8A7F76}
\colorlet{popgrey}{popgray}
% Alias tolérants (le modèle nomme parfois les couleurs différemment)
\colorlet{popwhite}{white}
\colorlet{popblack}{popink}
\colorlet{popred}{poppink}
\colorlet{popyellow}{popgold}
% Teintes claires (fonds de tableaux/encadrés) — le modèle nomme parfois
% les couleurs avec un suffixe L (ex. popblueL) sans les avoir définies.
\colorlet{poporangeL}{poporange!20}
\colorlet{popdarkL}{popdark!20}
\colorlet{poppurpleL}{poppurple!20}
\colorlet{popgreenL}{popgreen!20}
\colorlet{poppinkL}{poppink!20}
\colorlet{popblueL}{popblue!20}
\colorlet{popgoldL}{popgold!20}
\colorlet{popredL}{popred!20}
\colorlet{popgrayL}{popgray!20}
% Teintes claires pour fonds de tableaux / zones
\colorlet{popblueL}{popblue!10!white}
\colorlet{poporangeL}{poporange!14!white}
\colorlet{popgreenL}{popgreen!12!white}
\colorlet{poppurpleL}{poppurple!12!white}
\colorlet{poppinkL}{poppink!12!white}
\colorlet{popgoldL}{popgold!14!white}

\pagecolor{popcream}

% ============================================================
% Typographie
% ============================================================
\defaultfontfeatures{Ligatures=TeX}
\setmainfont{PlusJakartaSans-Medium.ttf}[Path=../../../fonts/,BoldFont=PlusJakartaSans-Bold.ttf]
% Maths en sans-serif (Fira Math) avec chiffres et lettres droites de Plus
% Jakarta, pour que les formules se fondent dans le texte.
\usepackage[math-style=ISO,bold-style=ISO]{unicode-math}
\setmathfont{FiraMath-Regular.otf}[Path=../../../fonts/]
\setmathfont{PlusJakartaSans-Medium.ttf}[Path=../../../fonts/,range={up/num,up/{Latin,latin},\mathup}]
\usepackage{icomma} % virgule décimale sans espace en mode math
% Caractères mathématiques Unicode absents des polices de texte (Plus Jakarta,
% Archivo Black) : rendus par leur équivalent mathématique, en texte comme en
% formule — sinon ils s'affichent en carré vide (ex. « Arithmétique dans ℤ »).
\usepackage{newunicodechar}
\newunicodechar{ℝ}{\ensuremath{\mathbb{R}}}
\newunicodechar{ℤ}{\ensuremath{\mathbb{Z}}}
\newunicodechar{ℂ}{\ensuremath{\mathbb{C}}}
\newunicodechar{ℕ}{\ensuremath{\mathbb{N}}}
\newunicodechar{ℚ}{\ensuremath{\mathbb{Q}}}
\newunicodechar{𝔻}{\ensuremath{\mathbb{D}}}
\newunicodechar{α}{\ensuremath{\alpha}}
\newunicodechar{β}{\ensuremath{\beta}}
\newunicodechar{γ}{\ensuremath{\gamma}}
\newunicodechar{δ}{\ensuremath{\delta}}
\newunicodechar{ε}{\ensuremath{\varepsilon}}
\newunicodechar{θ}{\ensuremath{\theta}}
\newunicodechar{λ}{\ensuremath{\lambda}}
\newunicodechar{μ}{\ensuremath{\mu}}
\newunicodechar{σ}{\ensuremath{\sigma}}
\newunicodechar{τ}{\ensuremath{\tau}}
\newunicodechar{φ}{\ensuremath{\varphi}}
\newunicodechar{ψ}{\ensuremath{\psi}}
\newunicodechar{ω}{\ensuremath{\omega}}
\newunicodechar{Σ}{\ensuremath{\Sigma}}
% majuscules produites par \MakeUppercase dans les titres (α-aminés -> Α)
\newunicodechar{Α}{\ensuremath{\alpha}}
\newunicodechar{Β}{\ensuremath{\beta}}
\newunicodechar{Γ}{\ensuremath{\Gamma}}
\newunicodechar{Δ}{\ensuremath{\Delta}}
\newunicodechar{Ε}{\ensuremath{\varepsilon}}
\newunicodechar{Θ}{\ensuremath{\Theta}}
\newunicodechar{Λ}{\ensuremath{\Lambda}}
\newunicodechar{Μ}{\ensuremath{\mu}}
\newunicodechar{Τ}{\ensuremath{\tau}}
\newunicodechar{Φ}{\ensuremath{\Phi}}
\newunicodechar{Ψ}{\ensuremath{\Psi}}
\newunicodechar{Ω}{\ensuremath{\Omega}}
\newunicodechar{↦}{\ensuremath{\mapsto}}
\newunicodechar{∈}{\ensuremath{\in}}
\newunicodechar{∉}{\ensuremath{\notin}}
\newunicodechar{∧}{\ensuremath{\wedge}}
\newunicodechar{⇔}{\ensuremath{\Leftrightarrow}}
\newunicodechar{⟺}{\ensuremath{\Longleftrightarrow}}
\newunicodechar{⇒}{\ensuremath{\Rightarrow}}
\newunicodechar{⟹}{\ensuremath{\Longrightarrow}}
\newunicodechar{∘}{\ensuremath{\circ}}
\newunicodechar{∪}{\ensuremath{\cup}}
\newunicodechar{∩}{\ensuremath{\cap}}
\newunicodechar{≡}{\ensuremath{\equiv}}
\newunicodechar{⊂}{\ensuremath{\subset}}
\newunicodechar{⊥}{\ensuremath{\perp}}
\newunicodechar{∃}{\ensuremath{\exists}}
\newunicodechar{∀}{\ensuremath{\forall}}
\newunicodechar{∛}{\ensuremath{\sqrt[3]{\,}}}
\newunicodechar{∜}{\ensuremath{\sqrt[4]{\,}}}
\newunicodechar{⃗}{\ensuremath{{}^{\to}}}
\newunicodechar{ⁿ}{\ifmmode^{n}\else\textsuperscript{\ensuremath{n}}\fi}
\newunicodechar{ˣ}{\ifmmode^{x}\else\textsuperscript{\ensuremath{x}}\fi}
\newunicodechar{ᵇ}{\ifmmode^{b}\else\textsuperscript{\ensuremath{b}}\fi}
\newunicodechar{ʳ}{\ifmmode^{r}\else\textsuperscript{\ensuremath{r}}\fi}
\newunicodechar{ᵘ}{\ifmmode^{u}\else\textsuperscript{\ensuremath{u}}\fi}
\newunicodechar{ʸ}{\ifmmode^{y}\else\textsuperscript{\ensuremath{y}}\fi}
\newunicodechar{ᵃ}{\ifmmode^{a}\else\textsuperscript{\ensuremath{a}}\fi}
\newunicodechar{ᵉ}{\ifmmode^{e}\else\textsuperscript{\ensuremath{e}}\fi}
\newunicodechar{ᵏ}{\ifmmode^{k}\else\textsuperscript{\ensuremath{k}}\fi}
\newunicodechar{ᵖ}{\ifmmode^{p}\else\textsuperscript{\ensuremath{p}}\fi}
\newunicodechar{ᴺ}{\ifmmode^{N}\else\textsuperscript{\ensuremath{N}}\fi}
\newunicodechar{ᶜ}{\ifmmode^{c}\else\textsuperscript{\ensuremath{c}}\fi}
\newunicodechar{ʰ}{\ifmmode^{h}\else\textsuperscript{\ensuremath{h}}\fi}
\newunicodechar{ᵠ}{\ifmmode^{q}\else\textsuperscript{\ensuremath{q}}\fi}
\newunicodechar{ₙ}{\ifmmode_{n}\else\textsubscript{\ensuremath{n}}\fi}
\newunicodechar{ₐ}{\ifmmode_{a}\else\textsubscript{\ensuremath{a}}\fi}
\newunicodechar{ᵢ}{\ifmmode_{i}\else\textsubscript{\ensuremath{i}}\fi}
\newunicodechar{ᵣ}{\ifmmode_{r}\else\textsubscript{\ensuremath{r}}\fi}
\newunicodechar{ₖ}{\ifmmode_{k}\else\textsubscript{\ensuremath{k}}\fi}
\newunicodechar{ᵦ}{\ifmmode_{\beta}\else\textsubscript{\ensuremath{\beta}}\fi}
\newunicodechar{ₓ}{\ifmmode_{x}\else\textsubscript{\ensuremath{x}}\fi}
\newfontfamily\popdisplay{ArchivoBlack-Regular.ttf}[Path=../../../fonts/]
\newfontfamily\poplogo{SpaceGrotesk-Bold.ttf}[Path=../../../fonts/]
\newfontfamily\popsemi{PlusJakartaSans-SemiBold.ttf}[Path=../../../fonts/]
\newfontfamily\popxbold{PlusJakartaSans-ExtraBold.ttf}[Path=../../../fonts/]
\newfontfamily\popxboldls{PlusJakartaSans-ExtraBold.ttf}[Path=../../../fonts/,LetterSpace=3.0]

\setlist[enumerate]{leftmargin=1.7em,itemsep=5pt,topsep=4pt}
\setlist[itemize]{leftmargin=1.7em,itemsep=4pt,topsep=4pt,label={\color{poporange}\textbullet}}
\setlength{\parindent}{0pt}
\setlength{\parskip}{5pt}
\renewcommand{\arraystretch}{1.25}

% pgfplots : style Pop par défaut
\pgfplotsset{
  every axis/.append style={axis line style={popink,line width=1pt},tick style={popink},
    grid style={popline,line width=0.5pt},label style={font=\small},tick label style={font=\footnotesize}},
  popcurve/.style={poporange,line width=1.8pt,no markers,smooth},
}
% Même style disponible dans un \draw TikZ simple (hors pgfplots)
\tikzset{popcurve/.style={poporange,line width=1.8pt}}
\tikzset{popgrid/.style={popline,very thin}}
\colorlet{popcurve}{poporange} % le nom du style est parfois utilisé comme couleur

% ============================================================
% Pill / squiggle
% ============================================================
\newcommand{\poppill}[3]{%
  \tikz[baseline=-0.55ex]\node[draw=popink,line width=1.2pt,rounded corners=2.6pt,%
    fill=#2,inner xsep=6.5pt,inner ysep=3pt]{%
    \popxboldls\fontsize{7}{7}\selectfont\color{#3}\MakeUppercase{#1}};%
}
\newcommand{\popsquiggle}[1]{%
  \tikz[baseline=-0.4ex]{
    \draw[#1,line width=2.6pt,line cap=round]
      plot [smooth] coordinates {(0,0.09) (0.34,0.01) (0.68,0.21) (1.02,0.01) (1.36,0.10)};
  }%
}
% Mot-clé surligné (marqueur orange sous le texte)
\newcommand{\cle}[1]{{\popxbold\color{popdark}#1}}

% ============================================================
% En-tête / pied
% ============================================================
\pagestyle{fancy}
\fancyhf{}
\renewcommand{\headrulewidth}{0pt}
\renewcommand{\footrulewidth}{0pt}
\lhead{\raisebox{-0.15ex}{\poplogo\fontsize{13}{13}\selectfont\color{popink}OnBuch\color{poporange}+}}
\rhead{\poppill{\DOCMATIERE\ \textbullet\ \DOCNIVEAU\ \textbullet\ Leçon \DOCLECON}{white}{popink}}
\lfoot{\popxbold\fontsize{7.6}{7.6}\selectfont\color{popmuted}\MakeUppercase{Cours signé OnBuch+}\ \textbullet\ {\popsemi\DOCTITRE}}
\rfoot{\tikz[baseline=-0.6ex]\node[rounded corners=8.5pt,fill=popink,minimum height=17pt,minimum width=17pt,inner xsep=5pt,inner sep=0]{\popxbold\fontsize{8}{8}\selectfont\color{popcream}\thepage\,/\,\pageref*{LastPage}};}

% ============================================================
% Titres
% ============================================================
\newcommand{\chaptitre}[2]{%
  \begin{tcolorbox}[enhanced,colback=poporange,colframe=popink,boxrule=2.6pt,arc=11pt,
      drop shadow=popink,left=15pt,right=15pt,top=12pt,bottom=11pt,before skip=3pt,after skip=18pt]
    \poppill{#2}{popink}{popcream}\par\vspace{9pt}
    {\raggedright\hyphenpenalty=10000\exhyphenpenalty=10000\popdisplay\fontsize{22}{24}\selectfont\color{popcream}\MakeUppercase{#1}\par}
  \end{tcolorbox}
}
% Section numérotée : pastille orange avec le numéro + titre Archivo
\newcounter{popsec}
\newcommand{\coursec}[1]{%
  \stepcounter{popsec}\setcounter{popsub}{0}%
  \par\vspace{16pt}\noindent
  \begin{minipage}{\linewidth}
  \tikz[baseline=-0.7ex]\node[draw=popink,line width=1.6pt,fill=poporange,rounded corners=4pt,minimum size=22pt,inner sep=0]{\popdisplay\fontsize{12}{12}\selectfont\color{popcream}\thepopsec};%
  \hspace{8pt}{\popdisplay\fontsize{15}{17}\selectfont\color{popink}\MakeUppercase{#1}}\par
  \vspace{4pt}\noindent{\color{poporange}\rule{36pt}{3.6pt}}
  \end{minipage}\par\nopagebreak\vspace{8pt}
}
\newcounter{popsub}
\newcommand{\courssub}[1]{%
  \stepcounter{popsub}%
  \par\vspace{9pt}\noindent{\popxbold\fontsize{11.5}{13}\selectfont\color{popink}{\color{poporange}\thepopsec.\thepopsub}\hspace{6pt}#1}\par\nopagebreak\vspace{4pt}
}

% ============================================================
% Boîtes pédagogiques — même recette : fond blanc, bordure épaisse colorée,
% ombre dure encre, badge plein. Titre optionnel [#1].
% ============================================================
\tcbset{popbase/.style={breakable,enhanced,colback=white,boxrule=2.3pt,arc=9pt,
  shadow={3pt}{-3pt}{0pt}{popink},left=14pt,right=14pt,top=11pt,bottom=11pt,before skip=11pt,after skip=15pt,
  fontupper=\popsemi\fontsize{10.6}{14.5}\selectfont\color{popink},
  fontlower=\popsemi\fontsize{10.6}{14.5}\selectfont\color{popink}}}
% Coche dessinée (le glyphe ✓ n'existe pas dans les polices du document)
\newcommand{\popcheck}[1]{\tikz[baseline=-0.5ex]\draw[#1,line width=1.6pt,line cap=round,line join=round](-0.13,0.0)--(-0.04,-0.09)--(0.13,0.1);}
\providecommand{\coche}{\popcheck{popgreen}}
\providecommand{\resultat}[1]{\par\smallskip\noindent\fcolorbox{popgreen}{popgreenL}{\parbox{\dimexpr\linewidth-2\fboxsep-2\fboxrule\relax}{\textbf{Résultat.} #1}}\par\smallskip}
\newcommand{\popboxhead}[3]{\poppill{#1}{#2}{white}\ifx\relax#3\relax\else\hspace{6pt}{\popxbold\fontsize{10.6}{12}\selectfont\color{popink}#3}\fi\par\vspace{7pt}}

\newtcolorbox{aretenir}[1][]{popbase,colframe=popgreen,before upper={\popboxhead{À retenir}{popgreen}{#1}}}
\newtcolorbox{exemplebox}[1][]{popbase,colframe=poporange,before upper={\popboxhead{Exemple}{poporange}{#1}}}
\newtcolorbox{definition}[1][]{popbase,colframe=popblue,before upper={\popboxhead{Définition}{popblue}{#1}}}
\newtcolorbox{propriete}[1][]{popbase,colframe=poppurple,before upper={\popboxhead{Propriété}{poppurple}{#1}}}
\newtcolorbox{methode}[1][]{popbase,colframe=popgold,before upper={\popboxhead{Méthode}{popgold}{#1}}}
\newtcolorbox{attention}[1][]{popbase,colframe=poppink,before upper={\popboxhead{Attention}{poppink}{#1}}}
\newtcolorbox{experience}[1][]{popbase,colframe=popblue,colback=popblueL,before upper={\popboxhead{Expérience}{popblue}{#1}}}
% « Le savais-tu ? » : carte violette pleine (contexte, culture, Cameroun)
\newtcolorbox{savaistu}[1][]{popbase,colback=poppurple,colframe=popink,
  fontupper=\popsemi\fontsize{10.4}{14.2}\selectfont\color{white},
  before upper={\poppill{Le savais-tu\,?}{white}{poppurple}\ifx\relax#1\relax\else\hspace{6pt}{\popxbold\color{white}#1}\fi\par\vspace{7pt}}}
% Exercice résolu : énoncé en haut, \tcblower puis la solution sur fond crème
\newcounter{popexo}\newcounter{popexr}
\newtcolorbox{exoresolu}[1][]{popbase,colframe=poporange,colbacklower=popcream,
  segmentation style={popink,line width=1.2pt,dashed},
  before upper={\stepcounter{popexr}\popboxhead{Exercice résolu \thepopexr}{poporange}{#1}},
  before lower={\poppill{Solution}{popink}{popcream}\par\vspace{6pt}}}
% Exercice (non résolu en place) — la correction est dans la partie Corrigés
\newtcolorbox{exercice}[1][]{popbase,colframe=popink,
  before upper={\stepcounter{popexo}\popboxhead{Exercice \thepopexo}{popink}{#1}}}
% Corrigé
\newtcolorbox{corrige}[1][]{popbase,colframe=popgreen,colback=popgreenL,
  before upper={\popboxhead{Corrigé}{popgreen}{#1}}}
% Objectifs / prérequis / bilan
\newtcolorbox{objectifs}{popbase,colframe=popink,colback=poporangeL,
  before upper={\popboxhead{Objectifs de la leçon}{popink}{}}}
\newtcolorbox{prerequis}{popbase,colframe=popink,colback=popblueL,
  before upper={\popboxhead{Prérequis}{popblue}{}}}
\newtcolorbox{bilan}{popbase,colframe=popink,colback=white,boxrule=2.8pt,
  before upper={\popboxhead{Fiche bilan}{poporange}{L'essentiel à connaître pour l'examen}}}
% Figure encadrée avec légende
\newtcolorbox{popfigure}[1][]{enhanced,breakable=false,colback=white,colframe=popline,boxrule=1.4pt,arc=8pt,
  left=10pt,right=10pt,top=10pt,bottom=8pt,before skip=10pt,after skip=12pt,halign=center,
  lower separated=false,halign lower=center,
  fontlower=\popsemi\fontsize{9}{11.5}\selectfont\color{popmuted},
  #1}
\newcommand{\legende}[1]{\par\vspace{6pt}{\popsemi\fontsize{9}{11.5}\selectfont\color{popmuted}#1}}

% ============================================================
% Signature OnBuch+
% ============================================================
\newcommand{\poplogoinline}[1]{{\poplogo\fontsize{#1}{#1}\selectfont\color{popink}OnBuch\color{poporange}+}}
\newcommand{\signatureonbuch}{%
  \begin{tcolorbox}[enhanced,colback=popink,colframe=popink,boxrule=2.2pt,arc=10pt,
      drop shadow=poporange,left=14pt,right=14pt,top=10pt,bottom=10pt,before skip=0pt,after skip=0pt]
    \tikz[baseline=-0.7ex]\node[circle,fill=popgreen,draw=popcream,line width=1.3pt,minimum size=22pt,inner sep=0]{\popcheck{white}};%
    \hspace{9pt}\begin{minipage}[c]{0.62\linewidth}
      {\popxbold\fontsize{12}{13}\selectfont\color{popcream}Signé\ \textbullet\ {\poplogo OnBuch\color{poporange}+}}\par\vspace{2pt}
      {\raggedright\popsemi\fontsize{8}{10.5}\selectfont\color{popline}\MakeUppercase{Équipe pédagogique OnBuch+}\\\MakeUppercase{Conforme au programme officiel MINESEC}\par}
    \end{minipage}\hfill
    {\poplogo\fontsize{22}{22}\selectfont\color{popcream}OnBuch\color{poporange}+}
  \end{tcolorbox}%
}

% ============================================================
% Page de garde
% #1 titre · #2 matière · #3 niveau · #4 module · #5 n° leçon · #6 durée ·
% #7 description / objectifs (texte ou itemize)
% ============================================================
\newcommand{\pagegarde}[7]{%
  \thispagestyle{empty}
  \newgeometry{a4paper,margin=1.7cm,top=1.6cm,bottom=1.4cm}%
  \begin{tikzpicture}[remember picture,overlay]
    \fill[poporange!18] ([xshift=-3.2cm,yshift=-3.6cm]current page.north east) circle (4.6cm);
    \fill[poppurple!14] ([xshift=2.4cm,yshift=7.6cm]current page.south west) circle (3.8cm);
  \end{tikzpicture}%
  {\poplogo\fontsize{30}{30}\selectfont\color{popink}OnBuch\color{poporange}+}\hfill
  \raisebox{0.6ex}{\poppill{Cours signé \textbullet\ Premium}{popink}{popcream}}\par
  \vspace{1pt}\popsquiggle{poppurple}\par
  \vspace{8pt}
  \begin{tcolorbox}[enhanced,colback=poporange,colframe=popink,boxrule=2.8pt,arc=13pt,
      drop shadow=popink,left=18pt,right=18pt,top=12pt,bottom=13pt,after skip=10pt]
    \poppill{#2 \textbullet\ #3}{popink}{popcream}\hspace{5pt}\poppill{#4}{white}{popink}\par\vspace{8pt}
    {\popdisplay\fontsize{14}{14}\selectfont\color{popink}LEÇON #5}\par\vspace{4pt}
    {\raggedright\hyphenpenalty=10000\exhyphenpenalty=10000\popdisplay\fontsize{25}{27}\selectfont\color{popcream}\MakeUppercase{#1}\par}
    \vspace{8pt}
    {\popxbold\fontsize{9.5}{9.5}\selectfont\color{popink}\MakeUppercase{Durée conseillée : #6}}
  \end{tcolorbox}
  \begin{tcolorbox}[enhanced,colback=white,colframe=popink,boxrule=2.2pt,arc=10pt,
      drop shadow=popink,left=16pt,right=16pt,top=10pt,bottom=10pt,after skip=8pt]
    \poppill{Dans cette leçon}{poppurple}{white}\par\vspace{5pt}
    \popsemi\fontsize{9.3}{12.6}\selectfont\color{popink}#7
  \end{tcolorbox}
  \noindent\begin{minipage}[t]{0.487\linewidth}
    \begin{tcolorbox}[enhanced,colback=popgreen,colframe=popink,boxrule=2.2pt,arc=10pt,
        drop shadow=popink,left=11pt,right=11pt,top=10pt,bottom=10pt,height=3.1cm,valign=center]
      \tcbox[colback=white,colframe=popink,boxrule=1.3pt,arc=5pt,boxsep=0pt,nobeforeafter,
        left=3pt,right=3pt,top=3pt,bottom=3pt,valign=center]{\qrcode[height=1.9cm]{https://play.google.com/store/apps/details?id=com.luvvix.onbuch&hl=fr}}%
      \hspace{8pt}\begin{minipage}[c]{0.5\linewidth}
        {\popdisplay\fontsize{11}{12.5}\selectfont\color{white}\MakeUppercase{Télécharge OnBuch}}\par\vspace{4pt}
        {\raggedright\popsemi\fontsize{8.4}{11}\selectfont\color{white}Gratuit \textbullet\ Google Play\\Tuteur IA, annales, concours, résultats\par}
      \end{minipage}
    \end{tcolorbox}
  \end{minipage}\hfill
  \begin{minipage}[t]{0.487\linewidth}
    \begin{tcolorbox}[enhanced,colback=poppurple,colframe=popink,boxrule=2.2pt,arc=10pt,
        drop shadow=popink,left=11pt,right=11pt,top=10pt,bottom=10pt,height=3.1cm,valign=center]
      \tikz[baseline=-0.7ex]\node[circle,fill=white,draw=popink,line width=1.3pt,minimum size=1.6cm,inner sep=0]{\popdisplay\fontsize{24}{24}\selectfont\color{poppurple}+};%
      \hspace{8pt}\begin{minipage}[c]{0.6\linewidth}
        {\popdisplay\fontsize{11}{12.5}\selectfont\color{white}\MakeUppercase{Passe à OnBuch+}}\par\vspace{4pt}
        {\raggedright\popsemi\fontsize{8.4}{11}\selectfont\color{white}Tous les cours signés, Tuteur IA illimité, fiches PDF — onglet OnBuch+ de l'app\par}
      \end{minipage}
    \end{tcolorbox}
  \end{minipage}
  \vspace{8pt}
  \popcontenu
  \vfill
  \signatureonbuch
  \restoregeometry
  \newpage
}

% Rangée « Ce cours contient » (page de garde)
\newcommand{\poptile}[3]{% #1 couleur #2 gros texte #3 légende
  \begin{tcolorbox}[enhanced,colback=#1,colframe=popink,boxrule=2pt,arc=9pt,shadow={2.5pt}{-2.5pt}{0pt}{popink},
      left=6pt,right=6pt,top=9pt,bottom=9pt,halign=center,valign=center,height=1.75cm,width=\linewidth,nobeforeafter]
    {\popdisplay\fontsize{12.5}{13}\selectfont\color{white}\MakeUppercase{#2}}\par\vspace{3pt}
    {\centering\popsemi\fontsize{7.8}{9.5}\selectfont\color{white}#3\par}
  \end{tcolorbox}}
\newcommand{\popcontenu}{%
  \par\noindent{\popxboldls\fontsize{7.5}{7.5}\selectfont\color{popmuted}CE COURS SIGNÉ CONTIENT}\par\vspace{7pt}
  \noindent\begin{minipage}[t]{0.235\linewidth}\poptile{popblue}{Cours}{complet, illustré, pas à pas}\end{minipage}\hfill
  \begin{minipage}[t]{0.235\linewidth}\poptile{poporange}{Méthodes}{et exercices résolus}\end{minipage}\hfill
  \begin{minipage}[t]{0.235\linewidth}\poptile{poppink}{Exercices}{type examen + corrigés}\end{minipage}\hfill
  \begin{minipage}[t]{0.235\linewidth}\poptile{popgreen}{Fiche bilan}{l'essentiel à retenir}\end{minipage}\par}

% Sommaire
\newtcolorbox{sommaire}{enhanced,colback=white,colframe=popink,boxrule=2.2pt,arc=10pt,
  drop shadow=popink,left=18pt,right=18pt,top=13pt,bottom=13pt,before skip=6pt,after skip=20pt,
  before upper={\poppill{Sommaire}{poppurple}{white}\par\vspace{9pt}\popsemi\fontsize{10.6}{14.5}\selectfont\color{popink}}}

% Signature de fin de cours — poussée tout en bas de la dernière page
\newcommand{\findecours}{%
  \par\vfill\nopagebreak
  \begin{center}\popsquiggle{poporange}\end{center}\vspace{4pt}
  \signatureonbuch
}
\AtBeginDocument{\def\llbracket{\mathopen{[\mkern-3.5mu[}}\def\rrbracket{\mathclose{]\mkern-3.5mu]}}} % intervalles d'entiers (glyphe ⟦ absent de la police)
% Glyphes absents de Fira Math : redéfinis à partir de glyphes disponibles
\AtBeginDocument{%
  \def\setminus{\mathbin{\backslash}}\let\smallsetminus\setminus
  \def\vdots{\mathord{\vbox{\baselineskip3.2pt\lineskiplimit0pt\kern4pt\hbox{.}\hbox{.}\hbox{.}}}}%
  \def\ddots{\mathinner{\mkern1mu\raise6.5pt\hbox{.}\mkern2mu\raise3.6pt\hbox{.}\mkern2mu\raise0.7pt\hbox{.}\mkern1mu}}%
  \def\triangle{\mathord{\scalebox{0.9}{$\symup{\Delta}$}}}%
  \def\nabla{\mathord{\raisebox{\depth}{\scalebox{1}[-1]{$\symup{\Delta}$}}}}%
  \def\checkmark{\popcheck{popdark}}%
  \def\star{\mathbin{\ast}}\def\bigstar{\mathord{\ast}}%
  \def\diamond{\mathbin{\scalebox{0.6}{$\Diamond$}}}%
  \def\bigcup{\mathop{\mathchoice{\raisebox{-0.3em}{\scalebox{1.7}{$\cup$}}}{\scalebox{1.25}{$\cup$}}{\cup}{\cup}}\displaylimits}%
  \def\bigcap{\mathop{\mathchoice{\raisebox{-0.3em}{\scalebox{1.7}{$\cap$}}}{\scalebox{1.25}{$\cap$}}{\cap}{\cap}}\displaylimits}%
  \def\lesseqqgtr{\mathrel{\lessgtr}}\def\bigoplus{\mathop{\oplus}\displaylimits}%
}
\providecommand{\ssi}{si et seulement si} % abréviation fréquente
\AtBeginDocument{\providecommand{\wideparen}{\overparen}} % arc de cercle

%% FIGFIX-BEGIN v4 — correctifs globaux des figures (docs/onbuch-generation/tools/figfix)
\usepackage{adjustbox}\usepackage{etoolbox}
% toute figure TikZ/pgfplots plus large que la ligne est réduite à la largeur disponible
\BeforeBeginEnvironment{tikzpicture}{\begin{adjustbox}{max width=\linewidth}}
\AfterEndEnvironment{tikzpicture}{\end{adjustbox}}
% légendes pgfplots lisibles par-dessus les courbes
\pgfplotsset{every axis/.append style={legend style={fill=white,fill opacity=0.92,text opacity=1,draw=black!15,inner sep=3pt}}}
% étiquettes d'arêtes (syntaxe edge["..."]) avec fond blanc
\tikzset{every edge quotes/.append style={fill=white,inner sep=1.2pt}}
%% FIGFIX-END
\begin{document}

\pagegarde{Produire le plan détaillé d’une dissertation}{Français}{Tle Tech}{Français – classe de Terminale technique}{N19}{12 séances (fiche de progression harmonisée)}{Cette leçon outille l'élève de Terminale technique pour analyser un sujet de dissertation, mobiliser ses connaissances, rechercher et organiser des idées, formuler une problématique et produire un plan détaillé conforme aux exigences du Baccalauréat technique. Elle s'appuie sur l'étude du Vieux nègre et la médaille de Ferdinand Oyono et sur des notions de langue (référent et substituts, sigles et emprunts) utiles à la rédaction.
\vspace{4pt}
\begin{multicols}{2}\small
\begin{itemize}
\item À la fin de cette leçon, je sais identifier et analyser les paratextes d'une œuvre littéraire pour préparer sa lecture
\item À la fin de cette leçon, je sais dégager le thème, la thèse et la visée d'un sujet de dissertation
\item À la fin de cette leçon, je sais mobiliser mes connaissances et rechercher des idées pertinentes par le brainstorming et le questionnement
\item À la fin de cette leçon, je sais formuler une problématique claire à partir d'un sujet donné
\item À la fin de cette leçon, je sais construire un plan détaillé (dialectique, thématique ou analytique) avec introduction, développement et conclusion
\item À la fin de cette leçon, je sais employer correctement les substituts du référent pour éviter les répétitions dans ma copie
\end{itemize}\end{multicols}}

\begin{sommaire}
\begin{multicols}{2}
\begin{enumerate}
\item Réception du texte : paratextes et lecture du texte ouvroir
\item Lecture méthodique n°1 et contrôle de lecture
\item Analyse du sujet de dissertation
\item Recherche des idées et formulation de la problématique
\item Élaboration et production du plan détaillé
\item Langue au service de la dissertation : référent et substituts, sigles et emprunts
\item Activité d'intégration
\item Méthodes et exercices résolus
\item Exercices
\item Corrigés des exercices
\item Fiche bilan
\end{enumerate}
\end{multicols}
\end{sommaire}

\begin{objectifs}
\begin{itemize}
\item À la fin de cette leçon, je sais identifier et analyser les paratextes d'une œuvre littéraire pour préparer sa lecture
\item À la fin de cette leçon, je sais dégager le thème, la thèse et la visée d'un sujet de dissertation
\item À la fin de cette leçon, je sais mobiliser mes connaissances et rechercher des idées pertinentes par le brainstorming et le questionnement
\item À la fin de cette leçon, je sais formuler une problématique claire à partir d'un sujet donné
\item À la fin de cette leçon, je sais construire un plan détaillé (dialectique, thématique ou analytique) avec introduction, développement et conclusion
\item À la fin de cette leçon, je sais employer correctement les substituts du référent pour éviter les répétitions dans ma copie
\item À la fin de cette leçon, je sais reconnaître et utiliser à bon escient les sigles et les emprunts dans un texte
\item À la fin de cette leçon, je sais rédiger et justifier une démarche, une réponse ou une conclusion dans une situation concrète
\item À la fin de cette leçon, je sais rendre compte d'une lecture et négocier un projet d'étude en groupe
\end{itemize}
\end{objectifs}

\begin{prerequis}
\begin{itemize}
\item Notions de base sur les genres littéraires (roman, essai) vues en Seconde et Première
\item La structure générale d'un texte argumentatif (thèse, arguments, exemples) vue en Première
\item Les connecteurs logiques et l'organisation du paragraphe argumentatif
\item La lecture méthodique d'un extrait de roman (Première)
\item Les notions de pronom, de reprise nominale et de champ lexical (classes antérieures)
\end{itemize}
\end{prerequis}

\newpage

\coursec{Réception du texte : paratextes et lecture du texte ouvroir}

Au lycée technique de Nkongsamba, les candidats au Probatoire et au Baccalauréat technique constatent chaque année que beaucoup de copies de dissertation obtiennent des notes médiocres, non pas par manque d'idées, mais parce que le plan est confus ou hors sujet. Un jury d'examen à Douala rapporte que 60 \% des échecs en dissertation viennent d'une mauvaise analyse du sujet. Comment transformer un sujet d'examen en un plan détaillé clair, logique et convaincant ? C'est toute l'ambition de cette leçon, qui s'appuie sur la lecture du \emph{Vieux nègre et la médaille} de Ferdinand Oyono, œuvre au programme.

Avant d'aborder la méthodologie de la dissertation, il est indispensable de maîtriser la réception des textes littéraires. La première étape de toute étude d'œuvre consiste à explorer son paratexte et à effectuer une lecture ouvroir (ou lecture de découverte) de l'incipit. Cette démarche permet de formuler des hypothèses de lecture pertinentes que l'on vérifiera au fil de l'étude approfondie.

\courssub{Qu'est-ce que le paratexte ?}

\begin{definition}[Paratexte]
Ensemble des éléments qui accompagnent et présentent un texte : titre, couverture, nom de l'auteur, préface, quatrième de couverture, dédicace, épigraphe, notes d'éditeur, etc. Le paratexte est le « seuil » par lequel on entre dans l'œuvre (Gérard Genette, \emph{Seuils}, 1987).
\end{definition}

Le paratexte se divise en deux catégories complémentaires :

\begin{itemize}
\item Le \textbf{péritexte} : éléments matériellement attachés au livre (couverture, titre, préface, 4e de couverture, dédicace, sommaire, notes de bas de page).
\item L'\textbf{épitexte} : éléments extérieurs au livre mais qui en parlent (critiques, interviews de l'auteur, bandeaux publicitaires, présentations en librairie, dossiers de presse, adaptations cinématographiques).
\end{itemize}

\begin{attention}[Erreur fréquente]
Ne pas confondre le résumé de la quatrième de couverture avec l'analyse de l'œuvre. Le paratexte sert à \textbf{anticiper} la lecture, à formuler des hypothèses, \textbf{pas à remplacer} l'étude du texte lui-même. Un élève qui recopie la 4e de couverture en guise d'introduction à son commentaire est hors sujet.
\end{attention}

\courssub{Analyse du paratexte du \emph{Vieux nègre et la médaille}}

\begin{popfigure}
\begin{tikzpicture}[
  node distance=1.2cm and 1.5cm,
  bubble/.style={draw=popblue, thick, fill=popblue!10, rounded corners=8pt, minimum width=3.2cm, minimum height=1.6cm, align=center, font=\small},
  centerbubble/.style={draw=poporange, thick, fill=poporange!15, rounded corners=12pt, minimum width=4.5cm, minimum height=2.2cm, align=center, font=\normalsize\bfseries},
  arrow/.style={-Stealth, popblue, thick}
]
\node[centerbubble, align=center] (center){\emph{Le Vieux nègre}\\\& la médaille};
\node[bubble, above left=of center] (titre) {Titre};
\node[bubble, above=of center, align=center] (auteur){Auteur :\\Ferdinand Oyono};
\node[bubble, above right=of center, align=center] (date){Publication :\\1956};
\node[bubble, right=of center, align=center] (editeur){Éditeur :\\René Julliard (Paris)};
\node[bubble, below right=of center, align=center] (quatre){4e de couv. :\\« Meka, chef de canton... »};
\node[bubble, below=of center, align=center] (genre){Genre :\\Roman / Nouvelles};
\node[bubble, below left=of center, align=center] (themes){Thèmes annoncés :\\Colonisation, ironie,\\récompense dérisoire};
\node[bubble, left=of center, align=center] (collection){Collection :\\« Lettres nouvelles »};

\draw[arrow] (center) -- (titre);
\draw[arrow] (center) -- (auteur);
\draw[arrow] (center) -- (date);
\draw[arrow] (center) -- (editeur);
\draw[arrow] (center) -- (quatre);
\draw[arrow] (center) -- (genre);
\draw[arrow] (center) -- (themes);
\draw[arrow] (center) -- (collection);
\end{tikzpicture}
\legende{Carte d'identité paratextuelle du roman : éléments du péritexte organisés autour de la couverture.}
\end{popfigure}

\begin{methode}[Analyser un paratexte : démarche en 4 étapes]
\begin{itemize}
\item \textbf{Observer} systématiquement chaque élément : titre, nom d'auteur, couverture, éditeur, collection, date, préface, 4e de couverture, dédicace, épigraphe.
\item \textbf{Relever les indices} signifiants : connotations du titre, réputation de l'auteur, maison d'édition (engagée ? commerciale ?), collection (littérature africaine ? classiques ?), résumé de la 4e de couverture (ton, informations données/cachées).
\item \textbf{Formuler des hypothèses de lecture} : quel genre ? quels thèmes ? quelle tonalité (ironique, tragique, réaliste) ? quel public visé ? quelle intention de l'auteur ?
\item \textbf{Prévoir de les vérifier} pendant la lecture intégrale : noter les passages qui confirment ou infirment chaque hypothèse.
\end{itemize}
\end{methode}

Appliquons cette méthode au \emph{Vieux nègre et la médaille} :

\begin{itemize}
\item \textbf{Titre} : « Le Vieux nègre » — terme péjoratif, colonial, qui renvoie à la déshumanisation ; « la médaille » — récompense officielle, symbole de la reconnaissance coloniale. L'association des deux suggère une \textbf{ironie} : la médaille offerte au « vieux nègre » est une récompense dérisoire pour une vie de soumission.
\item \textbf{Auteur} : Ferdinand Oyono, Camerounais, né en 1929. A déjà publié \emph{Une vie de boy} (1956), roman à charge contre l'administration coloniale. On s'attend à une œuvre \textbf{engagée}, critique.
\item \textbf{Date et éditeur} : 1956, chez René Julliard (Paris), collection « Lettres nouvelles » dirigée par Maurice Nadeau. Cette collection publie les voix nouvelles de la littérature francophone (Senghor, Césaire, Kateb Yacine). Le roman paraît \textbf{quatre ans avant l'indépendance du Cameroun} (1960), dans un contexte de montée des nationalismes africains.
\item \textbf{Quatrième de couverture} : Présente Meka, chef de canton décoré de la médaille de la Légion d'honneur. Le résumé insiste sur la cérémonie officielle et la fierté apparente de Meka, mais laisse deviner le regard critique du narrateur sur la mascarade.
\end{itemize}

\begin{exemplebox}[Hypothèses de lecture issues du paratexte]
\begin{enumerate}
\item \textbf{Thème central} : La colonisation française au Cameroun et ses relais locaux (chefs de canton).
\item \textbf{Tonalité} : Ironique, satirique, voire amère. Le titre même annonce un décalage entre la dignité humaine du personnage et l'appellation coloniale « vieux nègre ».
\item \textbf{Structure prévisible} : Une cérémonie de remise de médaille (point d'orgue), des flash-back sur la vie de Meka, une chute qui dévoile le caractère illusoire de la récompense.
\item \textbf{Personnage principal} : Meka, figure du « bon collaborateur » qui croit à la reconnaissance de l'administration, mais que le lecteur perçoit comme une victime consentante.
\end{enumerate}
\end{exemplebox}

\courssub{Présentation de l'auteur et contexte historique}

\begin{savaistu}[Ferdinand Oyono (1929--2010) : une double carrière]
Ferdinand Léopold Oyono n'est pas seulement romancier. Né à Ngoulémakong (région du Sud, Cameroun), il fait des études au séminaire de Yaoundé, puis en France (Sorbonne). Après l'indépendance, il devient \textbf{diplomate} (ambassadeur en Grande-Bretagne, aux Nations unies, en Algérie, au Maroc) et \textbf{ministre} camerounais (Fonction publique, puis Affaires étrangères). Ses trois romans — \emph{Une vie de boy} (1956), \emph{Le Vieux nègre et la médaille} (1956), \emph{Chemin d'Europe} (1960) — sont traduits en une vingtaine de langues et étudiés dans le monde entier. Il incarne la génération des « pères de la littérature africaine d'expression française » aux côtés de Mongo Beti, Sembène Ousmane, Cheikh Hamidou Kane.
\end{savaistu}

\begin{popfigure}
\begin{tikzpicture}[
  scale=0.9,
  event/.style={circle, fill=popblue, inner sep=2.5pt},
  eventlabel/.style={anchor=west, font=\small, text width=5.5cm, align=left},
  datelabel/.style={anchor=east, font=\small\bfseries, text width=2.2cm, align=right},
  axis/.style={popblue, thick}
]
\draw[axis] (0,0) -- (14,0);
\foreach \x/\year in {0/1929, 2.7/1956, 3.1/1960, 8.1/2010} {
  \draw[axis] (\x,0.15) -- (\x,-0.15) node[below, font=\small] {\year};
}
\node[event] (e1) at (0,0) {};
\node[datelabel] at (-0.3,0) {1929};
\node[eventlabel] at (0.3,0.4) {Naissance à Ngoulémakong (Cameroun)};
\node[event] (e2) at (2.7,0) {};
\node[datelabel] at (2.4,0) {1956};
\node[eventlabel] at (3.0,0.8) {Publication d'\emph{Une vie de boy} et du \emph{Vieux nègre et la médaille} (Julliard)};
\node[event] (e3) at (3.1,0) {};
\node[datelabel] at (2.8,0) {1960};
\node[eventlabel, fill=popgreen!20, rounded corners=3pt, inner sep=2pt] at (3.4,1.3) {\textbf{Indépendance du Cameroun (1er janv.)}};
\node[event] (e4) at (8.1,0) {};
\node[datelabel] at (7.8,0) {2010};
\node[eventlabel] at (8.4,0.4) {Décès à Yaoundé};
\end{tikzpicture}
\legende{Frise chronologique : vie de Ferdinand Oyono et repères historiques majeurs. L'œuvre paraît à la veille de l'indépendance, dans un contexte de décolonisation.}
\end{popfigure}

\begin{aretenir}[Contexte de publication (1956)]
\begin{itemize}
\item Cameroun sous administration coloniale française (mandat ONU puis tutelle).
\item Montée des mouvements nationalistes (UPC — Union des populations du Cameroun, interdite en 1955).
\item Littérature africaine d'expression française en plein essor : \emph{Contes d'Amadou Koumba} (Birago Diop, 1947), \emph{Le Maître de la parole} (Hampâté Bâ, 1955), \emph{Les Soleils des indépendances} (Ahmadou Kourouma, 1968 à venir).
\item Oyono publie deux romans la même année (1956) : coup d'éclat littéraire, entrée fracassante dans le paysage littéraire francophone.
\end{itemize}
\end{aretenir}

\courssub{Lecture du texte ouvroir (incipit)}

Le texte ouvroir correspond aux premières lignes du roman (incipit). Il donne les clés d'entrée dans l'œuvre : cadre spatio-temporel, personnages, point de vue narratif, tonalité.

\begin{exemplebox}[Incipit du \emph{Vieux nègre et la médaille} (extrait)]
\begin{quote}
\emph{« Meka se leva avant l'aube. Il faisait encore nuit dans la case. Il tâtonna pour trouver son pagne, l'enroula autour de ses reins et sortit. Dehors, l'air était frais. Le ciel était clair, constellé d'étoiles. Meka marcha vers la case de son fils aîné, Kelara. Il l'appela à voix basse : "Kelara, lève-toi. Aujourd'hui, c'est le grand jour." »}
\end{quote}
\end{exemplebox}

\begin{exoresolu}[Analyse guidée de l'incipit]
\textbf{Énoncé :} À partir de l'extrait ci-dessus, repérer le cadre spatio-temporel, les personnages présents, le point de vue narratif et la tonalité.
\tcblower
\textbf{Solution détaillée :}
\begin{enumerate}
\item \textbf{Cadre spatio-temporel} : « avant l'aube », « encore nuit », « air frais », « ciel constellé d'étoiles » $\rightarrow$ \textbf{matin très tôt, saison sèche} (nuits fraîches, ciel dégagé). Lieu : « la case » de Meka, puis « la case de son fils aîné Kelara » $\rightarrow$ \textbf{village camerounais, habitat traditionnel}.
\item \textbf{Personnages} : \textbf{Meka} (personnage principal, désigné par son prénom, verbes d'action « se leva », « marcha », « appela ») ; \textbf{Kelara} (fils aîné, mentionné seulement). Meka est le centre de la scène.
\item \textbf{Point de vue narratif} : Narrateur \textbf{hétérodiégétique} (ne participe pas à l'histoire), \textbf{focalisation interne} sur Meka (on suit ses gestes, ses sensations : « tâtonna », « l'air était frais » perçu par lui). Le narrateur sait ce que Meka fait et pense (« Aujourd'hui, c'est le grand jour »).
\item \textbf{Tonalité} : \textbf{Solennelle, rituelle} (« le grand jour »), contraste entre la simplicité du décor (case, pagne, nuit) et l'importance de l'événement annoncé. Aucune ironie explicite ici : l'ironie viendra du décalage entre la perception de Meka (fierté) et le regard du narrateur/lecteur.
\end{enumerate}
\end{exoresolu}

\courssub{Formulation et vérification des hypothèses de lecture}

À l'issue de la lecture ouvroir, on croise les indices du paratexte et ceux de l'incipit pour affiner ses hypothèses.

\begin{center}
\begin{tabularx}{\linewidth}{|l|X|X|}
\hline
\textbf{Indice paratexte / incipit} & \textbf{Hypothèse formulée} & \textbf{Mode de vérification pendant la lecture} \\
\hline
Titre : « Vieux nègre » (terme colonial) & Ironie du narrateur sur la déshumanisation & Repérer les occurrences de « vieux nègre » : qui le dit ? dans quelle bouche ? \\
Médailles / Légion d'honneur (4e couv.) & Critique de la collaboration avec le colonisateur & Suivre la cérémonie : discours, réactions de Meka, regard des villageois \\
Auteur : Oyono, \emph{Une vie de boy} (1956) & Tonalité satirique, engagement anti-colonial & Noter les passages où l'administration est tournée en dérision \\
Incipit : « le grand jour », focalisation sur Meka & Meka perçu comme naïf / victime consentante & Observer l'évolution de la conscience de Meka (éveil ? résignation ?) \\
Date : 1956 (avant indépendance) & Œuvre témoignant de la fin de la colonisation & Chercher les allusions aux événements politiques (UPC, répression) \\
\hline
\end{tabularx}
\end{center}

\begin{methode}[De la lecture ouvroir au projet d'étude]
\begin{itemize}
\item Lire le paratexte complet (couverture, 4e couv., préface s'il y a).
\item Lire l'incipit (première page ou premier chapitre) en annotant : cadre, personnages, narrateur, tonalité.
\item Dresser un tableau d'hypothèses (comme ci-dessus).
\item Lire l'œuvre intégrale en relevant, pour chaque hypothèse, les passages confirmatifs ou infirmatifs.
\item Rédiger un \textbf{compte rendu de lecture} structuré : fiche d'identité, résumé, analyse des personnages, thèmes, style, jugement personnel.
\item Négocier en classe un \textbf{projet d'étude} commun (axes d'analyse, passages clés à commenter, questions de dissertation potentielles).
\end{itemize}
\end{methode}

\begin{exercice}[Entraînement : paratexte et incipit]
On donne la quatrième de couverture suivante (fictive) :
\begin{quote}
\emph{« Dans un village du Nord-Cameroun, le chef traditionnel reçoit la visite d'un administrateur français. Autour d'un thé rituel, les non-dits s'accumulent. Qui manipule qui ? Un roman bref, incisif, où le rire côtoie la gravité. »}
\end{quote}
\begin{enumerate}
\item Relever les indices de genre, de thème, de tonalité, de cadre.
\item Formuler trois hypothèses de lecture.
\item Proposer un titre possible pour ce roman.
\end{enumerate}
\end{exercice}

\begin{corrige}[Exercice n°1]
\begin{enumerate}
\item \textbf{Indices} : « village du Nord-Cameroun » (cadre spatial précis, région sahélienne) ; « chef traditionnel » + « administrateur français » (confrontation pouvoir traditionnel / pouvoir colonial) ; « thé rituel » (coutume, diplomatie codifiée) ; « non-dits », « qui manipule qui ? » (tension, ironie, jeu de dupes) ; « rire côtoie la gravité » (tonalité satirique, tragi-comique).
\item \textbf{Hypothèses} : (1) Le roman met en scène l'imposture de la collaboration coloniale. (2) Le chef traditionnel use de ruse pour préserver son autorité réelle. (3) L'administrateur français est dépeint comme ignorant des codes locaux, donc vulnérable.
\item \textbf{Titre possible} : \emph{Le Thé du chef}, \emph{Non-dits sous le soleil}, \emph{La Main dans l'ombre}.
\end{enumerate}
\end{corrige}

\begin{attention}[Piège à éviter en examen]
Lors de l'épreuve de dissertation ou de commentaire, ne commencez \textbf{jamais} votre copie par un résumé de la quatrième de couverture. L'examinateur connaît l'œuvre. Montrez que vous avez \textbf{lu le texte} en citant des passages précis, en analysant la construction narrative, le style, les procédés d'ironie. Le paratexte n'est qu'un \textbf{outil de préparation}, pas le contenu de la copie.
\end{attention}

Cette première section a posé les bases de la réception de l'œuvre au programme. La maîtrise du paratexte et de la lecture ouvroir conditionne la qualité de l'analyse ultérieure : c'est en formulant des hypothèses rigoureuses dès l'abord du livre que l'on construit une lecture pertinente, capable de nourrir une dissertation solide. La section suivante abordera la lecture méthodique n°1 et le contrôle de lecture.

\coursec{Lecture méthodique n°1 et contrôle de lecture}

Dans la section précédente, nous avons franchi le seuil de l'œuvre de Ferdinand Oyono, \emph{Le Vieux nègre et la médaille} : l'examen des paratextes (couverture, quatrième de couverture, préface) et la lecture du texte ouvroir nous ont permis de formuler des hypothèses de lecture. Il s'agit maintenant de \emph{vérifier} ces hypothèses en entrant dans le texte lui-même. C'est le rôle de la \cle{lecture méthodique} : une lecture lente, active et questionnante, qui ne se contente pas de suivre l'histoire, mais qui cherche à comprendre \emph{comment} le texte produit son sens.

\courssub{Les étapes de la lecture méthodique}

Lire méthodiquement, ce n'est pas lire vite : c'est lire \emph{plusieurs fois}, avec un objectif différent à chaque passage. Intuitivement, on peut comparer cette démarche à celle d'un mécanicien qui examine un moteur : il ne se contente pas de constater que la voiture roule ; il ouvre le capot, identifie chaque pièce, et comprend le fonctionnement de l'ensemble.

\begin{methode}[Conduire une lecture méthodique]
\begin{enumerate}
\item \textbf{Lire silencieusement} l'extrait une première fois, crayon en main, pour saisir le sens global et souligner ce qui surprend ou résiste.
\item \textbf{Lire à voix haute} (lecture magistrale ou lecture par un élève) : la voix révèle les rythmes, les ruptures de ton, l'ironie.
\item \textbf{Questionner} le texte : qui parle ? à qui ? de quoi ? que se passe-t-il ? quels sentiments se dégagent ?
\item \textbf{Expliquer les mots difficiles} : deviner le sens par le contexte, puis vérifier dans le dictionnaire.
\item \textbf{Dégager le sens global} en une ou deux phrases, sans raconter.
\item \textbf{Vérifier les hypothèses} formulées lors de la lecture du texte ouvroir : sont-elles confirmées, infirmées, nuancées ?
\end{enumerate}
\end{methode}

\begin{attention}[L'erreur à éviter]
L'erreur la plus fréquente consiste à \emph{résumer} l'extrait au lieu de l'\emph{analyser}. Raconter ce qui arrive à Meka n'est pas faire une lecture méthodique : la lecture méthodique vise le \cle{sens} (ce que le texte veut dire) et les \cle{procédés} (comment il le dit), pas le simple récit des événements. Au contrôle, une copie qui résume ne peut pas obtenir la moyenne.
\end{attention}

\courssub{La situation d'énonciation : qui parle, à qui, où, quand ?}

Avant d'analyser le contenu, il faut situer la parole. Tout texte est prononcé par quelqu'un, dans des circonstances précises, et ces circonstances éclairent le sens.

\begin{definition}[Situation d'énonciation]
La \cle{situation d'énonciation} désigne l'ensemble des circonstances dans lesquelles un texte est produit : le \textbf{locuteur} (qui parle ?), le \textbf{destinataire} (à qui ?), le \textbf{lieu} (où ?), le \textbf{moment} (quand ?) et la \textbf{visée} (pour quoi faire ?).
\end{definition}

\begin{definition}[Focalisation narrative]
La \cle{focalisation} définit la perspective par laquelle l'histoire est présentée. On distingue : la \textbf{focalisation zéro} (narrateur omniscient, sait tout) ; la \textbf{focalisation interne} (le récit est filtré par la conscience d'un personnage, le lecteur ne sait que ce que ce personnage sait/perçoit) ; la \textbf{focalisation externe} (le narrateur se comporte comme une caméra, il ne donne pas accès aux pensées).
\end{definition}

Appliquons cette grille au premier extrait de \emph{Le Vieux nègre et la médaille} :

\begin{center}
\begin{tabularx}{\linewidth}{|l|X|}
\hline
\rowcolor{popblueL} \textbf{Question} & \textbf{Réponse dans l'extrait} \\
\hline
Qui parle ? & Un narrateur à la troisième personne (narrateur hétéradiégétique), qui adopte souvent le point de vue de Meka (\textbf{focalisation interne variable}). \\
\hline
À qui ? & Au lecteur, invité à partager le regard ironique du narrateur. \\
\hline
Où ? & À Doum, village africain sous administration coloniale. \\
\hline
Quand ? & À l'époque coloniale, la veille du 14 juillet, fête nationale française. \\
\hline
Pour quoi faire ? & Pour dénoncer, par l'ironie, l'absurdité de la récompense coloniale. \\
\hline
\end{tabularx}
\end{center}

Ce relevé fait apparaître un point essentiel : le narrateur ne dit jamais frontalement « la colonisation est injuste ». Il laisse le lecteur le comprendre. C'est toute la force de l'\cle{ironie}.

\courssub{Identification des personnages et réseau des relations}

La lecture méthodique exige de repérer les personnages et leurs relations. Trois groupes se dégagent autour du héros :

\begin{itemize}
\item \textbf{Meka}, le « vieux nègre » : protagoniste, vieillard du village de Doum, qui a donné ses terres à la mission et ses fils à la guerre. Il incarne l'Africain soumis qui a tout sacrifié à la colonisation.
\item \textbf{Les autorités coloniales} : le commandant de cercle, le père Vandermayer, les « Blancs » qui décident de décerner la médaille. Ils représentent le pouvoir.
\item \textbf{Les villageois} : Kelara (l'épouse de Meka), Engamba, les anciens et les gens de Doum, témoins admiratifs ou sceptiques de la récompense.
\end{itemize}

\begin{popfigure}
\begin{center}
\begin{tikzpicture}[scale=1, every node/.style={font=\small}]
% Noeud central
\node[draw, fill=popgoldL, rounded corners, thick, minimum width=3.2cm, minimum height=1cm] (meka) at (0,0) {\textbf{Meka}, le vieux nègre};
% Colonisateurs
\node[draw, fill=popblueL, rounded corners, minimum width=3.4cm] (cmdt) at (-4.5,2.2) {Le commandant de cercle};
\node[draw, fill=popblueL, rounded corners, minimum width=3.4cm] (pere) at (-4.5,-2.2) {Le père Vandermayer};
% Villageois
\node[draw, fill=popgreenL, rounded corners, minimum width=2.6cm] (kelara) at (4.5,2.2) {Kelara, son épouse};
\node[draw, fill=popgreenL, rounded corners, minimum width=2.6cm] (engamba) at (4.5,0) {Engamba, son ami};
\node[draw, fill=popgreenL, rounded corners, minimum width=2.6cm] (village) at (4.5,-2.2) {Les gens de Doum};
% Relations
\draw[<->, thick, popblue] (meka) -- (cmdt) node[midway, above, sloped, font=\scriptsize] {médaille / méfiance};
\draw[<->, thick, popblue] (meka) -- (pere) node[midway, below, sloped, font=\scriptsize] {terres données à la mission};
\draw[<->, thick, popgreen] (meka) -- (kelara) node[midway, above, sloped, font=\scriptsize] {vie commune, épreuves};
\draw[<->, thick, popgreen] (meka) -- (engamba) node[midway, above, font=\scriptsize] {amitié};
\draw[<->, thick, popgreen] (meka) -- (village) node[midway, below, sloped, font=\scriptsize] {admiration, jalousie};
\end{tikzpicture}
\end{center}
\legende{Réseau des personnages autour de Meka : à gauche, le monde des colonisateurs (bleu) ; à droite, la communauté villageoise (vert). Meka est pris entre les deux mondes.}
\end{popfigure}

Ce schéma fait apparaître la position centrale et inconfortable de Meka : il est le \emph{pont} entre deux univers qui ne se comprennent pas, et la médaille va précisément cristalliser cette tension.

\courssub{Dégagement du thème : la médaille et sa signification ironique}

Le thème d'un extrait n'est pas son sujet matériel (« on donne une médaille à un vieillard ») mais l'\emph{idée} que le texte développe. Ici, la médaille de l'administration coloniale, censée récompenser « cinquante ans de bons et loyaux services », apparaît comme une récompense \emph{dérisoire} : une vie entière de sacrifices — les terres, les fils morts à la guerre, la soumission — se solde par un morceau de métal et un ruban.

\begin{exemplebox}[La médaille, symbole ironique]
La médaille offerte à Meka symbolise la récompense dérisoire que la colonisation accorde à une vie de soumission. L'écart entre l'immensité du sacrifice (une existence entière) et l'insignifiance du prix (une médaille sans valeur) crée l'ironie : plus le narrateur insiste sur la solennité de la cérémonie, plus le lecteur mesure l'injustice de l'échange.
\end{exemplebox}

\courssub{Relevé des champs lexicaux : récompense et humiliation}

Le relevé du \cle{champ lexical} est un outil de lecture méthodique : on regroupe les mots qui appartiennent au même domaine de sens pour faire apparaître la cohérence — ou les tensions — du texte.

\begin{center}
\begin{tabularx}{\linewidth}{|X|X|}
\hline
\rowcolor{popblueL} \textbf{Champ lexical de la récompense} & \textbf{Champ lexical de l'humiliation et de la dérision} \\
\hline
médaille, honneur, décoration, ruban, cérémonie, fête, discours officiel, « bons et loyaux services » & vieux nègre, soumission, obéissance, terres perdues, fils morts, fatigue, vieillesse, ridicule, mensonge \\
\hline
\end{tabularx}
\end{center}

La confrontation des deux colonnes est révélatrice : le vocabulaire glorieux de la récompense est constamment \emph{contredit} par le vocabulaire de la servitude. Le titre même, qui jouxte « vieux nègre » (mépris) et « médaille » (gloire), annonce cette collision ironique.

\begin{exoresolu}[Repérer la situation d'énonciation et la focalisation]
Énoncé : dans l'extrait suivant, inspiré de l'incipit du roman — « Le vieux Meka, allongé sur sa natte, songeait que demain les Blancs lui donneraient une médaille » —, identifie le locuteur, le type de focalisation et l'effet produit sur le lecteur.
\tcblower
Solution détaillée : le locuteur est un narrateur extérieur (troisième personne : « le vieux Meka », « il ») : c'est un narrateur \emph{hétérodiégétique}. Cependant, la formulation « songeait que » donne accès aux pensées intimes du personnage : c'est une \cle{focalisation interne} (le récit est filtré par la conscience de Meka). Effet produit : le lecteur entre dans la tête de Meka et partage son attente naïve ; il pourra donc mesurer, mieux que le héros lui-même, le caractère dérisoire de la récompense. Le narrateur installe ainsi, dès l'incipit, un regard ironique sans jamais juger ouvertement.
\end{exoresolu}

\courssub{Le contrôle de lecture : questionnaire de compréhension}

Le contrôle de lecture vérifie que la lecture méthodique a porté ses fruits. Il comporte deux niveaux : la \textbf{compréhension globale} (le sens d'ensemble) et la \textbf{compréhension fine} (le détail significatif).

\begin{exercice}[Contrôle de lecture n°1]
\textbf{A. Compréhension globale}
\begin{enumerate}
\item En une phrase, sans raconter, dégage le thème de l'extrait étudié.
\item Quelle est la situation d'énonciation de l'extrait (qui parle, où, quand) ?
\item Quel sentiment général se dégage de l'évocation de la cérémonie ?
\end{enumerate}
\textbf{B. Compréhension fine}
\begin{enumerate}
\item Releve trois mots du champ lexical de la récompense et trois mots du champ lexical de l'humiliation.
\item Pourquoi la date du 14 juillet est-elle significative ?
\item Qu'a donné Meka à la mission ? Qu'est devenu son fils ?
\item Explique l'ironie contenue dans l'expression « bons et loyaux services ».
\end{enumerate}
\end{exercice}

\begin{corrige}[Contrôle de lecture n°1]
\textbf{A.} 1. Le thème : la récompense dérisoire accordée par l'administration coloniale à une vie de soumission. 2. Un narrateur à la troisième personne (hétérodiégétique) avec focalisation interne sur Meka, à Doum, la veille du 14 juillet, à l'époque coloniale. 3. Un mélange de solennité et de malaise : la fête cache une injustice.
\textbf{B.} 1. Récompense : médaille, honneur, cérémonie ; humiliation : soumission, terres perdues, fils morts. 2. Le 14 juillet célèbre la liberté du peuple français : décorer un Africain soumis ce jour-là est une contradiction ironique. 3. Meka a donné ses terres à la mission catholique ; son fils est mort à la guerre pour la France. 4. L'expression est le cliché administratif qui maquille l'exploitation en mérite : les « services » de Meka sont en réalité des sacrifices non reconnus.
\end{corrige}

\courssub{Compte rendu oral et négociation du projet d'étude}

La lecture méthodique débouche sur un \cle{compte rendu oral} : chaque groupe présente en cinq minutes les résultats de sa lecture (situation d'énonciation, personnages, thème, champs lexicaux, difficultés rencontrées). Ce travail collectif permet ensuite de \emph{négocier le projet d'étude} de l'œuvre : répartition des extraits entre les groupes, calendrier des lectures, supports à produire (fiches, affiches, exposés).

\begin{center}
\begin{tabularx}{\linewidth}{|l|X|l|l|}
\hline
\rowcolor{popblueL} \textbf{Séance} & \textbf{Extrait à lire} & \textbf{Groupe responsable} & \textbf{Support attendu} \\
\hline
Séance 3 & L'annonce de la médaille (chap. 1) & Groupe 1 & Fiche de lecture \\
\hline
Séance 5 & Les préparatifs de la fête (chap. 2) & Groupe 2 & Affiche + exposé \\
\hline
Séance 7 & La cérémonie du 14 juillet (chap. 3) & Groupe 3 & Fiche de lecture \\
\hline
Séance 9 & La nuit de Meka et la pluie (chap. 4) & Groupe 4 & Exposé oral \\
\hline
Séance 11 & Bilan de l'œuvre et synthèse & Tous & Synthèse collective \\
\hline
\end{tabularx}
\end{center}

\begin{savaistu}[Le compte rendu au Baccalauréat technique]
Le compte rendu de lecture n'est pas un exercice scolaire sans lendemain : il est évalué au Probatoire et au Baccalauréat technique sous forme de \emph{question de synthèse} ou de question de lecture. Savoir dégager un thème, relever un champ lexical et présenter clairement le sens d'un extrait est donc une compétence directement rentable le jour de l'examen.
\end{savaistu}

\courssub{Compte rendu et remédiation}

La dernière étape est la \cle{remédiation} : le professeur et la classe corrigent collectivement les erreurs de compréhension relevées pendant le contrôle. Trois erreurs reviennent souvent :

\begin{itemize}
\item \textbf{Confondre résumé et analyse} : raconter la cérémonie au lieu d'expliquer son ironie. Remède : reformuler chaque réponse en commençant par « le texte montre que... ».
\item \textbf{Ignorer le contexte colonial} : lire l'histoire comme un simple fait divers. Remède : revenir à la situation d'énonciation et à la date du 14 juillet.
\item \textbf{Mal expliquer les mots difficiles} : par exemple « dérisoire » (qui n'a presque aucune valeur) ou « énonciation ». Remède : deviner par le contexte, puis vérifier au dictionnaire.
\end{itemize}

\begin{aretenir}[Bilan de la section]
La lecture méthodique suit six étapes : lire silencieusement, lire à voix haute, questionner, expliquer les mots difficiles, dégager le sens global, vérifier les hypothèses. Elle repose sur la situation d'énonciation (qui ? à qui ? où ? quand ?), l'identification de la focalisation (ici interne sur Meka), l'identification des personnages (Meka entre colonisateurs et villageois), le dégagement du thème (la médaille, récompense dérisoire et ironique) et le relevé des champs lexicaux (récompense contre humiliation). Elle se conclut par un contrôle de lecture, un compte rendu oral et une remédiation collective.
\end{aretenir}

Cette première lecture méthodique nous a appris à interroger un texte au lieu de le subir. Nous sommes désormais prêts à transposer cette rigueur dans la production d'écrits : la prochaine étape consistera à analyser un sujet de dissertation avec la même exigence.

\coursec{Analyse du sujet de dissertation}

La lecture méthodique du \emph{Vieux nègre et la médaille} de Ferdinand Oyono nous a entraînés dans l'univers du roman~: nous avons rencontré Meka, vieillard décoré par l'administration coloniale, et nous avons compris que cette médaille, loin d'être un honneur, est un instrument d'humiliation. Mais au Baccalauréat technique, l'épreuve de dissertation ne vous demandera pas de raconter l'histoire de Meka~: elle vous demandera de \emph{réfléchir} à partir d'une question. Par exemple~: « La médaille récompense-t-elle vraiment le mérite~? » Avant d'écrire la moindre phrase, il faut donc franchir une étape décisive~: \cle{l'analyse du sujet}. C'est elle que nous étudions maintenant.

\courssub{Qu'est-ce qu'une dissertation~?}

\begin{definition}[La dissertation]
La \cle{dissertation} est un exercice écrit de \textbf{réflexion personnelle et argumentée} sur un sujet donné. Le candidat y expose, de manière organisée et progressive, sa réflexion sur une question, en s'appuyant sur ses connaissances littéraires, culturelles et son expérience de la vie.
\end{definition}

Trois mots de cette définition méritent d'être pesés. \emph{Personnelle}~: l'examinateur attend votre pensée, pas la récitation d'un corrigé appris par cœur. \emph{Argumentée}~: chaque idée doit être justifiée par un raisonnement et illustrée par un exemple (une œuvre, un fait historique, une situation de la vie courante). \emph{Sur un sujet donné}~: vous ne choisissez pas votre question~; vous devez répondre exactement à celle qui est posée. C'est ce dernier point qui rend l'analyse du sujet si importante.

\begin{attention}[La première cause d'échec au Bac]
Au Baccalauréat technique, la majorité des copies faibles ne manquent pas d'idées~: elles répondent \textbf{à côté de la question}. Un candidat qui traite brillamment un autre sujet que celui posé obtient une note médiocre. Analyser le sujet n'est donc pas une formalité~: c'est la moitié du travail. Prévoyez au brouillon 20 à 30 minutes pour cette étape.
\end{attention}

\courssub{Premier geste~: lire attentivement le sujet}

L'analyse commence par une lecture en deux temps. La \textbf{première lecture} donne une impression globale~: de quoi parle-t-on, en gros~? La \textbf{seconde lecture}, crayon en main, consiste à \cle{souligner les mots-clés}, c'est-à-dire les mots qui portent le sens du sujet et qu'on ne peut ni supprimer ni remplacer sans changer la question.

Reprenons notre sujet fil conducteur~: « La médaille récompense-t-elle vraiment le mérite~? » Les mots-clés sont~: \emph{médaille} (la récompense officielle, la distinction), \emph{vraiment} (adverbe qui installe le doute~: il invite à vérifier, à discuter), \emph{mérite} (la valeur personnelle, le travail accompli). Le mot « vraiment » est capital~: sans lui, le sujet serait une affirmation~; avec lui, c'est une question ouverte.

Il faut aussi identifier le \cle{domaine} du sujet, c'est-à-dire le champ de la vie ou de la pensée auquel il appartient~:

\begin{itemize}
\item le domaine \textbf{littéraire}~: « Le roman a-t-il pour seule fonction de divertir~? »~;
\item le domaine \textbf{social}~: « Le travail divise-t-il les hommes ou les unit-il~? »~;
\item le domaine \textbf{moral}~: « Faut-il toujours dire la vérité~? »~;
\item le domaine \textbf{politique}~: « La colonisation a-t-elle apporté le progrès aux peuples africains~? »
\end{itemize}

Notre sujet sur la médaille relève du domaine moral et social~: il interroge la justice des récompenses.

\courssub{Dégager le thème et la thèse}

\begin{definition}[Le thème]
Le \cle{thème} est le sujet général traité dans le texte ou posé par l'énoncé. Il répond à la question~: « De quoi parle le sujet~? » Exemples de thèmes~: la colonisation, le travail, la justice, la récompense, l'amitié.
\end{definition}

\begin{definition}[La thèse]
La \cle{thèse} est le point de vue défendu ou discuté dans le sujet. Elle répond à la question~: « Que dit ou que suppose le sujet à propos du thème~? » Elle peut être \textbf{explicite} (énoncée clairement dans le sujet) ou \textbf{implicite} (sous-entendue, à reconstruire par le candidat).
\end{definition}

La distinction est simple~: le thème tient en un ou deux mots~; la thèse est une phrase complète, une idée qu'on peut approuver ou contester.

\begin{exemplebox}[Thème et thèse dans notre sujet]
Sujet~: « La médaille récompense-t-elle vraiment le mérite~? »
\begin{itemize}
\item \textbf{Thème}~: la récompense (les distinctions honorifiques).
\item \textbf{Thèse à discuter}~: la récompense est proportionnée au mérite~; autrement dit, on ne décore que ceux qui le méritent. Cette thèse est implicite~: c'est l'idée que le commun des mortels attache aux médailles, et que le sujet nous invite à vérifier.
\end{itemize}
Dans \emph{Le Vieux nègre et la médaille}, Meka reçoit une médaille non pour son mérite personnel, mais pour servir la propagande coloniale~: voilà un contre-exemple qui permettra de discuter la thèse.
\end{exemplebox}

\courssub{Identifier le type de sujet et les consignes}

Tous les sujets de dissertation ne se traitent pas de la même manière. On distingue classiquement trois types de sujets, reconnaissables à leur formulation et à leur consigne.

\begin{center}
\begin{tabularx}{\linewidth}{|l|X|X|X|}
\hline
\rowcolor{popblueL} \textbf{Critère} & \textbf{Sujet d'opinion} & \textbf{Sujet de réflexion} & \textbf{Sujet de discussion} \\
\hline
Forme & Affirmation à juger & Question ouverte sur une notion & Formule polémique, deux points de vue possibles \\
\hline
Consigne typique & « Qu'en pensez-vous~? », « Partagez-vous ce point de vue~? » & « Commentez », « Réfléchissez sur... » & « Discutez », « Discutez cette opinion » \\
\hline
Attitude attendue & Prendre position clairement (oui / non / nuance) & Approfondir, définir, explorer la notion & Confronter deux thèses adverses avant de conclure \\
\hline
Exemple type Bac technique & « L'école est la seule voie de la réussite. Qu'en pensez-vous~? » & « Commentez cette pensée~: la lecture élargit l'esprit. » & « La médaille récompense-t-elle vraiment le mérite~? Discutez. » \\
\hline
\end{tabularx}
\end{center}

Le \cle{repérage des consignes} est donc une étape à part entière. Le verbe de la consigne est un ordre~: « Discutez » exige la confrontation d'arguments contraires~; « Commentez » exige l'explication approfondie d'une pensée~; « Qu'en pensez-vous~? » exige une prise de position personnelle et assumée. Ignorer la consigne, c'est déjà désobéir à l'énoncé.

\courssub{Le schéma de l'entonnoir~: du sujet à la problématique}

L'analyse du sujet fonctionne comme un entonnoir~: on part de l'énoncé brut, large, et on resserre progressivement jusqu'à la question précise qui guidera tout le devoir (la problématique, que nous étudierons en détail dans la section suivante).

\begin{popfigure}
\begin{center}
\begin{tikzpicture}[font=\small, >=Stealth]
% Niveau 1
\fill[popblue!25] (-4.5,4) -- (4.5,4) -- (3.2,2.6) -- (-3.2,2.6) -- cycle;
\node at (0,3.3) {\textbf{SUJET} : l'énoncé complet, lu deux fois};
% Flèche
\draw[->, thick, popblue] (0,2.6) -- (0,2.4);
% Niveau 2
\fill[popgreen!25] (-3.2,2.4) -- (3.2,2.4) -- (2.1,1.2) -- (-2.1,1.2) -- cycle;
\node at (0,1.8) {\textbf{MOTS-CLÉS} : soulignés, domaine identifié};
% Flèche
\draw[->, thick, popgreen] (0,1.2) -- (0,1.0);
% Niveau 3
\fill[poporange!25] (-2.1,1.0) -- (2.1,1.0) -- (1.2,-0.2) -- (-1.2,-0.2) -- cycle;
\node at (0,0.4) {\textbf{THÈME} : de quoi parle le sujet~?};
% Flèche
\draw[->, thick, poporange] (0,-0.2) -- (0,-0.4);
% Niveau 4
\fill[poppink!30] (-1.2,-0.4) -- (1.2,-0.4) -- (0.55,-1.5) -- (-0.55,-1.5) -- cycle;
\node at (0,-0.95) {\textbf{THÈSE} : quel point de vue est posé~?};
% Flèche
\draw[->, thick, poppink] (0,-1.5) -- (0,-1.7);
% Niveau 5
\fill[poppurple!25] (-0.55,-1.7) -- (0.55,-1.7) -- (0,-2.7) -- cycle;
\node[text=white] at (0,-2.15) {\textbf{PROBLÉMATIQUE}};
\end{tikzpicture}
\end{center}
\legende{Schéma en entonnoir de l'analyse du sujet~: chaque étape resserre la réflexion, de l'énoncé brut (en haut) jusqu'à la problématique (en bas), question directrice du devoir.}
\end{popfigure}

\courssub{Méthode~: analyser un sujet en quatre étapes}

\begin{methode}[Les quatre étapes de l'analyse du sujet]
\begin{enumerate}
\item \textbf{Lire deux fois} le sujet, lentement, sans rien écrire la première fois, crayon en main la seconde.
\item \textbf{Souligner les mots-clés} et identifier le domaine (littéraire, social, moral, politique) ainsi que le type de sujet et la consigne (« Discutez », « Commentez », « Qu'en pensez-vous~? »).
\item \textbf{Dégager le thème} (de quoi parle-t-on~?) puis \textbf{la thèse} explicite ou implicite (quelle idée est affirmée ou mise en doute~?).
\item \textbf{Reformuler le sujet} avec ses propres mots, en une ou deux phrases, pour vérifier qu'on l'a compris.
\end{enumerate}
\end{methode}

La quatrième étape est votre assurance contre le hors sujet. Si vous savez reformuler le sujet autrement, c'est que vous l'avez compris~; si vous ne pouvez que le répéter mot à mot, c'est qu'il reste opaque.

\begin{attention}[Le hors sujet partiel~: l'erreur la plus sanctionnée]
L'erreur fréquente consiste à \textbf{traiter le thème général au lieu du sujet précis}. Face à « La médaille récompense-t-elle vraiment le mérite~? », le candidat imprudent écrit tout ce qu'il sait sur les récompenses~: les prix scolaires, les décorations, les trophées sportifs... sans jamais répondre à la question posée (la récompense est-elle \emph{proportionnée au mérite}~?). C'est un \cle{hors sujet partiel}, lourdement sanctionné au Baccalauréat. Règle d'or~: chaque paragraphe du devoir doit pouvoir se rattacher visiblement à la question exacte du sujet.
\end{attention}

\courssub{Application~: analyse de sujets types du Baccalauréat technique}

\begin{exoresolu}[Analyse complète d'un sujet type Bac technique]
Énoncé~: « Un ami vous affirme~: "Les distinctions honorifiques ne sont jamais données aux plus méritants." Discutez cette opinion en vous appuyant sur vos lectures et votre expérience. »
\tcblower
\textbf{Solution détaillée, en suivant les quatre étapes.}
\begin{enumerate}
\item \textbf{Lectures}~: première lecture — il s'agit des récompenses et de leur justice. Seconde lecture — je souligne les mots-clés.
\item \textbf{Mots-clés et consigne}~: \emph{distinctions honorifiques} (= médailles, prix, décorations), \emph{jamais} (adverbe absolu, donc thèse radicale), \emph{plus méritants} (le critère du mérite). Consigne~: « Discutez » — il s'agit donc d'un \textbf{sujet de discussion}~: je devrai confronter le pour et le contre. Domaine~: moral et social.
\item \textbf{Thème}~: la récompense. \textbf{Thèse à discuter} (explicite cette fois)~: les récompenses sont distribuées injustement, jamais selon le mérite.
\item \textbf{Reformulation}~: « Les médailles et les prix sont-ils toujours attribués à ceux qui le méritent, ou servent-ils parfois d'autres intérêts (favoritisme, propagande, apparence)~? »
\end{enumerate}
\textbf{Vérification anti-hors sujet}~: je ne dois pas raconter l'histoire de Meka, mais m'en servir comme \emph{exemple}~: sa médaille récompense sa soumission à l'administration coloniale, non son mérite — ce qui soutient la thèse de l'ami. En revanche, un élève premier de sa classe décoré lors de la fête de la jeunesse illustre le cas où la récompense suit le mérite. La discussion est ainsi équilibrée.
\end{exoresolu}

\begin{exercice}[À vous de jouer]
Analysez les deux sujets suivants selon les quatre étapes de la méthode (mots-clés, domaine, type de sujet, thème, thèse, reformulation)~:
\begin{enumerate}
\item « Le travail est-il toujours une source de dignité pour l'homme~? »
\item « La colonisation a apporté le développement à l'Afrique. Qu'en pensez-vous~? »
\end{enumerate}
\end{exercice}

\begin{corrige}[Exercice]
\textbf{Sujet 1.} Mots-clés~: \emph{travail}, \emph{toujours}, \emph{dignité}. Domaine~: social et moral. Type~: sujet de discussion (question ouverte appelant pour et contre). Thème~: le travail. Thèse implicite~: le travail confère la dignité à l'homme. Reformulation~: « Le travail rend-il l'homme digne dans tous les cas, ou existe-t-il des travaux qui avilissent (travail forcé, exploitation)~? »
\textbf{Sujet 2.} Mots-clés~: \emph{colonisation}, \emph{développement}, \emph{Afrique}. Domaine~: politique et historique. Type~: sujet d'opinion (affirmation + « Qu'en pensez-vous~? »). Thème~: la colonisation. Thèse explicite~: la colonisation a été bénéfique à l'Afrique. Reformulation~: « Peut-on affirmer, comme le prétendent certains, que la présence coloniale a fait progresser l'Afrique, ou faut-il y voir d'abord une exploitation et une humiliation, comme le montre le sort de Meka dans \emph{Le Vieux nègre et la médaille}~? »
\end{corrige}

\begin{aretenir}[L'essentiel de l'analyse du sujet]
\begin{itemize}
\item La dissertation est une réflexion personnelle et argumentée sur un sujet \emph{donné}~: il faut répondre à la question posée, et à elle seule.
\item Analyser un sujet, c'est descendre l'entonnoir~: sujet $\rightarrow$ mots-clés $\rightarrow$ thème $\rightarrow$ thèse $\rightarrow$ problématique.
\item Le thème tient en un mot (la récompense)~; la thèse est une idée complète à discuter (la récompense suit le mérite).
\item La consigne (« Discutez », « Commentez », « Qu'en pensez-vous~? ») détermine le type de sujet et l'attitude à adopter.
\item La reformulation du sujet en ses propres mots est le meilleur rempart contre le hors sujet partiel.
\end{itemize}
\end{aretenir}

Le sujet est désormais démonté, compris, reformulé. Mais une question claire ne fait pas encore un devoir~: il faut maintenant mobiliser des idées, des arguments et des exemples pour y répondre. C'est l'objet de la section suivante~: la recherche des idées et la formulation de la problématique.

\coursec{Recherche des idées et formulation de la problématique}

Dans la section précédente, nous avons appris à \cle{analyser le sujet} de dissertation : repérer les mots-clés, dégager le thème, identifier le type de sujet et cerner la \emph{tension} qu'il contient. Mais une fois le sujet compris, une question se pose immédiatement : \emph{que vais-je écrire ?} Beaucoup d'élèves se jettent alors sur la copie et rédigent au fil de la plume. C'est la pire des stratégies. Entre l'analyse du sujet et la rédaction, il existe deux étapes décisives : la \cle{recherche des idées} et la \cle{formulation de la problématique}. C'est l'objet de cette section. Nous travaillerons sur le sujet fil conducteur de l'unité, né de notre lecture du \emph{Vieux nègre et la médaille} de Ferdinand Oyono : \og Faut-il récompenser les anciens serviteurs de l'État ? \fg.

\courssub{Les techniques de recherche des idées}

\textbf{Intuition.} Imagine un cultivateur de l'Ouest qui prépare son champ. Il ne sème pas n'importe où : il commence par ratisser large, il ramasse tout ce que la terre contient, puis seulement il trie les bonnes graines des mauvaises herbes. La recherche des idées fonctionne exactement ainsi : d'abord \emph{produire en quantité}, ensuite \emph{sélectionner en qualité}.

\begin{definition}[Le brainstorming ou remue-méninges]
Le \cle{brainstorming} (ou remue-méninges) est une technique qui consiste à noter, sans aucun jugement ni censure, toutes les idées, tous les mots, tous les exemples que le sujet fait naître dans l'esprit. Il se pratique en deux temps : d'abord \textbf{individuellement} (chacun couche sur le papier ses propres idées pendant cinq à dix minutes), puis \textbf{collectivement} (la classe ou le groupe met en commun, ce qui enrichit le stock initial).
\end{definition}

\begin{methode}[Conduire un brainstorming efficace]
\begin{enumerate}
\item Recopier le sujet en haut du brouillon et souligner ses mots-clés.
\item Pendant cinq minutes, noter \emph{tout} ce qui vient : idées, exemples, souvenirs, proverbes, noms d'œuvres. Ne rien rejeter à ce stade.
\item Passer au questionnement systématique (voir ci-dessous) pour traquer les idées oubliées.
\item Mettre en commun avec le groupe : chaque idée d'un camarade peut en déclencher une nouvelle chez toi.
\item Seulement après cette récolte, passer au tri.
\end{enumerate}
\end{methode}

\textbf{Le questionnement : l'outil qui débloque les idées.} Quand le cerveau reste muet devant le sujet, il faut l'interroger avec les questions classiques du questionnaire : \emph{qui ? quoi ? où ? quand ? comment ? pourquoi ?} Appliquons-les à notre sujet :

\begin{center}
\begin{tabularx}{\linewidth}{|l|X|}\hline
\rowcolor{popblueL}\textbf{Question} & \textbf{Idées déclenchées sur le sujet} \\ \hline
Qui ? & Les anciens fonctionnaires, les militaires, les enseignants à la retraite, Meka le vieux nègre, l'État, le gouverneur, le Président. \\ \hline
Quoi ? & Médailles, pensions, distinctions honorifiques, prix d'excellence, trophées sportifs, simples remerciements. \\ \hline
Où ? & Cérémonies officielles (20 mai, fêtes nationales), écoles, stades, villages. \\ \hline
Quand ? & En fin de carrière, à la retraite, parfois à titre posthume. \\ \hline
Comment ? & Par décret, lors de cérémonies publiques, avec ou sans avantages matériels. \\ \hline
Pourquoi ? & Récompenser le sacrifice, encourager les jeunes, honorer la mémoire, mais aussi parfois : flatter, manipuler, se donner bonne conscience. \\ \hline
\end{tabularx}
\end{center}

\textbf{Mobiliser ses connaissances.} Un bon candidat puise ses idées dans cinq réservoirs :
\begin{itemize}
\item \textbf{L'expérience personnelle} : tu as vu des cérémonies de remise de médailles, des enseignants primés, des retraités de ton quartier.
\item \textbf{L'actualité camerounaise} : les distinctions de la fête nationale, les récompenses des Lions Indomptables, les prix d'excellence au concours.
\item \textbf{Les œuvres lues} : ici, \emph{Le Vieux nègre et la médaille} est une mine d'or — Meka reçoit une médaille pour ses terres données aux missionnaires, mais cette médaille lui apporte plus d'humiliations que de bonheur.
\item \textbf{L'histoire} : les anciens combattants, les décorations coloniales, les grands serviteurs de l'État camerounais.
\item \textbf{Les proverbes} : \og On ne reconnaît le bien que lorsqu'on l'a perdu \fg, \og La reconnaissance est la mémoire du cœur \fg.
\end{itemize}

\begin{popfigure}
\begin{tikzpicture}[mindmap, grow cyclic, every node/.style={concept, font=\small},
root concept/.append style={concept color=popblue, font=\small\bfseries},
level 1/.append style={level distance=3.4cm, sibling angle=72},
level 2/.append style={level distance=2.1cm, sibling angle=45}]
\node[root concept]{Faut-il récompenser les anciens serviteurs de l'État ?}
child[concept color=popgreen] { node {Justice et reconnaissance}
  child { node[font=\footnotesize] {Une vie de sacrifice mérite honneur} }
  child { node[font=\footnotesize] {Proverbe : la reconnaissance, mémoire du cœur} } }
child[concept color=poporange] { node {Utilité sociale}
  child { node[font=\footnotesize] {Encourager les jeunes à servir} }
  child { node[font=\footnotesize] {Exemple : prix d'excellence scolaires} } }
child[concept color=poppurple] { node {Limites et dérives}
  child { node[font=\footnotesize] {Médaille de Meka : honneur illusoire} }
  child { node[font=\footnotesize] {Distinctions parfois politiques} } }
child[concept color=poppink] { node {Formes de récompenses}
  child { node[font=\footnotesize] {Médailles, pensions, trophées} }
  child { node[font=\footnotesize] {Médailles sportives (Lions Indomptables)} } }
child[concept color=popgold] { node {Conditions d'un vrai hommage}
  child { node[font=\footnotesize] {Récompense sincère et accompagnée} }
  child { node[font=\footnotesize] {Pension décente plutôt que ruban} } };
\end{tikzpicture}
\legende{Carte mentale du brainstorming collectif autour du sujet : chaque branche principale regroupe les idées par famille ; les petites branches donnent les arguments et exemples.}
\end{popfigure}

\courssub{Le tri et le classement des idées}

\textbf{Intuition.} Après la récolte, le cultivateur trie. Après le brainstorming, le candidat doit faire de même : toutes les idées notées ne se valent pas.

\begin{methode}[Trier et classer ses idées]
\begin{enumerate}
\item \textbf{Évaluer la pertinence} : chaque idée répond-elle vraiment au sujet posé ? Une idée hors sujet, même brillante, doit être écartée.
\item \textbf{Éliminer les redondances} : deux idées qui disent la même chose n'en font qu'une.
\item \textbf{Regrouper par axes} : rassembler les idées qui vont dans le même sens (pour / contre / conditions, ou bien aspects sociaux / moraux / politiques). Ces regroupements préfigurent les parties du plan.
\item \textbf{Associer à chaque idée un argument et un exemple} : une idée sans illustration est une affirmation gratuite.
\end{enumerate}
\end{methode}

\begin{attention}[Le piège de l'accumulation]
L'erreur la plus fréquente est d'accumuler les idées sans les trier. \textbf{Dix idées désordonnées valent moins que quatre idées hiérarchisées.} Le correcteur ne note pas la quantité de ce que tu as pensé, mais la clarté de ce que tu as organisé. Une copie qui aligne douze arguments sans ordre donne l'impression d'un marché bruyant ; une copie qui en développe quatre, bien choisis et bien illustrés, donne l'impression d'une démonstration.
\end{attention}

Appliquons le tri à notre brainstorming :

\begin{center}
\begin{tabularx}{\linewidth}{|X|X|}\hline
\rowcolor{popblueL}\textbf{Idées retenues (avec justification)} & \textbf{Idées écartées (avec justification)} \\ \hline
La récompense est un devoir de justice envers qui a sacrifié sa vie (idée centrale, directement liée au sujet). & Les salaires des ministres (hors sujet : il ne s'agit pas de rémunération mais de récompense honorifique). \\ \hline
Elle encourage les jeunes générations à servir l'État (argument social fort, illustrable par les prix d'excellence). & La météo le jour des cérémonies (anecdote sans valeur argumentative). \\ \hline
La médaille de Meka montre qu'un honneur vide peut devenir une humiliation (exemple littéraire majeur, au cœur du sujet). & La liste de toutes les médailles existantes (inventaire inutile, redondant). \\ \hline
Une vraie récompense doit s'accompagner d'avantages concrets : pension, soins, considération (idée de synthèse qui prépare la troisième partie). & \og Les médailles sont belles \fg{} (jugement esthétique sans portée argumentative). \\ \hline
\end{tabularx}
\end{center}

\textbf{Choix des arguments et des exemples.} Pour chaque idée retenue, il faut un \cle{argument} (la raison qui la justifie) et un \cle{exemple} (le fait qui la prouve). Privilégie les situations camerounaises : les distinctions honorifiques du 20 mai, les récompenses scolaires (prix d'excellence, major de concours), les médailles sportives (Lions Indomptables, Françoise Mbango), et bien sûr la médaille de Meka dans le roman d'Oyono.

\courssub{La problématique : la boussole de la dissertation}

\textbf{Intuition.} Un voyageur qui part de Douala pour Ngaoundéré sans carte risque de se retrouver à Bafoussam. La problématique est la carte du devoir : une fois posée, chaque paragraphe sait où il va.

\begin{definition}[Problématique]
La \cle{problématique} est la question — ou l'ensemble de questions — qui traduit le problème posé par le sujet et qui guide tout le développement. Elle naît de la \emph{tension} repérée lors de l'analyse du sujet : un sujet de dissertation contient toujours un désaccord possible, un doute, une opposition (par exemple : récompenser est-il un devoir ou une hypocrisie ?). La problématique transforme cette tension en question explicite.
\end{definition}

\begin{methode}[Formuler une problématique en trois étapes]
\begin{enumerate}
\item \textbf{Repérer la tension du sujet} : quels sont les deux termes qui semblent s'opposer ? (Ici : \emph{récompenser} / mais la récompense peut être vaine ou injuste.)
\item \textbf{Transformer cette tension en question ouverte} : une question qui ne se résout pas par \og oui \fg{} ou \og non \fg{} sec, mais qui appelle une discussion. Utiliser des formes comme \og Dans quelle mesure... ? \fg, \og ... suffit-elle à... ? \fg, \og Faut-il... ou plutôt... ? \fg.
\item \textbf{Vérifier qu'elle appelle un plan} : la problématique doit pouvoir se répondre en deux ou trois parties (thèse / antithèse, ou thèse / antithèse / synthèse). Si elle n'appelle qu'une seule réponse évidente, elle est mauvaise.
\end{enumerate}
\end{methode}

\begin{exemplebox}[Problématiques sur le sujet de la médaille de Meka]
\begin{itemize}
\item \og Une récompense officielle suffit-elle à honorer une vie de sacrifice ? \fg
\item \og La médaille offerte à un vieux serviteur est-elle un hommage sincère ou une manière de se débarrasser de lui à peu de frais ? \fg
\item \og Dans quelle mesure les distinctions honorifiques rendent-elles justice aux anciens serviteurs de l'État ? \fg
\end{itemize}
Chacune de ces questions couvre tout le sujet et annonce un plan : oui, la récompense est nécessaire (partie 1) ; mais elle peut être illusoire ou hypocrite (partie 2) ; à condition d'être sincère et accompagnée, elle honore vraiment (partie 3).
\end{exemplebox}

\begin{propriete}[Critères de validité d'une problématique]
Une bonne problématique satisfait trois conditions : (1) elle \textbf{couvre tout le sujet} (aucun mot-clé n'est oublié) ; (2) elle est \textbf{ouverte} (elle appelle une discussion, pas une évidence) ; (3) elle \textbf{annonce le plan} (la réponse s'organise naturellement en deux ou trois parties). Au brouillon, vérifie toujours ces trois points avant de rédiger l'introduction.
\end{propriete}

\begin{exoresolu}[Du brainstorming à la problématique]
Énoncé : à partir du sujet \og Faut-il récompenser les anciens serviteurs de l'État ? \fg, un élève a noté les idées suivantes : (a) les retraités ont sacrifié leur jeunesse ; (b) Meka est humilié après sa médaille ; (c) les jeunes ont besoin de modèles ; (d) certaines médailles sont distribuées par favoritisme ; (e) une pension décente vaut mieux qu'un ruban ; (f) les stades sont pleins le 20 mai. Trie ces idées et formule une problématique.
\tcblower
\textbf{Solution détaillée.} \emph{Tri} : l'idée (f) est écartée (constat sans valeur argumentative). Les idées (a) et (c) se regroupent : la récompense est un devoir de justice et un exemple pour la jeunesse (axe 1). Les idées (b) et (d) se regroupent : la récompense peut être illusoire ou injuste (axe 2). L'idée (e) ouvre une synthèse : la vraie récompense doit être concrète (axe 3). \emph{Problématique} : la tension oppose la nécessité de récompenser et le risque d'une récompense vide. D'où la question : \og Une distinction honorifique suffit-elle à rendre justice à celui qui a donné sa vie à l'État ? \fg. Cette question couvre tout le sujet et annonce un plan en trois parties : nécessité de la récompense, limites et dérives, conditions d'un hommage véritable.
\end{exoresolu}

\begin{savaistu}[La problématique au Baccalauréat technique]
Au Baccalauréat technique, la problématique figure obligatoirement dans l'introduction, juste avant l'annonce du plan. Les correcteurs la repèrent immédiatement : elle conditionne la note de \emph{cohérence} de la copie. Une problématique absente ou fausse fait perdre des points même si le développement est riche, car le correcteur ne retrouve pas de fil directeur. À l'inverse, une problématique claire et une annonce de plan fidèle donnent à ta copie une architecture visible — et les copies bien architecturées sont toujours mieux notées.
\end{savaistu}

\begin{exercice}[À toi de jouer]
Sujet : \og Le sport est-il la seule voie de réussite pour la jeunesse camerounaise ? \fg
\begin{enumerate}
\item Conduis un brainstorming individuel de cinq minutes, puis complète-le par le questionnement (qui ? quoi ? où ? quand ? comment ? pourquoi ?).
\item Trie tes idées : écarte deux idées non pertinentes en justifiant, regroupe les autres en deux ou trois axes.
\item Formule une problématique et vérifie qu'elle satisfait les trois critères de validité.
\end{enumerate}
\end{exercice}

\begin{corrige}[Exercice 1]
\emph{Éléments de réponse.} Brainstorming possible : Lions Indomptables, Samuel Eto'o, médecins et ingénieurs, prix d'excellence, chômage des diplômés, écoles et universités, artistes musiciens, agriculture, concours de la fonction publique. Idées écartables : par exemple \og les stades sont beaux \fg{} (constat sans argument) et \og je n'aime pas le football \fg{} (goût personnel, pas un argument). Regroupement : axe 1 — le sport, voie réelle de réussite (Eto'o, Mbango : gloire et richesse) ; axe 2 — mais d'autres voies existent (études, concours, agriculture, arts) et le sport seul est risqué (peu de carrières réussies, blessures) ; axe 3 — la réussite suppose d'abord le travail et la formation, quel que soit le domaine. Problématique possible : \og Le sport suffit-il à assurer la réussite de la jeunesse, ou celle-ci exige-t-elle d'autres voies et d'autres valeurs ? \fg{} Elle couvre le sujet, elle est ouverte, elle annonce un plan en trois parties.
\end{corrige}

\begin{aretenir}[L'essentiel de la section]
\begin{itemize}
\item La recherche des idées se fait par \cle{brainstorming} (individuel puis collectif), aidé du \cle{questionnement} : qui ? quoi ? où ? quand ? comment ? pourquoi ?
\item On puise dans cinq réservoirs : expérience, actualité camerounaise, œuvres lues (dont \emph{Le Vieux nègre et la médaille}), histoire, proverbes.
\item Les idées doivent être \cle{triées} (pertinence, redondances) puis \cle{regroupées par axes}, chaque idée étant doublée d'un argument et d'un exemple.
\item La \cle{problématique} est la question née de la tension du sujet ; elle doit couvrir tout le sujet, être ouverte et annoncer le plan.
\end{itemize}
\end{aretenir}

\coursec{Élaboration et production du plan détaillé}

Dans la section précédente, nous avons appris à analyser le sujet, à mobiliser nos idées et à formuler une problématique claire. Mais une problématique, aussi pertinente soit-elle, reste une simple question si elle ne débouche pas sur une organisation. C'est précisément le rôle du plan détaillé : transformer un amas d'idées en architecture rigoureuse. C'est l'étape décisive qui sépare la copie brouillonne de la copie remarquable.

\courssub{Qu'est-ce qu'un plan détaillé ?}

\begin{definition}[Le plan détaillé]
Le \cle{plan détaillé} est l'organisation écrite et hiérarchisée des idées d'une dissertation. Il présente, dans l'ordre où elles seront traitées, les \cle{parties} (I, II, III), les \cle{sous-parties} (A, B), les \cle{arguments} (1, 2) et les \cle{exemples} qui illustreront chaque argument. Il se rédige au brouillon avant toute rédaction au propre.
\end{definition}

L'intuition est simple : personne ne construit une maison sans plan. De même, personne ne rédige une dissertation solide sans avoir d'abord dessiné sa charpente. Le plan détaillé est cette charpente : il garantit que chaque idée trouvera sa place, que rien ne sera oublié et que l'ensemble tiendra debout.

\begin{propriete}[Les trois qualités d'un bon plan]
Un bon plan détaillé est :
\begin{itemize}
\item \cle{équilibré} : les parties ont un poids comparable (on ne consacre pas dix lignes à la partie I et trois pages à la partie II) ;
\item \cle{progressif} : l'ordre des parties obéit à une logique (du plus simple au plus complexe, du constat aux solutions, de la thèse à la synthèse) ;
\item \cle{exhaustif} : il couvre la totalité du sujet, sans hors-sujet ni oubli.
\end{itemize}
\end{propriete}

\begin{popfigure}
\begin{center}
\begin{tikzpicture}[scale=0.9, every node/.style={transform shape}]
% Coordonnées de la pyramide
\coordinate (baseL) at (0,0);
\coordinate (baseR) at (6,0);
\coordinate (top) at (3,4.5);
% Découpes horizontales
\coordinate (intBaseL) at (0.6,1.3);
\coordinate (intBaseR) at (5.4,1.3);
\coordinate (devBaseL) at (1.2,2.6);
\coordinate (devBaseR) at (4.8,2.6);

% Remplissage des couches (bas -> haut : Conclusion, Développement, Introduction)
\fill[popblueL] (baseL) -- (baseR) -- (intBaseR) -- (intBaseL) -- cycle;
\fill[poporangeL] (intBaseL) -- (intBaseR) -- (devBaseR) -- (devBaseL) -- cycle;
\fill[popgreenL] (devBaseL) -- (devBaseR) -- (top) -- cycle;

% Contours
\draw[popdark, thick] (baseL) -- (baseR) -- (top) -- cycle;
\draw[popdark] (intBaseL) -- (intBaseR);
\draw[popdark] (devBaseL) -- (devBaseR);

% Légendes internes
\node[popdark, font=\small\bfseries, align=center] at (3,0.65) {Conclusion\\ Bilan + Réponse + Ouverture};
\node[popdark, font=\small\bfseries, align=center] at (3,2.0) {Développement\\ 2 ou 3 parties équilibrées};
\node[popdark, font=\small\bfseries, align=center, text width=3cm] at (3,3.6) {Introduction\\ Amorce, Sujet, Problématique, Annonce du plan};

% Annotations à droite
\node[popink, font=\scriptsize, align=left, text width=4.2cm, anchor=west] (annInt) at (6.5,3.6) {\textbf{Introduction} : accroche, définition des termes, problématique, annonce du plan.};
\node[popink, font=\scriptsize, align=left, text width=4.2cm, anchor=west] (annDev) at (6.5,2.0) {\textbf{Développement} : chaque paragraphe = 1 argument + explication + exemple précis. Connecteurs logiques.};
\node[popink, font=\scriptsize, align=left, text width=4.2cm, anchor=west] (annConcl) at (6.5,0.65) {\textbf{Conclusion} : bilan synthétique, réponse tranchée à la problématique, ouverture.};

% Flèches depuis les arêtes droites de la pyramide
\draw[popink, ->, thick] (5.4,1.3) -- (annConcl.west);
\draw[popink, ->, thick] (4.8,2.6) -- (annDev.west);
\draw[popink, ->, thick] (3.6,3.8) -- (annInt.west);
\end{tikzpicture}
\legende{L'architecture de la dissertation : l'introduction ouvre le propos, le développement (partie la plus volumineuse) déploie les arguments, la conclusion resserre et clôt.}
\end{center}
\end{popfigure}

\courssub{Les trois types de plans}

Il n'existe pas un plan unique, mais trois grands modèles. Le choix dépend du type de sujet posé.

\begin{definition}[Le plan dialectique]
Le \cle{plan dialectique} est un plan en trois temps : la \cle{thèse} (une première position), l'\cle{antithèse} (la position opposée ou les limites de la thèse) et la \cle{synthèse} (le dépassement des deux positions). Il est adapté aux sujets de \emph{discussion}, c'est-à-dire aux sujets qui opposent deux points de vue ou qui interrogent une opinion.
\end{definition}

Le \cle{plan thématique} examine successivement les différents \emph{aspects} d'une notion, d'un personnage ou d'un phénomène (par exemple : les causes, puis les manifestations, puis les valeurs défendues). Il convient aux sujets descriptifs ou d'appréciation portant sur une œuvre ou un thème.

Le \cle{plan analytique} (ou explicatif) suit la démarche de l'enquête : \emph{constat} (on décrit le problème), \emph{causes} (on explique pourquoi), \emph{conséquences} ou \emph{solutions} (on envisage la suite). Il convient aux sujets portant sur un problème de société.

\begin{center}
\begin{tabularx}{\linewidth}{|l|X|X|X|}
\hline
\rowcolor{popblueL} \textbf{Critère} & \textbf{Plan dialectique} & \textbf{Plan thématique} & \textbf{Plan analytique} \\
\hline
Structure & Thèse / Antithèse / Synthèse & Aspect 1 / Aspect 2 / Aspect 3 & Constat / Causes / Conséquences-Solutions \\
\hline
Type de sujet & Sujet de discussion, opinion à débattre & Sujet sur une œuvre, un thème, un personnage & Sujet sur un problème de société \\
\hline
Formulation typique & « Faut-il... ? », « Peut-on dire que... ? » & « Montrez que... », « Quelle image de... ? » & « Que pensez-vous de... ? » (phénomène social) \\
\hline
Exemple de sujet & « La médaille récompense-t-elle toujours le mérite ? » & « La satire coloniale dans \emph{Le Vieux nègre et la médaille} » & « L'aliénation des élites africaines colonisées » \\
\hline
\end{tabularx}
\end{center}

\begin{attention}[Le critère de choix]
Ne choisis jamais un plan « par habitude ». Interroge d'abord le sujet : s'il contient une opposition ou invite à discuter une opinion, pense dialectique ; s'il demande de décrire ou d'apprécier, pense thématique ; s'il présente un problème à expliquer, pense analytique. Un plan dialectique plaqué sur un sujet thématique produit un hors-sujet partiel.
\end{attention}

\courssub{La structure interne de la dissertation}

\textbf{L'introduction} comporte quatre moments obligatoires :
\begin{enumerate}
\item l'\cle{amorce} (ou accroche) : une phrase générale qui situe le thème ;
\item la \cle{présentation du sujet} : on reformule et on délimite la question ;
\item la \cle{problématique} : la question directrice, souvent introduite par « on peut se demander si... » ;
\item l'\cle{annonce du plan} : on présente les grandes parties, sans les détailler.
\end{enumerate}

\textbf{Le développement} est composé de parties équilibrées. Chaque paragraphe suit la loi des trois temps : un \cle{argument} (l'idée), une \cle{explication} (on développe, on justifie) et un \cle{exemple} (on illustre par une référence précise : œuvre, fait, personnage). Les paragraphes sont reliés par des \cle{connecteurs logiques} (d'abord, en effet, cependant, par conséquent...).

\textbf{La conclusion} comporte trois moments : le \cle{bilan} (on résume le chemin parcouru), la \cle{réponse à la problématique} (on tranche clairement) et l'\cle{ouverture} (une question ou une perspective nouvelle, sans idée neuve développée).

\begin{methode}[Produire un plan détaillé en six étapes]
\begin{enumerate}
\item Reprendre la \cle{problématique} formulée à l'étape précédente : c'est elle qui commande tout le plan.
\item Déterminer les \cle{grandes parties} (2 ou 3) en choisissant le type de plan adapté au sujet.
\item Diviser chaque partie en \cle{sous-parties} (A, B) qui précisent l'angle de traitement.
\item Sous chaque sous-partie, formuler les \cle{arguments} (1, 2) sous forme de phrases complètes.
\item Associer à chaque argument un \cle{exemple} précis et vérifiable.
\item \cle{Vérifier la cohérence} : équilibre des parties, progression logique, couverture totale du sujet, adéquation avec la problématique.
\end{enumerate}
\end{methode}

\begin{savaistu}[Le temps du brouillon]
Au Baccalauréat technique, les correcteurs et les expériences de classe concordent : l'élaboration du plan détaillé au brouillon prend entre 20 et 30 minutes sur les 4 heures de l'épreuve. C'est un investissement, pas une perte de temps : une copie construite sur un plan solide se rédige ensuite presque sans hésitation, tandis qu'une copie commencée « au fil de la plume » s'enlise et se contredit.
\end{savaistu}

\courssub{Exemple complet : un plan détaillé sur \emph{Le Vieux nègre et la médaille}}

\begin{exoresolu}[Plan détaillé rédigé]
\textbf{Sujet :} « Dans \emph{Le Vieux nègre et la médaille} de Ferdinand Oyono, la médaille offerte à Meka est-elle une récompense ou un instrument de domination ? »

\textbf{Problématique :} la médaille, symbole apparent de reconnaissance, ne cache-t-elle pas un instrument de la domination coloniale ?
\tcblower
\textbf{Type de plan retenu :} dialectique (le sujet oppose deux lectures : récompense / instrument).

\textbf{Introduction (au brouillon) :} amorce sur les décorations coloniales ; présentation de l'œuvre et de la scène de remise ; problématique ; annonce : « Nous verrons d'abord que la médaille apparaît comme une récompense, puis qu'elle fonctionne comme un instrument de domination, enfin qu'Oyono propose une prise de conscience qui dépasse cette alternative. »

\textbf{I. La médaille comme récompense apparente}
\begin{itemize}
\item \textbf{A. Une reconnaissance officielle}
\begin{enumerate}
\item Meka a « donné » ses terres à la mission et ses fils à la guerre : la médaille couronne une vie de sacrifices. \emph{Exemple : le discours du commandant lors de la cérémonie.}
\item La fête organisée à Doum célèbre le vieil homme : honneur social. \emph{Exemple : l'émulation et la jalousie des villageois.}
\end{enumerate}
\item \textbf{B. La fierté sincère de Meka}
\begin{enumerate}
\item Meka croit à cette reconnaissance : il se sent enfin « quelqu'un ». \emph{Exemple : sa joie, ses préparatifs, ses rêves de prestige.}
\end{enumerate}
\end{itemize}

\textbf{II. La médaille comme instrument de domination}
\begin{itemize}
\item \textbf{A. Une récompense dérisoire}
\begin{enumerate}
\item Un simple bout de métal contre des terres et des vies : l'échange est inégal. \emph{Exemple : l'ironie du narrateur sur la « grande » médaille.}
\item La médaille ne change rien à la condition de Meka : il reste pauvre et méprisé. \emph{Exemple : son arrestation et sa nuit au poste le soir même de la fête.}
\end{enumerate}
\item \textbf{B. Un outil de propagande coloniale}
\begin{enumerate}
\item La cérémonie sert l'image de la colonisation « généreuse ». \emph{Exemple : les discours officiels, la présence des notables blancs.}
\item Elle divise les Africains en créant de faux privilégiés. \emph{Exemple : les rivalités entre villageois autour de Meka.}
\end{enumerate}
\end{itemize}

\textbf{III. Le dépassement : la prise de conscience}
\begin{itemize}
\item \textbf{A. La désillusion de Meka}
\begin{enumerate}
\item L'humiliation au poste de police détruit l'illusion. \emph{Exemple : la médaille ne le protège pas des coups.}
\item Meka comprend la vanité de la récompense. \emph{Exemple : ses réflexions amères après sa libération.}
\end{enumerate}
\item \textbf{B. La satire d'Oyono}
\begin{enumerate}
\item L'ironie dénonce l'hypocrisie du système colonial. \emph{Exemple : le comique amer des scènes officielles.}
\item La solidarité villageoise comme vraie valeur. \emph{Exemple : l'entraide des gens de Doum.}
\end{enumerate}
\end{itemize}

\textbf{Conclusion (au brouillon) :} bilan — la médaille est une récompense en apparence seulement ; réponse — c'est un instrument de domination que la prise de conscience de Meka révèle ; ouverture — la littérature africaine comme arme de dénonciation.
\end{exoresolu}

\begin{exemplebox}[Ce qui est en jeu au Baccalauréat technique]
La dissertation est notée sur 20. La \cle{cohérence du plan} (respect du sujet, logique des parties, équilibre) et la \cle{pertinence des arguments et des exemples} pèsent lourd dans le barème, souvent davantage que la seule correction de la langue. Un plan déséquilibré ou contradictoire avec l'annonce faite en introduction fait perdre plusieurs points, même si le style est bon.
\end{exemplebox}

\begin{attention}[L'annonce du plan engage toute la copie]
Erreur fréquente et grave : annoncer trois parties dans l'introduction et n'en traiter que deux — ou traiter les parties dans un ordre différent de celui annoncé. Le correcteur lit l'introduction comme un contrat : toute partie annoncée et non traitée est sanctionnée comme un manquement. Rédige l'annonce du plan \emph{après} avoir stabilisé ton plan au brouillon, et vérifie en fin de copie que le contrat est respecté.
\end{attention}

\courssub{Auto-évaluation du plan détaillé}

Avant de passer à la rédaction, soumets ton plan à une grille de contrôle. Cette relecture critique prend cinq minutes et peut sauver la copie.

\begin{center}
\begin{tabularx}{\linewidth}{|l|X|X|}
\hline
\rowcolor{popblueL} \textbf{Critère} & \textbf{Question à se poser} & \textbf{Si la réponse est non...} \\
\hline
Respect du sujet & Chaque partie répond-elle au sujet posé, sans hors-sujet ? & Supprimer ou reformuler la partie fautive. \\
\hline
Problématique & Le plan répond-il exactement à la problématique formulée ? & Ajuster la problématique ou les parties. \\
\hline
Logique & L'ordre des parties est-il progressif et justifié ? & Réordonner (constat avant causes, thèse avant synthèse). \\
\hline
Équilibre & Les parties ont-elles un poids comparable ? & Fusionner deux sous-parties ou en subdiviser une. \\
\hline
Exemples & Chaque argument est-il illustré par un exemple précis ? & Chercher une référence dans l'œuvre ou la culture. \\
\hline
Annonce & L'introduction annonce-t-elle exactement ce que la copie traitera ? & Réécrire l'annonce du plan. \\
\hline
\end{tabularx}
\end{center}

\begin{exercice}[À toi de jouer]
Sujet : « Le rire, dans \emph{Le Vieux nègre et la médaille}, n'est-il qu'un divertissement ? »
\begin{enumerate}
\item Identifie le type de plan le plus adapté et justifie ton choix.
\item Formule une problématique.
\item Rédige le squelette complet du plan détaillé (I. A. 1. 2. B. ... II. ...), avec un exemple précis pour chaque argument.
\item Contrôle ton plan à l'aide de la grille d'auto-évaluation ci-dessus.
\end{enumerate}
\end{exercice}

\begin{corrige}[Exercice 1]
\textbf{1.} Le sujet invite à discuter une opinion restrictive (« n'est-il \emph{que}... ») : le \cle{plan dialectique} s'impose (thèse : le rire divertit ; antithèse : il dénonce ; synthèse : un rire qui éveille les consciences).

\textbf{2.} Problématique possible : le comique d'Oyono se réduit-il au divertissement, ou constitue-t-il une arme de critique sociale ?

\textbf{3.} Squelette attendu :
\textbf{I.} Le rire comme divertissement — A. les scènes comiques (ex. : les préparatifs désordonnés de la fête) ; B. les personnages ridicules (ex. : le catéchiste vaniteux).
\textbf{II.} Le rire comme dénonciation — A. la satire des colons (ex. : les discours pompeux du commandant) ; B. la satire des Africains aliénés (ex. : les notables qui singent les Blancs).
\textbf{III.} Un rire qui éveille — A. le rire amer après l'arrestation de Meka ; B. la solidarité villageoise, vraie leçon du roman.

\textbf{4.} Vérification : sujet respecté (le rire est bien le centre), progression (du plaisir à la critique), équilibre (deux sous-parties par partie), exemples présents.
\end{corrige}

\begin{aretenir}[L'essentiel]
Le plan détaillé est la charpente de la dissertation : parties, sous-parties, arguments, exemples. On choisit le plan \cle{dialectique} (discussion), \cle{thématique} (aspects) ou \cle{analytique} (problème) selon le sujet. Un bon plan est équilibré, progressif et exhaustif ; l'annonce du plan engage toute la copie. Consacrer 20 à 30 minutes au brouillon, c'est sécuriser l'ensemble de la rédaction.
\end{aretenir}

\coursec{Langue au service de la dissertation : référent et substituts, sigles et emprunts}

Dans la section précédente, nous avons appris à construire un \cle{plan détaillé} : problématique claire, parties équilibrées, arguments et exemples hiérarchisés. Mais un bon plan ne suffit pas : encore faut-il le \emph{rédiger} dans une langue claire, fluide et correcte. Or deux défauts ruinent souvent les copies de Terminale : les \cle{répétitions} maladroites (le même mot revient dix fois) et l'emploi abusif de \cle{sigles} ou d'\cle{emprunts} familiers. Cette dernière section met donc la langue française au service de ta dissertation : tu apprendras à repérer le \cle{référent} et ses \cle{substituts}, puis à manier avec discernement sigles et emprunts, notamment dans le français en usage au Cameroun.

\courssub{Le référent : ce dont on parle}

Commençons par une observation simple. Lis cette phrase : « Meka reçut sa médaille. Le vieux nègre la fixa sur sa poitrine. Il sourit. Le vieillard était fier. » Quatre expressions différentes — \emph{Meka}, \emph{le vieux nègre}, \emph{il}, \emph{le vieillard} — et pourtant tu n'as jamais eu l'impression qu'on parlait de quatre personnes. Pourquoi ? Parce que ces quatre mots désignent une seule et même réalité : le héros du roman de Ferdinand Oyono.

\begin{definition}[Le référent]
Le \cle{référent} est la réalité extralinguistique — une personne, un objet, un animal, une notion — désignée par un mot ou un groupe de mots dans un texte. Dans \emph{Le Vieux nègre et la médaille}, le mot \emph{Meka} a pour référent le personnage principal du roman, un vieillard camerounais décoré par l'administration coloniale.
\end{definition}

Il faut bien distinguer le \emph{mot} (ce qui est écrit sur la page) et le \emph{référent} (la réalité désignée). Plusieurs mots différents peuvent avoir le même référent : c'est précisément ce qui rend possible la variété d'expression.

\begin{attention}[Ne pas confondre référent et sens]
Le référent n'est pas le dictionnaire ! \emph{Le vieillard} et \emph{Meka} n'ont pas le même sens dans le dictionnaire (l'un est un nom commun, l'autre un nom propre), mais ils ont le \textbf{même référent} dans le roman. Le référent dépend toujours du \cle{contexte} : dans un autre texte, \emph{le vieillard} désignera quelqu'un d'autre.
\end{attention}

\courssub{Les substituts du référent}

Pour éviter de répéter sans cesse le même nom, la langue dispose de \cle{substituts} : des mots qui « remplacent » le nom du référent tout en continuant à le désigner.

\begin{definition}[Chaîne de référence]
Une \cle{chaîne de référence} est l'ensemble des mots différents qui désignent le même référent dans un texte. Elle comprend en général un \emph{terme introducteur} (la première mention, souvent un nom propre ou un groupe nominal complet) suivi de \emph{reprises} (pronoms, synonymes, reprises nominales).
\end{definition}

Recensons les principaux types de substituts.

\begin{center}
\begin{tabularx}{\linewidth}{|l|X|X|}
\hline
\rowcolor{popblueL} \textbf{Type de substitut} & \textbf{Définition} & \textbf{Exemple (référent : Meka)} \\
\hline
Pronom personnel & Remplace un nom déjà mentionné (il, elle, lui, le, en, y...) & \emph{Il} attendit la médaille. \\
\hline
Pronom démonstratif & Reprend en montrant (celui, celle, ce, cela) & \emph{Celui} qui avait tant attendu... \\
\hline
Reprise nominale & Nom commun précédé d'un déterminant, qui décrit le référent & \emph{le vieux nègre}, \emph{le vieillard} \\
\hline
Synonyme & Mot de sens très proche & \emph{le vieux} pour \emph{le vieillard} \\
\hline
Hyperonyme & Terme plus général englobant le référent & \emph{l'homme}, \emph{le personnage} \\
\hline
\end{tabularx}
\end{center}

\begin{popfigure}
\begin{center}
\begin{tikzpicture}[
  node distance=1.1cm,
  ref/.style={rectangle, rounded corners, draw=popblue, fill=popblueL, thick, minimum width=2.6cm, minimum height=0.9cm, font=\small},
  lab/.style={font=\scriptsize\itshape, text=popmuted}
]
\node[ref, fill=popgoldL, draw=popgold] (meka) {\textbf{Meka}};
\node[ref, right=of meka] (negre) {le vieux nègre};
\node[ref, right=of negre] (il) {il};
\node[ref, right=of il] (vieillard) {le vieillard};
\node[lab, below=0.15cm of meka] {terme introducteur};
\node[lab, below=0.15cm of negre] {reprise nominale};
\node[lab, below=0.15cm of il] {pronom personnel};
\node[lab, below=0.15cm of vieillard] {reprise nominale};
\draw[->, thick, popblue] (meka) -- (negre);
\draw[->, thick, popblue] (negre) -- (il);
\draw[->, thick, popblue] (il) -- (vieillard);
\node[draw=poppurple, dashed, thick, ellipse, fit=(meka)(vieillard), inner sep=0.35cm, label={[poppurple, font=\small\bfseries]above:un seul référent : le héros du roman}] {};
\end{tikzpicture}
\end{center}
\legende{Chaîne de référence extraite de l'univers du \emph{Vieux nègre et la médaille} : quatre expressions différentes (1. Meka, 2. le vieux nègre, 3. il, 4. le vieillard) renvoient à un seul et même référent, le héros du roman.}
\end{popfigure}

\begin{aretenir}[Rôle des substituts]
Les substituts remplissent deux fonctions essentielles dans ta dissertation :
\begin{itemize}
\item ils \textbf{évitent les répétitions} qui alourdissent le style ;
\item ils assurent la \cle{cohésion} du texte : le lecteur suit le même personnage, le même objet ou la même notion d'une phrase à l'autre sans se perdre.
\end{itemize}
Un texte bien « enchaîné » est un texte agréable à lire et mieux noté.
\end{aretenir}

\begin{methode}[Varier les reprises sans perdre la clarté]
Pour éviter les répétitions dans ta copie :
\begin{enumerate}
\item Repère le mot qui se répète (par exemple \emph{la médaille} quatre fois dans un paragraphe).
\item À la première occurrence, écris le nom complet : c'est le terme introducteur.
\item Aux occurrences suivantes, alterne les substituts : pronom (\emph{elle}, \emph{la}), reprise nominale (\emph{cette distinction}, \emph{le ruban}), synonyme ou hyperonyme (\emph{cette récompense}).
\item Vérifie toujours que chaque substitut renvoie \textbf{sans ambiguïté} au bon référent : si deux personnages masculins sont présents, le pronom \emph{il} peut créer une confusion ; reprends alors le nom.
\end{enumerate}
Règle d'or : la variété ne doit jamais se faire au détriment de la \cle{clarté}.
\end{methode}

\begin{exoresolu}[Repérage d'une chaîne de référence]
\textbf{Énoncé.} Dans le passage suivant, inspiré de l'univers du \emph{Vieux nègre et la médaille}, souligne tous les éléments de la chaîne de référence dont le référent est Meka, puis classe chaque reprise selon son type : « Meka attendait depuis longtemps cette médaille. Le vieux nègre avait donné sa terre et ses fils à la guerre. Il croyait que la France le récompenserait enfin. Le vieillard, ce jour-là, portait son plus beau pagne. »
\tcblower
\textbf{Solution détaillée.}
\begin{itemize}
\item \emph{Meka} : terme introducteur (nom propre).
\item \emph{Le vieux nègre} : reprise nominale (nom commun + déterminant).
\item \emph{Il} et \emph{le} (dans « la France \textbf{le} récompenserait ») : pronoms personnels.
\item \emph{Le vieillard} : reprise nominale, proche du synonyme de « vieux nègre ».
\end{itemize}
La chaîne compte donc cinq maillons pour un seul référent. On remarque qu'Oyono (et le rédacteur qui l'imite) évite ainsi la lourdeur de répéter « Meka » cinq fois, tout en donnant des informations nouvelles : la reprise nominale \emph{le vieux nègre} insiste sur l'âge et la condition coloniale du personnage.
\end{exoresolu}

\courssub{Les sigles : écrire court sans se faire mal comprendre}

Ta vie quotidienne au Cameroun est pleine de sigles : tu passes le probatoire organisé sous la tutelle du MINESEC, tu écoutes la CRTV, tes parents paient peut-être des factures à la SNH ou à ENEO. Ces abréviations sont pratiques, mais leur usage à l'écrit obéit à des règles précises.

\begin{definition}[Le sigle]
Un \cle{sigle} est une abréviation constituée des initiales d'un groupe de mots. Il se prononce de deux façons :
\begin{itemize}
\item \textbf{lettre par lettre} : CRTV se dit « cé-èr-té-vé », SNH se dit « esse-ène-ache » ;
\item \textbf{comme un mot ordinaire} : on parle alors d'\cle{acronyme}, comme MINESEC (« minéssec ») ou CNU (« cé-nu » peut aussi se lire lettre par lettre selon les usages).
\end{itemize}
Le sigle s'écrit en \textbf{majuscules}, généralement \textbf{sans points} entre les lettres et sans accent.
\end{definition}

\begin{propriete}[Règles d'écriture et de lecture des sigles]
\begin{enumerate}
\item Le sigle s'écrit en capitales : MINESEC, GCE, CRTV.
\item Dans un texte rédigé (dissertation, résumé, commentaire), on écrit le nom \textbf{en entier à la première occurrence}, suivi du sigle entre parenthèses : « le ministère des Enseignements secondaires (MINESEC) » ; on peut ensuite utiliser le sigle seul.
\item On ne crée jamais de sigle de sa propre invention dans une copie.
\item Un sigle inconnu du lecteur et non expliqué est une faute de communication.
\end{enumerate}
\end{propriete}

\courssub{Les emprunts : quand le français voyage}

\begin{definition}[L'emprunt]
Un \cle{emprunt} est un mot qu'une langue prend à une langue étrangère et intègre à son vocabulaire. Le français a emprunté à l'anglais (\emph{week-end}, \emph{business}, \emph{football}), à l'arabe (\emph{algèbre}, \emph{sucre}), et le français en usage au Cameroun emprunte largement aux langues nationales : \emph{ndolè} (plat de feuilles), \emph{miondo} (bâton de manioc), \emph{mboa} (habitant du pays, en duala), \emph{sawa}...
\end{definition}

\begin{savaistu}[Le français camerounais, une richesse]
Le français parlé au Cameroun intègre des centaines d'emprunts aux langues nationales : \emph{ndolè}, \emph{miondo}, \emph{eru}, \emph{koki}, \emph{mbol}... Ces mots traduisent des réalités culturelles que le français de France ne nomme pas : ils constituent une richesse identitaire, étudiée par les linguistes sous le nom de « camfranglais » ou de français camerounais quand ils se mêlent à l'anglais et au pidgin. Mais cette richesse a ses lieux : conversation, chanson, littérature orale — non la copie d'examen, où l'on attend un français normé et compréhensible par tout correcteur.
\end{savaistu}

\begin{center}
\begin{tabularx}{\linewidth}{|l|X|X|}
\hline
\rowcolor{popgreenL} \textbf{Catégorie} & \textbf{Exemple} & \textbf{Origine / observation} \\
\hline
Sigle & MINESEC & Ministère des Enseignements secondaires (Cameroun) \\
\hline
Sigle & GCE & \emph{General Certificate of Education} (examen anglophone) \\
\hline
Sigle & CRTV & Cameroon Radio Television \\
\hline
Emprunt & week-end & anglais, intégré au français courant \\
\hline
Emprunt & business & anglais ; en copie, préférer \emph{affaires}, \emph{commerce} \\
\hline
Emprunt & ndolè & langues du littoral ; réalité culturelle camerounaise \\
\hline
Emprunt & miondo & langue ewondo ; mets national \\
\hline
\end{tabularx}
\end{center}

\begin{attention}[Sigles et emprunts dans une copie d'examen]
Deux erreurs fréquentes coûtent des points au Probatoire et au Baccalauréat technique :
\begin{itemize}
\item Employer un \textbf{sigle non expliqué} : le correcteur ne doit pas deviner. Écris toujours le nom complet à la première occurrence.
\item Employer un \textbf{emprunt familier ou local} (\emph{business}, \emph{ndolè}, \emph{mboa}) dans une dissertation : c'est un hors-registre. Remplace-le par un terme de français standard (\emph{commerce}, \emph{plat traditionnel}, \emph{compatriote}), sauf si tu cites volontairement une réalité culturelle — auquel cas tu peux gloser le mot : « le \emph{ndolè}, mets traditionnel du littoral camerounais ».
\end{itemize}
\end{attention}

\courssub{Application : réécrire un paragraphe}

\begin{exoresolu}[Améliorer les reprises et corriger sigles et emprunts]
\textbf{Énoncé.} Réécris le paragraphe suivant en (a) supprimant les répétitions grâce à des substituts variés, (b) corrigeant l'usage du sigle et des emprunts : « Meka veut faire business avec sa médaille. Meka pense que la médaille va enrichir Meka. La médaille, pour Meka, c'est comme un diplôme du MINESEC. Meka invite tout le village à manger le ndolè. »
\tcblower
\textbf{Solution détaillée.} « Meka veut tirer profit de sa médaille : il pense que cette distinction va l'enrichir. Pour le vieillard, elle vaut un diplôme du ministère des Enseignements secondaires (MINESEC). C'est pourquoi il invite tout le village à partager un festin traditionnel. »
Corrections effectuées :
\begin{itemize}
\item \emph{business} $\rightarrow$ \emph{tirer profit} (emprunt familier supprimé) ;
\item répétitions de \emph{Meka} et de \emph{la médaille} $\rightarrow$ pronoms (\emph{il}, \emph{elle}), reprises nominales (\emph{le vieillard}, \emph{cette distinction}) ;
\item \emph{MINESEC} $\rightarrow$ nom développé à la première occurrence, sigle entre parenthèses ;
\item \emph{ndolè} $\rightarrow$ \emph{festin traditionnel} (ou « le \emph{ndolè}, mets traditionnel du littoral » si l'on veut garder la couleur locale).
\end{itemize}
Le paragraphe gagne en fluidité (cohésion assurée par les substituts) et en correction (registre soutenu adapté à l'examen).
\end{exoresolu}

\begin{exercice}[À toi de jouer]
1. Dans le passage suivant, relève la chaîne de référence dont le référent est \emph{la médaille} et classe chaque maillon : « La médaille brillait au soleil. Elle était accrochée à un ruban rouge. Cette récompense, Meka l'attendait depuis des années. L'objet symbolisait pour lui la reconnaissance de la France. »
2. Développe puis écris correctement, pour une première occurrence dans une dissertation, les sigles suivants : CRTV, SNH, GCE.
3. Remplace les emprunts soulignés par des termes de français standard : « Après la cérémonie, Meka a organisé un grand \emph{break} ; tout le village a dansé, le \emph{business} de la médaille l'avait rendu célèbre. »
\end{exercice}

\begin{corrige}[Exercice — Langue au service de la dissertation]
1. Chaîne de référence : \emph{La médaille} (terme introducteur) $\rightarrow$ \emph{Elle} (pronom personnel) $\rightarrow$ \emph{Cette récompense} (reprise nominale / hyperonyme) $\rightarrow$ \emph{l'} dans « l'attendait » (pronom personnel) $\rightarrow$ \emph{L'objet} (reprise nominale / hyperonyme). Cinq maillons, un seul référent.
2. CRTV : « la Cameroon Radio Television (CRTV) » ; SNH : « la Société nationale des hydrocarbures (SNH) » ; GCE : « le \emph{General Certificate of Education} (GCE) ». Le nom complet précède le sigle entre parenthèses.
3. « Après la cérémonie, Meka a organisé une grande \textbf{réception} ; tout le village a dansé : le \textbf{commerce} (ou : l'exploitation) de la médaille l'avait rendu célèbre. »
\end{corrige}

\begin{aretenir}[L'essentiel de la section]
\begin{itemize}
\item Le \cle{référent} est la réalité désignée par un mot ; la \cle{chaîne de référence} réunit tous les mots qui le désignent.
\item Les substituts (pronoms, reprises nominales, synonymes, hyperonymes) évitent les répétitions et assurent la cohésion de la dissertation — sans jamais sacrifier la clarté.
\item Un \cle{sigle} s'écrit en majuscules et se développe à la première occurrence ; un \cle{emprunt} familier ou local n'a pas sa place dans une copie d'examen, sauf citation glosée.
\end{itemize}
Tu disposes désormais de tous les outils : analyser un sujet, formuler une problématique, bâtir un plan détaillé et le rédiger dans une langue claire, variée et correcte. À toi de jouer le jour de l'examen !
\end{aretenir}

\newpage

\coursec{Activité d'intégration}

\coursec{Leçon 19 : Produire le plan détaillé d'une dissertation}

\courssub{Objectifs de l'activité}
Mobiliser l'ensemble des savoir-faire acquis au cours de la leçon (analyse de sujet, recherche d'idées, formulation de la problématique, élaboration du plan détaillé, maîtrise des outils de langue : chaînes de référence, sigles, emprunts) pour produire collectivement un plan détaillé de dissertation complet, cohérent et prêt à être rédigé, dans les conditions de l'examen du Baccalauréat technique.

\courssub{Situation}
Dans le cadre de la préparation à l'épreuve écrite de français du Baccalauréat technique, session 2027, votre classe de Terminale (séries F, G, H, BT, etc.) organise un atelier d'intégration. Le sujet proposé est un sujet « zéro » authentique, construit sur le modèle des sujets officiels MINESEC et ancré dans l'œuvre au programme : \emph{Le Vieux nègre et la médaille} de Ferdinand Oyono.

\textbf{Sujet :} « La médaille décernée à Meka dans \emph{Le Vieux nègre et la médaille} : une juste récompense ou une dérision ? Vous prendrez position en vous appuyant sur l'œuvre et sur des exemples de la vie camerounaise. »

Ce sujet vous place en situation de \cle{réflexion critique} : il ne s'agit pas de raconter l'histoire, mais d'analyser la portée symbolique et politique de la scène de la remise de médaille (chapitre 10) et de la confronter à la réalité sociologique du Cameroun d'hier (époque coloniale) et d'aujourd'hui (gestion des distinctions honorifiques, rapports entre élites et pouvoir). Vous disposez de 2 heures pour réaliser le travail en groupe (4 élèves), présenter votre plan au tableau et participer à l'évaluation croisée.

\courssub{Consignes}
Travail en groupe de 4 élèves. Chaque groupe rend une copie unique (le plan détaillé) et désigne un rapporteur pour la présentation orale (5 min max).

\begin{enumerate}
    \item \textbf{Analyse du sujet (15 min) :} Identifiez et définissez les \cle{mots-clés} (médaille, Meka, juste récompense, dérision, vie camerounaise). Dégagez le \cle{thème} et la \cle{thèse} implicite. Précisez le type de sujet (choix / position à prendre) et les contraintes (appui sur l'œuvre + exemples camerounais).
    \item \textbf{Brainstorming et tri des idées (25 min) :} Sur une feuille de brouillon, notez toutes les idées, arguments, exemples textuels (citations, scènes) et exemples extérieurs (histoire, actualité, vie locale). Organisez-les dans un \textbf{tableau à deux colonnes} : « Pour la thèse : juste récompense » / « Contre la thèse : dérision / ironie ». Éliminez le hors-sujet.
    \item \textbf{Formulation de la problématique (10 min) :} Rédigez \textbf{une seule question problématisée} qui exprime la tension du sujet et guide le développement. Elle doit intégrer la dimension littéraire et la dimension sociologique.
    \item \textbf{Choix et justification du type de plan (10 min) :} Choisissez entre plan dialectique (thèse/antithèse/synthèse), plan analytique (causes/conséquences, aspects), ou plan thématique. Justifiez votre choix en 2 lignes.
    \item \textbf{Rédaction du plan détaillé complet (45 min) :} Rédigez le plan détaillé structuré ainsi :
    \begin{itemize}
        \item \textbf{Introduction :} Accroche (citation ou fait de société), présentation de l'œuvre/auteur, énoncé du sujet, \textbf{problématique}, \textbf{annonce du plan} (formule type : « Nous examinerons d'abord... puis nous analyserons... enfin nous montrerons... »).
        \item \textbf{Développement :} 2 ou 3 parties principales (I, II, [III]), chacune subdivisée en 2 ou 3 arguments (A, B, [C]). Pour chaque argument : \textbf{idée directrice} (phrase complète), \textbf{explication/analyse}, \textbf{exemple textuel précis} (référence chapitre/scène ou citation courte), \textbf{exemple camerounais} (historique, sociologique, actualité), \textbf{transition} vers l'argument suivant.
        \item \textbf{Conclusion :} \textbf{Bilan synthétique} (réponse à la problématique), \textbf{ouverture} (sur une autre œuvre, l'actualité, une question philosophique).
    \end{itemize}
    \item \textbf{Vérification linguistique (15 min) :} Relisez le plan pour vérifier : la cohésion par les \cle{chaînes de référence} (reprises pronominales, nominatives, synonymiques) ; l'usage correct des \cle{sigles} (ex. : MINESEC, ONU, RDPC, FECAFOOT — développement à la première occurrence) et des \cle{emprunts} lexicalisés (ex. : \emph{software}, \emph{leader}, \emph{businessman} — italique si non lexicalisés, roman si francisés).
    \item \textbf{Présentation et évaluation (20 min) :} Le rapporteur présente le plan. La classe note sur une grille d'évaluation (Respect du sujet / Cohérence logique / Équilibre des parties / Pertinence des exemples / Qualité de la langue). Correction collective harmonisée par l'enseignant.
\end{enumerate}

\courssub{Ressources à mobiliser}
\begin{itemize}
    \item \textbf{Méthode d'analyse du sujet :} Mots-clés $\to$ Thème $\to$ Thèse $\to$ Type de sujet $\to$ Contraintes.
    \item \textbf{Méthode de recherche d'idées :} QQOQCP (Qui ? Quoi ? Où ? Quand ? Comment ? Pourquoi ?) + Tableau pour/contre + Sélection (pertinence, originalité).
    \item \textbf{Problématique :} Question ouverte, issue de la tension du sujet, annonciatrice du plan. Éviter la question fermée (Oui/Non) ou la paraphrase du sujet.
    \item \textbf{Types de plans :}
    \begin{itemize}
        \item \emph{Dialectique :} Thèse (arguments pour) / Antithèse (arguments contre) / Synthèse (dépassement, nuance). Idéal pour « Pour ou contre ? », « Juste ou dérision ? ».
        \item \emph{Analytique :} Causes / Conséquences ; ou Aspects (politique, social, psychologique). Idéal pour « Quelles sont les causes/conséquences ? », « Analysez... ».
        \item \emph{Thématique :} Regroupement par grands thèmes. Moins recommandé pour une prise de position.
    \end{itemize}
    \item \textbf{Structure du plan détaillé :} Introduction (4 temps) / Développement (Parties $\to$ Arguments $\to$ Exemples) / Conclusion (2 temps).
    \item \textbf{Outils de langue (Unité 19) :}
    \begin{itemize}
        \item \cle{Référent et substituts :} Assurer la cohésion (ex. : \emph{Meka, le vieux nègre, le récipiendaire, ce dernier, l'ancien chef}).
        \item \cle{Sigles :} Écrire en majuscules, développer à la 1ère occurrence (ex. : \emph{Union des Populations du Cameroun (UPC)}).
        \item \cle{Emprunts :} Respecter la graphie (italique/roman) et l'accord (ex. : \emph{des leaders} / \emph{des businessmen}).
    \end{itemize}
    \item \textbf{Connaissances sur l'œuvre :} \emph{Le Vieux nègre et la médaille}, Ferdinand Oyono, 1956. Contexte : Cameroun colonial (années 50). Personnages : Meka, Kelara, le Commandant, le Prêtre (Père Vandamme), les émissaires. Scène clé : Ch. 10 (remise médaille).
    \item \textbf{Contexte camerounais :} Ordres nationaux (Ordre de la Valeur, Ordre du Mérite), critères d'attribution, controverses médiatiques, figures historiques (Ruben Um Nyobé, Félix-Roland Moumié, Ernest Ouandié, Ahmadou Ahidjo, Paul Biya), distinctions traditionnelles (chefferies).
\end{itemize}

\begin{exoresolu}[Démarche et corrigé commenté]
\textbf{Énoncé de la tâche (Rappel)}
Sujet : « La médaille décernée à Meka dans \emph{Le Vieux nègre et la médaille} : une juste récompense ou une dérision ? Vous prendrez position en vous appuyant sur l'œuvre et sur des exemples de la vie camerounaise. »
Consignes : Analyse du sujet (mots-clés, thème, thèse, type, contraintes), Brainstorming (tableau pour/contre), Problématique, Choix du plan justifié, Plan détaillé complet (Introduction, Développement 3 parties, Conclusion), Vérification linguistique (chaînes de référence, sigles, emprunts).
\tcblower
\textbf{Étape 1 : Analyse du sujet (Modélisation)}

\begin{center}
\begin{tabularx}{\linewidth}{|l|X|}\hline
\textbf{Élément} & \textbf{Analyse} \\ \hline
Mots-clés & \textbf{Médaille} : distinction honorifique, symbole de reconnaissance officielle. \textbf{Meka} : protagoniste, « vieux nègre », ancien chef, figure de la collaboration. \textbf{Juste récompense} : thèse positive, légitimité, mérite perçu par le pouvoir colonial. \textbf{Dérision} : thèse négative, ironie, humiliation déguisée, manipulation. \textbf{Vie camerounaise} : ouverture hors texte, ancrage réel (passé/présent). \\ \hline
Thème & La signification politique et symbolique des distinctions honorifiques dans le rapport colonisateur/colonisé (et par extension dans la société camerounaise post-coloniale). \\ \hline
Thèse implicite & Le sujet suggère une opposition binaire (juste vs dérision) mais invite à dépasser le manichéisme pour montrer l'ambiguïté : ce qui paraît « juste » pour l'administration est une « dérision » pour la dignité humaine et l'histoire. \\ \hline
Type de sujet & Sujet à \textbf{thèse imposée avec nuance} / Prise de position. Contrainte double : texte (œuvre) + hors texte (vie camerounaise). \\ \hline
\end{tabularx}
\end{center}

\textbf{Étape 2 : Brainstorming et Tri (Extrait du tableau groupe-type)}

\begin{center}
\begin{tabularx}{\linewidth}{|X|X|}\hline
\rowcolor{popblueL} \textbf{Arguments « Juste récompense » (Thèse apparente / Vision coloniale)} & \textbf{Arguments « Dérision / Ironie » (Antithèse / Vision critique / Réalité)} \\ \hline
- Meka a servi loyalement l'administration 30 ans (Ch. 1). & - La médaille (médaille de « mérite civique ») arrive trop tard, à la veille de l'indépendance, pour « acheter » le silence. \\
- Reconnaissance officielle : cérémonie solennelle, discours du Commandant (Ch. 10). & - Le Commandant ne connaît même pas son nom correctement (« Meka ? Meka... »), signe du mépris. \\
- Fierté apparente de Meka et Kelara (préparatifs, tenue neuve). & - Fierté naïve, illusion : Meka « ne comprenait pas » la portée politique (Ch. 10). \\
- Exemple camerounais : Chefs de canton de l'époque du mandat français (ex. Charles Atangana) décorés de la Légion d'honneur ou de la Médaille coloniale pour « services rendus ». & - Exemple texte : Le prêtre (Père Vandamme) refuse de bénir la médaille, la qualifiant de « médaille de Judas » (Ch. 9). \\
- Exemple actuel : Ordre de la Valeur décerné à des serviteurs de l'État (enseignants, médecins). & - Exemple camerounais : Distinctions données à des « courtisans » du pouvoir, excluant les vrais résistants (UPC) ou les compétences critiques. \\
& - Ironie finale : Meka meurt peu après, la médaille sur la poitrine, symbole vain. \\ \hline
\end{tabularx}
\end{center}

\textbf{Étape 3 : Formulation de la problématique}
\emph{« En quoi la médaille remise à Meka, présentée comme une reconnaissance officielle par le pouvoir colonial, se révèle-t-elle in fine comme une dérision cruelle qui dévoile l'aliénation du collaborateur et interroge, par-delà la fiction, la portée réelle des distinctions honorifiques dans l'histoire politique du Cameroun ? »}

\textit{Justification :} Cette problématique articule les deux pôles du sujet (reconnaissance vs dérision), intègre l'œuvre (Meka, pouvoir colonial) et l'ouverture (histoire politique camerounaise), et annonce un mouvement dialectique (apparence $\to$ réalité $\to$ portée générale).

\textbf{Étape 4 : Choix du plan}
\emph{Plan dialectique (Thèse / Antithèse / Synthèse) justifié :} Le sujet pose une alternative binaire (« juste récompense OU dérision »). Le plan dialectique permet d'examiner la thèse (vision officielle/naïve), l'antithèse (réalité ironique/critique) et de synthétiser en montrant que la « dérision » l'emporte structurellement, tout en ouvrant sur la complexité actuelle des décorations.

\textbf{Étape 5 : Plan détaillé complet (Production attendue)}

\medskip
\noindent \textbf{INTRODUCTION}
\begin{itemize}
    \item \textbf{Accroche :} « Les médailles sont des piqûres de morphine pour la vanité humaine », écrivait Paul Valéry. Au Cameroun, comme ailleurs, les distinctions honorifiques rythment la vie politique.
    \item \textbf{Présentation de l'œuvre :} \emph{Le Vieux nègre et la médaille} (1956) de Ferdinand Oyono, figure majeure de la littérature camerounaise d'expression française, dépeint les derniers soubresauts de l'administration coloniale à travers le destin de Meka, ancien chef de canton.
    \item \textbf{Énoncé du sujet :} La scène finale de la remise de la médaille de « mérite civique » à Meka (Ch. 10) cristallise l'ambiguïté de cet honneur.
    \item \textbf{Problématique :} [Celle formulée à l'étape 3].
    \item \textbf{Annonce du plan :} Nous analyserons d'abord comment la cérémonie construit l'illusion d'une \textbf{juste récompense} (I), avant de démontrer que le texte et l'histoire révèlent une \textbf{dérision systémique} (II). Nous montrerons enfin que cette dérision interroge la \textbf{portée symbolique des décorations au Cameroun d'hier à aujourd'hui} (III).
\end{itemize}

\medskip
\noindent \textbf{DÉVELOPPEMENT}

\medskip
\noindent \textbf{I. L'illusion d'une juste récompense : la mise en scène du mérite officiel}
\begin{itemize}
    \item \textbf{A. Une carrière au service de l'administration coloniale.}
    \begin{itemize}
        \item \textit{Idée directrice :} Meka incarne le collaborateur modèle, fidèle pendant trois décennies.
        \item \textit{Explication :} Le roman ouvre sur son éviction (Ch. 1), mais les flashbacks montrent son zèle : perception des impôts, recrutement pour les travaux forcés, répression des récalcitrants. Il est le « bras droit » du Commandant.
        \item \textit{Exemple textuel :} Ch. 1 : « Il avait servi le blanc pendant trente ans... » ; Ch. 10 : Le discours du Commandant énumère ses « services éminents ».
        \item \textit{Exemple camerounais :} Les chefs de canton de l'époque du mandat français (ex. Charles Atangana à Yaoundé) décorés de la Légion d'honneur ou de la Médaille coloniale pour « pacification » et administration.
        \item \textit{Transition :} Cette fidélité administrative fonde la légitimité apparente de la médaille.
    \end{itemize}
    \item \textbf{B. Une cérémonie solennelle, rituel de la reconnaissance.}
    \begin{itemize}
        \item \textit{Idée directrice :} Le protocole colonial investit l'événement d'une solennité qui valide le mérite.
        \item \textit{Explication :} Déploiement de la garde indigène, présence du Prêtre, discours officiel, remise de l'insigne, banquet. Meka revêt son « uniforme » (veste kaki, casque), symbole de son statut.
        \item \textit{Exemple textuel :} Ch. 10 : Description minutieuse du décor, du « soleil qui darde », du Commandant « tout blanc dans son uniforme ». La phrase : « Meka s'avança, raide, l'œil fixe... »
        \item \textit{Exemple camerounais :} Les cérémonies actuelles du 20 mai (Fête nationale) ou du 10 mai (Fête du Travail) au Palais de l'Unité ou dans les chefs-lieux de région, où le MINESEC, le MINSANTE, etc., décorent des agents.
        \item \textit{Transition :} Mais cette théâtralisation masque une réalité tout autre, perçue par le regard lucide de l'auteur.
    \end{itemize}
\end{itemize}

\medskip
\noindent \textbf{II. La réalité de la dérision : ironie, manipulation et aliénation}
\begin{itemize}
    \item \textbf{A. L'ironie narrative : le décalage entre le signifiant et le signifié.}
    \begin{itemize}
        \item \textit{Idée directrice :} L'écriture d'Oyono instaure une distance ironique qui déconstruit l'hommage.
        \item \textit{Explication :} Le narrateur adopte le point de vue de Meka (focalisation interne) pour mieux en montrer la naïveté. Le Commandant lit un texte préfabriqué, ne regardant pas Meka. La médaille est une « petite chose métallique ».
        \item \textit{Exemple textuel :} Ch. 10 : « Le Commandant lut d'une voix monotone... » ; « Meka ne comprenait pas... » ; La comparaison de la médaille à une « breloque » (Ch. 11). Le refus du Père Vandamme (Ch. 9) : « C'est la médaille de Judas. »
        \item \textit{Exemple camerounais :} La remise de médailles à des personnalités controversées (affaires de corruption, détournements) qui discrédite l'institution (ex. Opération \emph{Épervier} où des décorés sont ensuite poursuivis).
        \item \textit{Transition :} Cette ironie révèle que la médaille est un outil de manipulation politique.
    \end{itemize}
    \item \textbf{B. Un instrument de domination : « diviser pour régner » à la veille de l'indépendance.}
    \begin{itemize}
        \item \textit{Idée directrice :} La médaille vise à neutraliser les élites traditionnelles face au nationalisme montant (UPC).
        \item \textit{Explication :} Contexte 1956 : insurrection nationaliste. L'administration a besoin de la loyauté des chefs. La médaille « achète » le silence de Meka et son ralliement public.
        \item \textit{Exemple textuel :} Ch. 10 : Le discours du Commandant insiste sur « l'ordre », « la paix », « la France généreuse » $\to$ propagande anti-indépendantiste. Kelara (l'épouse) est la seule à voir la vérité : « Ils se moquent de toi. »
        \item \textit{Exemple camerounais :} L'histoire de l'UPC (Union des Populations du Cameroun) : les chefs décorés par l'administration coloniale (ou post-coloniale) étaient souvent ceux qui combattaient les « maquisards » (Ruben Um Nyobé, Ernest Ouandié), créant une fracture mémorielle.
        \textit{Transition :} Au-delà du contexte colonial, cette dérision questionne la valeur même de la distinction honorifique dans la société camerounaise.
    \end{itemize}
\end{itemize}

\medskip
\noindent \textbf{III. Portée symbolique : les distinctions honorifiques au Cameroun, entre mérite et clientélisme}
\begin{itemize}
    \item \textbf{A. La banalisation des ordres nationaux : un enjeu de reconnaissance citoyenne.}
    \begin{itemize}
        \item \textit{Idée directrice :} Les Ordres de la Valeur et du Mérite (créés 1960, réformés 1972, 1992) peinent à garantir l'exemplarité.
        \item \textit{Explication :} Critères officiels : « services éminents rendus à la Nation ». Réalité : quotas par ministère, lobbying, clientélisme politique. L'\emph{Ordre de la Valeur} perd de sa valeur symbolique.
        \item \textit{Exemple camerounais :} Polémiques récurrentes dans la presse (ex. \emph{Mutations}, \emph{Le Messager}) sur les listes du 20 mai : absence de figures de la société civile (enseignants, médecins de brousse, artistes) au profit de hauts fonctionnaires en fin de carrière ou de militants du parti au pouvoir (\emph{Rassemblement Démocratique du Peuple Camerounais (RDPC)}).
        \item \textit{Transition :} Pourtant, la distinction conserve une fonction sociale forte.
    \end{itemize}
    \item \textbf{B. La distinction comme levier de cohésion ou de fracture mémorielle.}
    \begin{itemize}
        \item \textit{Idée directrice :} Réhabiliter les « oubliés » de l'histoire (résistants) vs entretenir le clientélisme.
        \item \textit{Explication :} La médaille peut réparer l'injustice (réhabilitation posthume de Ruben Um Nyobé, Félix-Roland Moumié, décoration de vétérans des deux guerres mondiales). Mais elle peut aussi figer les hiérarchies.
        \item \textit{Exemple camerounais :} Cérémonie de 2010 (50e anniversaire) : décorations de vétérans et de figures de la résistance. Débat public sur les critères : mérite objectif (résultats sportifs \emph{Fédération Camerounaise de Football (FECAFOOT)}, scientifiques) vs allégeance politique.
        \item \textit{Exemple textuel (renvoi) :} La mort de Meka avec sa médaille (Ch. 11) symbolise l'inanité d'un honneur sans conscience historique.
    \end{itemize}
\end{itemize}

\medskip
\noindent \textbf{CONCLUSION}
\begin{itemize}
    \item \textbf{Bilan synthétique :} La médaille de Meka n'est pas une « juste récompense » au sens moral, mais une \textbf{dérision tragique}. Elle sanctionne non le mérite, mais l'aliénation consentie. Le plan dialectique a permis de passer de l'apparence (cérémonie) à la réalité (manipulation), pour aboutir à une critique structurelle de l'institution honorifique.
    \item \textbf{Ouverture :} Au-delà d'Oyono, cette analyse rejoint \emph{Ville cruelle} de Mongo Beti ou \emph{L'Aventure ambiguë} de Cheikh Hamidou Kane : la question de la collaboration et de la reconnaissance par l'Autre. Aujourd'hui, le Cameroun gagnerait à instituer un \emph{Collège des distinctions} indépendant (sur le modèle du \emph{Conseil de l'Ordre de la Légion d'honneur} en France) pour restaurer la crédibilité de la médaille. « Une médaille ne fait pas le mérite, c'est le mérite qui fait la médaille. »
\end{itemize}

\textbf{Étape 6 : Vérification linguistique (Exemples corrigés dans le plan)}
\begin{itemize}
    \item \textbf{Chaînes de référence (Meka) :} \emph{Meka (Ch. 1) $\to$ le vieux nègre (titre/Ch. 2) $\to$ l'ancien chef (Ch. 5) $\to$ le récipiendaire (Ch. 10) $\to$ ce dernier (Ch. 10) $\to$ le décoré (Ch. 11).} Variété lexicale assurée, pas de répétition lourde.
    \item \textbf{Sigles :} \emph{Union des Populations du Cameroun (UPC)} [1ère occ.] $\to$ \emph{l'UPC} [suivantes]. \emph{Rassemblement Démocratique du Peuple Camerounais (RDPC)}. \emph{Fédération Camerounaise de Football (FECAFOOT)}. \emph{Ministère des Enseignements Secondaires (MINESEC)}. Majuscules respectées, développement fait.
    \item \textbf{Emprunts :} \emph{leader} (francisé, roman, pluriel \emph{leaders}) ; \emph{software} (non francisé, italique) — non utilisé ici ; \emph{businessman} (francisé \emph{businessman} ou \emph{hommes d'affaires} préféré) ; \emph{quota} (francisé). Dans le plan : \emph{clientélisme} (francisé), \emph{lobbying} (italique \emph{lobbying} car anglicisme récent).
\end{itemize}
\end{exoresolu}

\courssub{Bilan de l'activité}
Cet atelier d'intégration a permis de vérifier l'acquisition opérationnelle des compétences ciblées par la leçon 19 :
\begin{enumerate}
    \item \textbf{Maîtrise de la méthode :} Les élèves ont appliqué rigoureusement les 4 temps de l'analyse de sujet, le brainstorming structuré (tableau pour/contre), la formulation d'une problématique « moteur » et le choix justifié du plan dialectique, seul adapté à l'alternative du sujet.
    \item \textbf{Articulation texte / hors-texte :} La contrainte « œuvre + vie camerounaise » a été respectée : chaque argument littéraire (scène, citation, focalisation) est relayé par un exemple historique ou sociologique précis (chefferies, UPC, Ordres nationaux, actualité 20 mai), évitant le hors-sujet ou la dissertation hors texte.
    \item \textbf{Qualité du plan détaillé :} Les productions montrent une progression nette : les idées directrices sont des phrases complètes (arguments), les transitions explicites, l'annonce de plan standardisée. L'équilibre des parties (I : 2 args, II : 2 args, III : 2 args) est respecté.
    \item \textbf{Outils de langue au service de la cohérence :} La vérification ciblée a révélé des progrès sur les chaînes de référence (diversification des substituts de \emph{Meka}) et le traitement des sigles (développement systématique à la première occurrence). L'usage des emprunts (\emph{lobbying}, \emph{leader}) est maîtrisé (italique/accord).
    \item \textbf{Difficultés résiduelles (remédiation) :} Certains groupes ont encore tendance à rédiger des « résumés d'arguments » au lieu d'\emph{idées directrices argumentatives} (ex. : « Meka est content » au lieu de « La cérémonie valide l'illusion du mérite »). La distinction \emph{exemple textuel} / \emph{exemple camerounais} reste floue pour quelques élèves qui mélangent les deux. La synthèse (partie III) peine parfois à dépasser le simple « pour/contre » pour offrir une véritable hauteur de vue. Ces points feront l'objet d'exercices ciblés en séance de remédiation.
\end{enumerate}
L'activité, par son format « examen blanc » et son évaluation par les pairs (grille), a ancré les automatismes nécessaires pour l'épreuve du Probatoire / Baccalauréat technique.

\newpage

\coursec{Méthodes et exercices résolus}

\coursec{Produire le plan détaillé d'une dissertation}

\courssub{Réception du texte : Paratextes et lecture ouvroir de \emph{Le Vieux nègre et la médaille}}

\begin{methode}[Analyser les paratextes d'une œuvre]
L'analyse des paratextes (tout ce qui entoure le texte principal) permet d'entrer dans l'œuvre, de formuler des hypothèses de lecture et de situer le contrat de lecture.
\begin{enumerate}
\item \textbf{Le titre} : \emph{Le Vieux nègre et la médaille}. Identifier la structure (nom + complément), le sens des mots (\cle{vieux} : âge, expérience, usure ; \cle{nègre} : terme colonial péjoratif assumé par l'auteur pour dénoncer ; \cle{médaille} : symbole de reconnaissance officielle, ironie annoncée). Confronter titre et sous-titre s'il y a lieu.
\item \textbf{L'auteur} : Ferdinand Oyono (1929--2010), écrivain et diplomate camerounais. Figure majeure de la littérature africaine d'expression française. Œuvres : \emph{Une vie de boy} (1956), \emph{Le Vieux nègre et la médaille} (1956), \emph{Chemin d'Europe} (1960). Thèmes : critique du colonialisme, aliénation, identité.
\item \textbf{La couverture / illustration} : Observer l'image (souvent un vieillard, une médaille, un décor rural). Analyser les couleurs, la composition, le regard du personnage.
\item \textbf{La quatrième de couverture / résumé éditeur} : Repérer les informations données sur l'intrigue (Meka, la médaille, la déception) et la tonalité (ironie, drame).
\item \textbf{Les mentions éditoriales} : Collection (souvent "L'école des lettres" ou "Classiques africains"), date de publication originale (1956), éditeur. Situer dans l'histoire littéraire (pré-indépendance, littérature de combat).
\item \textbf{Le sommaire / structure} : Roman court, pas de chapitres numérotés mais des séquences narratives. Repérer les grandes étapes : l'attente, la cérémonie, le retour, la nuit de révélation.
\end{enumerate}
\end{methode}

\begin{exemplebox}[Exploitation des paratextes en classe]
\textbf{Activité guidée :} À partir de la page de couverture et de la quatrième de couverture de l'édition de référence (ex. CLE International ou Hatier) :
\begin{enumerate}
\item Relevez les informations explicites (auteur, titre, éditeur, collection, ISBN).
\item Formulez 3 hypothèses sur l'intrigue à partir du titre seul.
\item Relevez les mots-clés de la quatrième de couverture qui annoncent le thème de la \cle{colonisation} et de l'\cle{aliénation}.
\item Quelle est la fonction du mot "nègre" dans le titre ? (Provocation, réalisme historique, reprise du vocabulaire colonial pour le retourner).
\end{enumerate}
\end{exemplebox}

\begin{methode}[La lecture ouvroir (lecture cursive / découverte)]
La lecture ouvroir est une lecture silencieuse, individuelle et continue de l'œuvre intégrale, réalisée \textbf{avant} l'étude analytique en classe. Objectifs : saisir l'intrigue globale, identifier les personnages, ressentir l'émotion esthétique, repérer les passages marquants.
\begin{enumerate}
\item \textbf{Préparation} : Distribuer l'œuvre (ou extrait long) une semaine avant. Donner une fiche de lecture guidée (personnages, lieux, temps, schéma narratif).
\item \textbf{Consignes} : Lire sans s'arrêter au dictionnaire. Noter en marge : \cle{?} (incompréhension), \cle{!} (surprise/émotion), \cle{Thème} (colonisation, famille, mort, honneur).
\item \textbf{Exploitation en classe (1 séance)} :
\begin{itemize}
\item Tour de table : impressions premières (ennui, colère, pitié, admiration pour l'écriture ?).
\item Mise en place du \cle{schéma narratif} : Situation initiale (Meka attend le commandant), Élément déclencheur (invitation à la cérémonie), Péripéties (préparatifs, discours, remise de médaille), Élément de résolution (retour au village, nuit de réflexion), Situation finale (Meka enterre la médaille / lucidité).
\item Identification des \cle{personnages principaux} : Meka (protagoniste, vieux cultivateur), Kelara (épouse, lucide), Le Commandant (représentant de l'autorité coloniale), Le Prêtre (complicité religieuse), Les fils (absents, partis à la ville/brousse), Les villageois (chorale, témoins).
\item Repérage des \cle{lieux et temps} : Village de Doum (fictif, Cameroun), années 1930-40 (période coloniale française), durée : quelques jours.
\end{itemize}
\end{enumerate}
\end{methode}

\begin{aretenir}[Fiche de lecture synthétique -- Le Vieux nègre et la médaille]
\begin{itemize}
\item \textbf{Auteur} : Ferdinand Oyono (Cameroun).
\item \textbf{Genre} : Roman court / Nouvelle longue / Récit initiatique inversé (désillusion).
\item \textbf{Thèmes majeurs} : Critique du système colonial (assimilation, exploitation, hypocrisie), Aliénation / Prise de conscience, Valeurs traditionnelles vs valeurs importées, Vanité des honneurs coloniaux, Relation père/fils, Place de la femme (Kelara, voix de la raison).
\item \textbf{Style} : Ironie, sous-entendu, focalisation interne (vision de Meka), langue française "africanisée" (proverbes, rythme oral), réalisme descriptif.
\end{itemize}
\end{aretenir}

\courssub{Lecture méthodique n°1 et contrôle de lecture}

\begin{methode}[Conduire une lecture méthodique (explication de texte)]
La lecture méthodique porte sur un extrait signifiant (ex. : la scène de la remise de médaille, pp. 60-70 env., ou la nuit finale de Meka). Elle suit trois mouvements :
\begin{enumerate}
\item \textbf{Approche globale} : Situer l'extrait dans l'œuvre (début, milieu, fin), identifier le thème dominant, le ton (ironique, pathétique, dramatique), le type de discours (narratif, descriptif, dialogué).
\item \textbf{Analyse détaillée} (par paragraphes ou mouvements) :
\begin{itemize}
\item \emph{Le contenu} : Qui ? Quoi ? Où ? Quand ? Comment ? Pourquoi ?
\item \emph{La forme} : Figures de style (ironie, antiphrase, métaphore, comparaison), champs lexicaux, syntaxe (phrases courtes/nouvelles, interrogatives), temps verbaux, focalisation.
\item \emph{L'interprétation} : Que révèle ce passage sur les personnages, la thèse de l'auteur, la condition coloniale ?
\end{itemize}
\item \textbf{Synthèse} : Résumer l'apport de l'extrait à la compréhension de l'œuvre entière. Lien avec la problématique de la séquence (ex. : "L'aliénation coloniale est-elle irréversible ?").
\end{enumerate}
\end{methode}

\begin{exoresolu}[Lecture méthodique : Extrait de la cérémonie (Remise de la médaille)]
Énoncé : Étudier le passage où le Commandant épingle la médaille sur la poitrine de Meka (éd. CLE, p. 65-66 environ). Analyser l'ironie de la situation.
\tcblower
\textbf{1. Situation :} Cérémonie officielle, place publique, foule villageoise. Meka est au centre, "comme un roi".
\textbf{2. Analyse de l'ironie :}
\begin{itemize}
\item \emph{Antiphrase} : "Meka était heureux. Il était fier." (Le lecteur sait que c'est une mascarade).
\item \emph{Contraste} : La solennité du rituel (drapeau, hymne, sabre, discours ampoulé) vs la réalité de la médaille (métal clinquant, "breloque", "chose sans valeur").
\item \emph{Le discours du Commandant} : Langue de bois, formules toutes faites ("dévouement", "loyauté", "France généreuse"). Meka ne comprend pas le français, il subit la forme sans saisir le fond vide.
\item \emph{Le regard de Kelara} : Seule, elle ne sourit pas. Elle voit le ridicule. Elle incarne la lucidité que Meka n'a pas encore.
\end{itemize}
\textbf{3. Synthèse :} Cet extrait est le point culminant de l'aliénation de Meka. L'ironie dramatique (le lecteur sait, le personnage ignore) prépare la chute finale. La médaille devient le symbole de la duperie coloniale : on décore l'exploité pour mieux l'asservir.
\end{exoresolu}

\begin{exercice}[Contrôle de lecture -- Questionnaire type Bac]
\begin{enumerate}
\item \textbf{Question de compréhension (4 pts) :} Pourquoi Meka a-t-il reçu la médaille ? Citez deux "mérites" invoqués par l'administration coloniale dans le roman.
\item \textbf{Question d'analyse (6 pts) :} Étudiez le personnage de Kelara. En quoi représente-t-elle la conscience lucide du village ? Appuyez votre réponse sur deux passages précis du texte.
\item \textbf{Question d'interprétation (5 pts) :} "La médaille est un leurre." Commentez cette affirmation en vous appuyant sur la fin du roman (l'enterrement de la médaille).
\item \textbf{Question de langue (5 pts) :} Relevez dans le texte trois emprunts à l'oralité (proverbes, syntaxe, vocabulaire) et expliquez leur effet de sens.
\end{enumerate}
\end{exercice}

\begin{corrige}[Exercice n°1 -- Contrôle de lecture]
\begin{enumerate}
\item \textbf{Compréhension :} Meka reçoit la médaille pour sa "loyauté" à l'administration française : il a fourni des porteurs, payé l'impôt sans réclamer, envoyé ses fils à l'école/au travail forcé, dénoncé les "agitateurs". (2 pts par mérite cité).
\item \textbf{Analyse Kelara :} Elle refuse la fête (p. 20 : "Cette médaille ne nous rendra pas nos fils"), elle voit la misère derrière l'apparat (p. 60 : elle ne danse pas), elle pousse Meka à la réflexion finale. Elle est la gardienne des valeurs réelles (terre, lignée, dignité).
\item \textbf{Interprétation :} La médaille est un leurre car elle masque la spoliation (terres prises, impôts, travail forcé). L'enterrement final (Meka la jette dans la latrine / la cache) symbolise le rejet de l'honneur colonial et la reprise de sa dignité d'homme libre. C'est un acte symbolique de décolonisation mentale.
\item \textbf{Langue :} Ex. : "Le blanc est comme le fromage : il a des trous" (métaphore/proverbe, critique de l'incohérence coloniale) ; "La parole est un œuf" (proverbe, valeur de la parole) ; Syntaxes elliptiques, présent de narration oralisé. Effet : ancrage culturel, distance ironique vis-à-vis du français officiel.
\end{enumerate}
\end{corrige}

\courssub{Analyser le sujet de dissertation}

\begin{methode}[Analyser un sujet de dissertation]
Pour analyser correctement un sujet de dissertation, on procède en étapes :
\begin{enumerate}
\item \textbf{Lire le sujet plusieurs fois} et souligner les mots-clés (les termes porteurs de sens).
\item \textbf{Définir chaque mot-clé} : sens propre, sens figuré, champ lexical associé.
\item \textbf{Repérer le type de sujet} : sujet de réflexion (question ouverte), sujet d'opinion (faut-il... ?), sujet de type affirmation à discuter.
\item \textbf{Dégager les présupposés} : ce que le sujet suppose vrai et qu'il faut vérifier ou nuancer.
\item \textbf{Formuler la tension} : opposer deux lectures possibles du sujet (pour / contre ; oui / mais).
\item \textbf{Transformer cette tension en problématique} : une question (ou deux) qui commande tout le devoir.
\end{enumerate}
Réflexe : ne jamais commencer à rédiger avant d'avoir écrit la problématique au brouillon.
\end{methode}

\begin{exoresolu}[Analyse d'un sujet type Bac]
Énoncé : « L'école est-elle la seule voie de la réussite sociale ? » Analysez ce sujet en suivant la méthode (mots-clés, type de sujet, présupposé, tension, problématique).
\tcblower
\textbf{Étape 1 — Mots-clés :} \cle{école}, \cle{seule voie}, \cle{réussite sociale}.

\textbf{Étape 2 — Définitions :}
\begin{itemize}
\item \emph{École} : institution d'instruction, mais aussi toute formation structurée (apprentissage technique, professionnel).
\item \emph{Seule voie} : exclusivité ; le sujet affirme qu'il n'existerait pas d'autre chemin.
\item \emph{Réussite sociale} : ascension sociale, insertion professionnelle, reconnaissance, bien-être matériel et moral.
\end{itemize}

\textbf{Étape 3 — Type de sujet :} sujet d'opinion (question fermée « est-elle... ? » appelant oui / non / nuance).

\textbf{Étape 4 — Présupposé :} le sujet suppose que l'école mène à la réussite ; or le chômage des diplômés et la réussite d'artisans ou d'entrepreneurs non diplômés invitent à nuancer.

\textbf{Étape 5 — Tension :} d'un côté, l'école donne des savoirs, des diplômes, ouvre les portes de l'emploi salarié ; de l'autre, l'expérience, le talent, l'apprentissage sur le tas et l'entrepreneuriat permettent aussi de réussir.

\textbf{Étape 6 — Problématique :} « L'école est-elle une condition indispensable de la réussite sociale, ou n'est-elle qu'une voie parmi d'autres ? »

\textbf{Conclusion :} le sujet appelle un plan dialectique (thèse / antithèse / synthèse) ou un plan thématique nuancé.
\end{exoresolu}

\courssub{Rechercher les idées et formuler la problématique}

\begin{methode}[Rechercher les idées]
\begin{enumerate}
\item \textbf{Faire un remue-méninges} (brainstorming) : noter au brouillon tout ce qui vient à l'esprit autour des mots-clés, sans trier.
\item \textbf{Mobiliser trois sources d'idées} : les textes étudiés en classe (ex. \emph{Le Vieux nègre et la médaille} de Ferdinand Oyono), la culture générale (histoire, actualité camerounaise et africaine), l'expérience personnelle et l'observation de la vie courante.
\item \textbf{Illustrer chaque idée} : une idée sans exemple (fait, texte, personnage, statistique simple) est une idée morte.
\item \textbf{Trier et regrouper} : éliminer les hors-sujets, regrouper les idées voisines en deux ou trois grands ensembles.
\item \textbf{Vérifier la couverture du sujet} : chaque mot-clé doit être traité.
\item \textbf{Formuler la problématique} : une question claire, sans « est-ce que » lourd si possible, qui annonce les grandes parties.
\end{enumerate}
Formule à retenir : \textbf{1 idée = 1 argument + 1 exemple + 1 analyse.}
\end{methode}

\begin{exoresolu}[Recherche d'idées et problématique]
Énoncé : Sujet : « La lecture est-elle encore utile à l'ère du téléphone portable ? » Faites la recherche d'idées (au moins six idées avec exemples), triez-les et formulez la problématique.
\tcblower
\textbf{Étape 1 — Remue-méninges :} lecture, évasion, vocabulaire, réussite scolaire, distraction, réseaux sociaux, information rapide, superficialité, coût des livres, lecture sur écran.

\textbf{Étape 2 — Idées retenues et illustrées :}
\begin{itemize}
\item La lecture développe le vocabulaire et l'expression (ex. : un élève lecteur réussit mieux la dissertation au Bac).
\item La lecture forme l'esprit critique et l'imagination (ex. : les romans d'Oyono font comprendre la colonisation mieux qu'un discours).
\item La lecture donne accès aux savoirs durables (ex. : manuels, romans classiques, presse d'analyse).
\item Le téléphone portable offre une information rapide mais fragmentée (ex. : messages courts, vidéos, réseaux sociaux).
\item L'usage excessif du portable nuit à la concentration (ex. : élèves qui zappent et ne finissent plus un texte long).
\item Le portable peut aussi servir à lire (ex. : livres numériques, applications éducatives).
\end{itemize}

\textbf{Étape 3 — Tri et regroupement :} ensemble A : les bienfaits irremplaçables de la lecture (idées 1, 2, 3) ; ensemble B : les limites du téléphone portable (idées 4, 5) ; ensemble C : une complémentarité possible (idée 6).

\textbf{Étape 4 — Problématique :} « Malgré la concurrence du téléphone portable, la lecture conserve-t-elle une utilité propre, et comment concilier les deux ? »

\textbf{Conclusion :} ces trois ensembles annoncent un plan en trois parties : l'utilité de la lecture, les limites du portable, leur complémentarité.
\end{exoresolu}

\courssub{Élaborer le plan détaillé}

\begin{methode}[Construire un plan détaillé]
\begin{enumerate}
\item \textbf{Choisir le type de plan} adapté au sujet :
\begin{itemize}
\item \emph{Plan dialectique} (thèse / antithèse / synthèse) pour les sujets d'opinion ;
\item \emph{Plan thématique} (aspects successifs d'une notion) pour les sujets de réflexion ;
\item \emph{Plan analytique} (constat / causes / conséquences ou solutions) pour les sujets de faits sociaux.
\end{itemize}
\item \textbf{Rédiger les intitulés de parties} en phrases complètes ou en formules nominales claires, jamais en un seul mot vague.
\item \textbf{Diviser chaque partie en deux ou trois sous-parties} (A, B, C), chacune portant un argument.
\item \textbf{Associer à chaque sous-partie un exemple précis} (texte étudié, fait de la vie courante, actualité).
\item \textbf{Prévoir les transitions} : une phrase qui relie chaque partie à la suivante.
\item \textbf{Vérifier l'équilibre} : des parties de longueur comparable, aucune idée répétée, aucun hors-sujet.
\end{enumerate}
\end{methode}

\begin{exoresolu}[Production d'un plan détaillé]
Énoncé : Sujet : « Faut-il condamner la colonisation ? » (sujet en lien avec \emph{Le Vieux nègre et la médaille}). Élaborez un plan détaillé complet (parties, sous-parties, exemples, transitions).
\tcblower
\textbf{Choix du plan :} sujet d'opinion $\Rightarrow$ plan dialectique.

\textbf{Problématique :} « La colonisation, entre domination et apports matériels, mérite-t-elle d'être condamnée ? »

\textbf{Introduction (au brouillon) :} amorce (l'Afrique a connu la colonisation) — présentation du sujet — problématique — annonce du plan.

\textbf{I. La colonisation est condamnable car elle fut une domination (thèse)}
\begin{itemize}
\item A. Elle humilie et déshumanise les colonisés. \emph{Exemple :} dans \emph{Le Vieux nègre et la médaille}, Meka reçoit une médaille dérisoire après avoir perdu ses terres et ses fils.
\item B. Elle spolie les terres et exploite les hommes (travail forcé, impôts). \emph{Exemple :} les corvées et la réquisition des terres des ancêtres de Meka.
\end{itemize}
\emph{Transition :} « Certains objectent cependant que la colonisation a laissé des réalisations matérielles. »

\textbf{II. La colonisation a cependant eu des apports (antithèse)}
\begin{itemize}
\item A. Infrastructures : routes, écoles, hôpitaux. \emph{Exemple :} les premières écoles au Cameroun.
\item B. Ouverture au monde : langues, techniques, administration moderne.
\end{itemize}
\emph{Transition :} « Mais ces apports, limités et intéressés, effacent-ils le bilan de la domination ? »

\textbf{III. Synthèse : une condamnation lucide, tournée vers l'avenir}
\begin{itemize}
\item A. Les apports servaient d'abord l'exploitation : la balance penche vers la condamnation.
\item B. L'essentiel est de construire l'indépendance réelle : éducation, mémoire, développement. \emph{Exemple :} la lucidité finale de Meka, qui découvre la vanité de la médaille.
\end{itemize}

\textbf{Conclusion :} réponse à la problématique (oui, il faut condamner la colonisation comme système de domination) et ouverture (la mémoire au service de l'unité africaine).
\end{exoresolu}

\courssub{Le référent et ses substituts : outils de la cohérence}

\begin{methode}[Éviter les répétitions grâce aux substituts]
Dans une dissertation, un même \cle{référent} (la personne ou la chose dont on parle) ne doit pas être répété tel quel. On utilise des \cle{substituts} :
\begin{enumerate}
\item \textbf{Le pronom personnel} : il, elle, ils (ex. : Meka $\rightarrow$ il).
\item \textbf{Le pronom démonstratif} : celui-ci, celui-là, ce dernier.
\item \textbf{Le nom générique} : le personnage, le vieillard, le héros, cet homme.
\item \textbf{La périphrase} : le protagoniste du roman, l'époux de Kelara.
\item \textbf{Le synonyme ou quasi-synonyme} : le colonisateur $\rightarrow$ l'administrateur, le Blanc.
\end{enumerate}
Réflexe : à la relecture, souligner les noms propres répétés et remplacer une occurrence sur deux par un substitut adapté. Attention : le substitut doit être \textbf{clair} (on doit savoir de qui on parle) et \textbf{varié}.
\end{methode}

\begin{exoresolu}[Reprise d'un paragraphe répétitif]
Énoncé : Réécrivez ce paragraphe en remplaçant les répétitions par des substituts variés : « Meka reçoit une médaille. Meka est fier de cette médaille. Meka comprend vite que la médaille ne rend ni ses terres ni ses fils. Meka rentre dans son village, déçu. »
\tcblower
\textbf{Étape 1 — Identifier le référent :} Meka (4 occurrences) ; secondairement « la médaille » (3 occurrences).

\textbf{Étape 2 — Choisir des substituts :} pour Meka : \emph{il}, \emph{le vieillard}, \emph{le héros d'Oyono} ; pour la médaille : \emph{cette récompense}, \emph{elle}, \emph{ce ruban de métal}.

\textbf{Étape 3 — Réécriture :} « Meka reçoit une médaille. Le vieillard est d'abord fier de cette récompense. Mais il comprend vite qu'elle ne lui rend ni ses terres ni ses fils. C'est en homme déçu que le héros d'Oyono rentre dans son village. »

\textbf{Vérification :} chaque substitut renvoie clairement à son référent ; aucune répétition inutile ne subsiste ; le paragraphe gagne en fluidité.

\textbf{Conclusion :} la chaîne de référence (Meka $\rightarrow$ le vieillard $\rightarrow$ il $\rightarrow$ le héros d'Oyono) assure la cohérence du texte sans lourdeur.
\end{exoresolu}

\courssub{Sigles et emprunts : les employer correctement}

\begin{methode}[Utiliser sigles et emprunts dans une copie]
\begin{enumerate}
\item \textbf{Sigle} : mot formé des initiales d'un groupe de mots, prononcé lettre par lettre (ONU, MINESEC) ; \textbf{acronyme} : sigle prononcé comme un mot (SIDA, CNUED... à prononcer selon l'usage). Règle : à la \textbf{première occurrence}, écrire le sigle en capitales et, si nécessaire, donner la signification entre parenthèses : « l'ONU (Organisation des Nations unies) ».
\item \textbf{Emprunt} : mot venant d'une langue étrangère (ex. : \emph{week-end}, \emph{business}, \emph{développement durable} pour \emph{sustainable development}). Règle : préférer l'équivalent français quand il existe (\emph{fin de semaine}, \emph{affaires}) ; si l'emprunt est conservé, l'écrire en italique ou entre guillemets.
\item \textbf{Accord} : les sigles sont en général invariables (des ONG) ; l'accord se fait selon le nom principal pour le genre (la SNC, car \emph{société}).
\item \textbf{Pertinence} : n'employer un sigle ou un emprunt que s'il est connu du lecteur ou expliqué.
\end{enumerate}
\end{methode}

\begin{exoresolu}[Correction d'un passage]
Énoncé : Corrigez et justifiez : « Les ong luttent contre le sida. Le business du bois attire des investisseurs qui organisent des brainstorming pour moderniser la filière. »
\tcblower
\textbf{Correction 1 — « ong » :} un sigle s'écrit en capitales : \textbf{ONG}. À la première occurrence, on peut préciser : « les ONG (organisations non gouvernementales) ». Invariable au pluriel : « des ONG ».

\textbf{Correction 2 — « sida » :} acronyme (se prononce comme un mot) : l'usage courant admet la minuscule, mais la forme rigoureuse est \textbf{SIDA} (syndrome d'immunodéficience acquise). Dans une copie, écrire « le SIDA » ou « le sida » de façon constante.

\textbf{Correction 3 — « business » :} emprunt à l'anglais ; l'équivalent français existe : « le \textbf{commerce} du bois » ou « l'\textbf{exploitation} du bois ». Si l'on garde l'emprunt : \emph{business} en italique.

\textbf{Correction 4 — « brainstorming » :} emprunt récent ; préférer « \textbf{remue-méninges} » ou « des séances de réflexion collective ».

\textbf{Texte corrigé :} « Les ONG (organisations non gouvernementales) luttent contre le SIDA. Le commerce du bois attire des investisseurs qui organisent des remue-méninges pour moderniser la filière. »

\textbf{Conclusion :} en français soigné, on écrit les sigles en capitales, on les explicite à la première occurrence et on préfère les équivalents français aux emprunts.
\end{exoresolu}

\courssub{Appliquer la démarche complète à une situation concrète}

\begin{methode}[Rédiger et justifier une démarche complète]
Pour traiter un sujet de dissertation de A à Z au brouillon :
\begin{enumerate}
\item Analyser le sujet (mots-clés, type, présupposé).
\item Rechercher et trier les idées (3 sources : textes, culture, expérience).
\item Formuler la problématique.
\item Choisir le type de plan et rédiger le plan détaillé (parties, sous-parties, exemples, transitions).
\item Justifier chaque choix : pourquoi ce plan ? pourquoi cet exemple ? quelle idée sert quelle partie ?
\end{enumerate}
Au Bac, le plan détaillé au brouillon est votre \textbf{contrat de rédaction} : la copie doit le suivre fidèlement.
\end{methode}

\begin{exoresolu}[Sujet complet type Baccalauréat technique]
Énoncé : Sujet : « Le travail manuel est-il dévalorisé dans notre société ? » Produisez l'analyse du sujet, la problématique et le plan détaillé, puis justifiez vos choix.
\tcblower
\textbf{1. Analyse du sujet.} Mots-clés : \emph{travail manuel} (métiers de l'artisanat, de la technique, de la terre), \emph{dévalorisé} (méprisé, mal payé, peu considéré), \emph{notre société} (le Cameroun contemporain). Type : sujet d'opinion. Présupposé : le travail intellectuel serait préféré au travail manuel — à vérifier.

\textbf{2. Problématique :} « Le travail manuel souffre-t-il d'un mépris réel, et comment lui rendre sa valeur ? »

\textbf{3. Plan détaillé (dialectique).}

\textbf{I. Le travail manuel est effectivement dévalorisé (thèse)}
\begin{itemize}
\item A. Mépris social : on pousse les élèves vers les filières « nobles » ; l'artisan est regardé de haut. \emph{Exemple :} familles qui rougissent d'un fils menuisier.
\item B. Faible rémunération et absence de protection. \emph{Exemple :} journaliers du bâtiment sans contrat.
\end{itemize}
\emph{Transition :} « Pourtant, cette dévalorisation est injuste et de moins en moins vraie. »

\textbf{II. Le travail manuel est en réalité indispensable et revalorisé (antithèse)}
\begin{itemize}
\item A. Nul développement sans techniciens : routes, maisons, énergie. \emph{Exemple :} les chantiers de construction à Douala et Yaoundé.
\item B. Les métiers techniques rapportent parfois plus que les emplois de bureau. \emph{Exemple :} un bon plombier ou électricien gagne mieux sa vie qu'un employé débutant.
\end{itemize}
\emph{Transition :} « Comment alors réconcilier la société avec ses mains ? »

\textbf{III. Synthèse : revaloriser le travail manuel}
\begin{itemize}
\item A. Par l'éducation : valoriser l'enseignement technique dès le collège. \emph{Exemple :} les lycées techniques du Cameroun et leurs lauréats.
\item B. Par l'exemple et la politique : reconnaître les artisans, sécuriser leurs revenus.
\end{itemize}

\textbf{4. Justification des choix :} le plan dialectique convient car le sujet appelle une opinion nuancée ; les exemples sont tirés de la vie courante camerounaise (savoir-faire exigible : appliquer les notions à des situations concrètes) ; la synthèse dépasse l'opposition en proposant des solutions, ce qui donne une conclusion constructive.

\textbf{Conclusion :} le devoir répondra « oui, mais » : le travail manuel est dévalorisé dans les esprits, alors qu'il est la condition du développement ; il faut donc le revaloriser.
\end{exoresolu}

\begin{attention}[Les 5 erreurs qui coûtent des points au Bac]
\begin{enumerate}
\item \textbf{Le hors-sujet} : rédiger avant d'avoir défini les mots-clés et formulé la problématique. Toute partie qui ne répond pas à la problématique est sanctionnée, même bien écrite.
\item \textbf{Le plan déséquilibré ou incomplet} : une partie de dix lignes et une autre de trois pages ; oublier la synthèse dans un plan dialectique ; annoncer trois parties et n'en traiter que deux.
\item \textbf{Les idées sans exemples} : aligner des affirmations générales sans aucun texte, fait ou personnage pour les appuyer. Une idée non illustrée vaut à peine la moitié de sa valeur.
\item \textbf{Les répétitions de référents} : répéter dix fois « Meka » ou « l'auteur » au lieu d'utiliser pronoms, noms génériques et synonymes ; ou, inversement, employer des substituts ambigus qui brouillent la compréhension.
\item \textbf{Les fautes sur les sigles et emprunts} : écrire « les ongs », accorder un sigle invariable, multiplier les anglicismes (\emph{business}, \emph{challenge}) sans équivalent ni guillemets, ce qui donne une impression de négligence linguistique.
\end{enumerate}
Réflexe final : réserver les dix dernières minutes à la relecture (accords, répétitions, cohérence avec le plan).
\end{attention}

\begin{savaistu}[Ferdinand Oyono et la littérature camerounaise]
Ferdinand Oyono (Ngoulemakong, 1929 -- Yaoundé, 2010) est un pilier des lettres camerounaises. Après des études au séminaire de Yaoundé puis en France (Sorbonne), il publie ses trois romans majeurs entre 1956 et 1960. \emph{Le Vieux nègre et la médaille} reçoit le prix Sainte-Beuve 1956. Diplomate, il fut ambassadeur du Cameroun en France, au Royaume-Uni, en Algérie, en Libye, et représentant permanent à l'ONU. Son œuvre, brève mais dense, utilise l'ironie comme arme politique : elle déconstruit le mythe colonial de la "mission civilisatrice" en montrant, de l'intérieur, l'absurdité et la violence du système. Il est, avec Mongo Beti et Ferdinand Léopold Oyono (son homonyme, souvent confondu), l'un des pères du roman camerounais d'expression française.
\end{savaistu}

\newpage

\coursec{Exercices}

\coursec{Leçon 19 : Produire le plan détaillé d'une dissertation}

\courssub{Rappels de cours et méthodes}

\begin{methode}{Analyser un sujet de dissertation}
\begin{enumerate}
\item \textbf{Identifier et souligner les mots-clés} : termes porteurs de sens qui délimitent le sujet (sujet, verbes, adverbes de portée, connecteurs).
\item \textbf{Définir chaque mot-clé} : donner une définition précise, contextuelle et nuancée (dictionnaire, sens commun, sens philosophique/littéraire).
\item \textbf{Dégager l'enjeu (la tension)} : expliquer pourquoi le sujet pose problème (opposition entre deux thèses plausibles, paradoxe, actualité).
\item \textbf{Fixer les limites} : préciser le cadre spatio-temporel, le domaine (littérature, société, histoire) ou les exclusions implicites.
\item \textbf{Formuler la problématique} : transformer le sujet en une question principale (souvent « Dans quelle mesure... ? », « Comment... ? ») et deux questions secondaires qui structurent le débat (thèse / antithèse).
\item \textbf{Annoncer le plan} : indiquer la logique dialectique (Thèse, Antithèse, Synthèse) ou thématique.
\end{enumerate}
\end{methode}

\begin{methode}{Rechercher les idées (Remue-méninges et tri)}
\begin{enumerate}
\item \textbf{Noter sans trier} : sur une feuille de brouillard, écrire toutes les idées, arguments, exemples, citations, références littéraires (Le Vieux nègre et la médaille, autres œuvres du programme), faits sociaux (Cameroun), historiques qui viennent à l'esprit.
\item \textbf{Classifier} : regrouper les idées par affinité (pour/contre, causes/conséquences, individuel/collectif, court terme/long terme).
\item \textbf{Sélectionner et hiérarchiser} : retenir les arguments les plus forts, les plus originaux, ceux qui s'illustrent par des exemples précis. Écarter les lieux communs, les hors-sujet, les redites.
\item \textbf{Structurer en trois parties} : répartir les arguments retenus dans l'architecture Thèse / Antithèse / Synthèse (ou Autre structure / Autre structure / Dépassement).
\end{enumerate}
\end{methode}

\begin{methode}{Élaborer un plan détaillé dialectique}
\begin{enumerate}
\item \textbf{Partie I (Thèse) :} Affirmer la réponse la plus évidente ou la plus courante au sujet.
  \begin{itemize}
  \item Argument 1 + Exemple précis (littéraire, historique, vie courante camerounaise).
  \item Argument 2 + Exemple précis.
  \item (Argument 3 + Exemple si possible).
  \end{itemize}
\item \textbf{Partie II (Antithèse) :} Nuancer, contredire ou complexifier la thèse en apportant le point de vue inverse ou les limites.
  \begin{itemize}
  \item Argument 1 + Exemple précis.
  \item Argument 2 + Exemple précis.
  \end{itemize}
\item \textbf{Partie III (Synthèse) :} Dépasser l'opposition en proposant une réponse nuancée, conditionnelle, ou une réconciliation des deux pôles (souvent : « Oui, mais à condition que... », « Ni l'un ni l'autre, mais... »).
  \begin{itemize}
  \item Argument 1 (dépassement) + Exemple précis.
  \item Argument 2 (ouverture) + Exemple précis.
  \end{itemize}
\item \textbf{Rédiger les transitions} : une phrase de liaison entre chaque partie (ex: « Cependant, cette vision idéale se heurte à la réalité... » ; « Au terme de cette analyse, il apparaît que... »).
\end{enumerate}
\end{methode}

\begin{definition}{Problématique}
Question centrale, issue de l'analyse du sujet, qui met en tension au moins deux réponses possibles (thèse et antithèse) et guide tout le développement de la dissertation. Elle ne se contente pas de répéter le sujet.
\end{definition}

\begin{definition}{Plan détaillé}
Squelette complet de la dissertation : pour chaque grande partie (I, II, III), on indique les arguments (sous-parties) et, pour chaque argument, l'exemple précis qui l'illustre. Il précède la rédaction de l'introduction et du développement.
\end{definition}

\begin{aretenir}{Les substituts du référent (Cohésion textuelle)}
Le \cle{référent} est l'entité (personne, objet, idée) désignée la première fois par un nom propre ou un groupe nominal. Les \cle{substituts} (pronoms personnels, possessifs, démonstratifs, relatifs ; périphrases nominales ; synonymes) reprennent ce référent pour éviter les répétitions et enrichir le portrait (ex: Meka $\rightarrow$ le vieux nègre, le futur médaillé, lui, son). Maîtriser la chaîne de référence est indispensable pour une rédaction fluide et élégante.
\end{aretenir}

\begin{aretenir}{Sigles et emprunts : usage modéré en examen}
\begin{itemize}
\item \cle{Sigle} : sigle = ensemble de lettres initiales prononcé lettre par lettre (ex: \emph{MINESEC}, \emph{ONG}, \emph{CNUCED}, \emph{BEPC}). À définir à la première occurrence.
\item \cle{Acronyme} : sigle prononcé comme un mot (ex: \emph{UNESCO}, \emph{NASA}, \emph{OMS}). À définir aussi.
\item \cle{Emprunt} : mot venu d'une langue étrangère (anglicisme : \emph{week-end}, \emph{business}, \emph{football} ; camfranglais : \emph{tchop}, \emph{djo}). En copie d'examen, privilégier l'équivalent français soutenu (\emph{fin de semaine}, \emph{affaires}, \emph{football} $\rightarrow$ ballon rond / sport ; \emph{tchop} $\rightarrow$ repas, nourriture).
\end{itemize}
\end{aretenir}

\courssub{Je vérifie mes connaissances}

\begin{exercice}[QCM : L'analyse du sujet]
Pour chaque question, choisis la bonne réponse.
\begin{enumerate}
\item La première étape du travail sur un sujet de dissertation est :
\begin{enumerate}
\item[a)] rédiger l'introduction ;
\item[b)] définir les mots-clés du sujet ;
\item[c)] chercher des exemples littéraires.
\end{enumerate}
\item Dans le sujet « Le travail est-il toujours source de dignité ? », le mot-clé qui appelle une définition prioritaire est :
\begin{enumerate}
\item[a)] « toujours » ; \item[b)] « dignité » ; \item[c)] « est-il ».
\end{enumerate}
\item La problématique se présente généralement sous la forme :
\begin{enumerate}
\item[a)] d'une citation ; \item[b)] d'une ou plusieurs questions ; \item[c)] d'un résumé du sujet.
\end{enumerate}
\item Un plan dialectique comporte :
\begin{enumerate}
\item[a)] deux parties ; \item[b)] trois parties (thèse, antithèse, synthèse) ; \item[c)] quatre parties.
\end{enumerate}
\end{enumerate}
\end{exercice}

\begin{exercice}[Vrai ou faux — Justifier]
Dis si chaque affirmation est vraie ou fausse, puis justifie ta réponse en une ou deux phrases.
\begin{enumerate}
\item Le paratexte d'une œuvre comprend le titre, le nom de l'auteur, la couverture et la quatrième de couverture.
\item Dans une dissertation, chaque partie du plan détaillé doit contenir au moins deux arguments illustrés chacun d'un exemple.
\item Un sigle est un mot emprunté à une langue étrangère.
\item Dans la phrase « Meka reçut sa médaille. Le vieil homme la serra contre son cœur », « le vieil homme » et « son » sont des substituts du référent « Meka ».
\end{enumerate}
\end{exercice}

\begin{exercice}[Définitions]
Définis en deux ou trois phrases chacune des notions suivantes, telles qu'elles s'appliquent à la dissertation :
\begin{enumerate}
\item le mot-clé ;
\item la problématique ;
\item le plan détaillé ;
\item le référent et le substitut.
\end{enumerate}
\end{exercice}

\begin{exercice}[Sigles et emprunts]
\begin{enumerate}
\item Donne la signification complète des sigles suivants : \emph{MINESEC}, \emph{CMR} (code ISO du Cameroun), \emph{ONG}, \emph{BAC}.
\item Classe les mots suivants en deux colonnes (Sigles / Emprunts) : \emph{week-end}, \emph{CNUCED}, \emph{football}, \emph{BEPC}, \emph{business}, \emph{tchop} (emprunt local camfranglais).
\item Pourquoi doit-on limiter l'emploi des sigles et des emprunts dans une copie d'examen ?
\end{enumerate}
\end{exercice}

\courssub{Je m'entraîne}

\begin{exercice}[Analyser un sujet]
On donne le sujet : \emph{« L'école est-elle le seul chemin de la réussite ? »}
\begin{enumerate}
\item Souligne les mots-clés et définis chacun d'eux en une phrase.
\item Dégage l'enjeu du sujet (ce qui pose problème).
\item Formule la problématique sous forme d'une question principale et de deux questions secondaires.
\end{enumerate}
\end{exercice}

\begin{exercice}[Recherche des idées]
À partir du sujet : \emph{« Faut-il honorer les anciens combattants ? »}
\begin{enumerate}
\item Par la technique du remue-méninges, note au moins huit idées (arguments ou exemples) liées au sujet.
\item Classe ces idées en trois groupes correspondant aux trois parties d'un plan dialectique (Thèse, Antithèse, Synthèse).
\item Rédige la problématique du sujet.
\end{enumerate}
\end{exercice}

\begin{exercice}[Chaîne de référence]
Lis l'extrait suivant, inspiré de la situation du \emph{Vieux nègre et la médaille} de Ferdinand Oyono :

\emph{« Meka attendait depuis l'aube. Le vieux nègre avait revêtu son plus beau pagne. Le futur médaillé serrait contre sa poitrine la lettre du commandant. Lui, si longtemps méprisé, allait enfin être honoré devant tout le village. »}

\begin{enumerate}
\item Relève tous les substituts qui désignent Meka.
\item Indique la nature de chaque substitut (nom, pronom, périphrase).
\item Explique l'effet produit par le choix de ces substituts sur le portrait du personnage.
\end{enumerate}
\end{exercice}

\begin{exercice}[Corriger un plan défaillant]
Un élève propose le plan suivant pour le sujet : \emph{« Le sport est-il utile à la jeunesse camerounaise ? »}

I. L'histoire du football au Cameroun. — II. Le sport fatigue les élèves et les empêche d'étudier. — III. Le sport est utile à la jeunesse.

\begin{enumerate}
\item Relève trois erreurs commises par cet élève (pertinence, équilibre, logique du plan).
\item Propose un plan dialectique corrigé, avec pour chaque partie deux arguments et un exemple.
\item Justifie tes corrections en deux ou trois phrases.
\end{enumerate}
\end{exercice}

\courssub{Je me prépare au Bac}

\begin{exercice}[Sujet type Baccalauréat : Le travail et la dignité]
Sujet : \emph{« Le travail est-il toujours source de dignité ? »}
\begin{enumerate}
\item Analyse le sujet : définis les mots-clés, dégage l'enjeu et les limites éventuelles du sujet.
\item Formule la problématique (une question principale et deux questions secondaires).
\item Recherche les idées : note au moins six arguments et six exemples (vie quotidienne camerounaise, littérature, histoire).
\item Produis le plan détaillé d'un plan dialectique en trois parties, chaque partie comportant au moins deux arguments illustrés.
\item Rédige l'introduction complète de la dissertation (amène-sujet, présentation du sujet, problématique, annonce du plan).
\end{enumerate}
\end{exercice}

\begin{exercice}[Récompenses officielles et motivation des citoyens]
Sujet de discussion : \emph{« Les récompenses officielles (médailles, distinctions) motivent-elles réellement les citoyens camerounais ? »}

Tu t'appuieras sur ta lecture du \emph{Vieux nègre et la médaille} de Ferdinand Oyono et sur tes observations de la vie sociale au Cameroun.
\begin{enumerate}
\item Dégage la thèse (les récompenses motivent), l'antithèse (elles déçoivent ou humilient parfois) et la synthèse possible.
\item Formule la problématique.
\item Produis le plan détaillé complet : trois parties, chacune avec deux arguments et un exemple précis (dont un tiré de l'œuvre étudiée).
\item Rédige l'introduction de cette dissertation en veillant à éviter toute répétition grâce aux substituts du référent.
\end{enumerate}
\end{exercice}

\begin{exercice}[Question de synthèse type Probatoire : Compte rendu de lecture]
Lis attentivement l'extrait suivant :

\emph{« Le jour de la cérémonie, Meka se leva avant le soleil. Il enfila son pagne neuf, chaussa ses vieux souliers cirés et prit le chemin de la place du village. Toute la journée, sous un soleil de feu, il attendit. Les discours se succédaient, les fanfares jouaient, mais son nom ne venait pas. Quand enfin l'administrateur l'appela, le vieil homme s'avança, les jambes tremblantes, le cœur battant. On lui épingla la médaille sur la poitrine. La foule applaudit. Pourtant, au fond de lui, une étrange tristesse grandissait : cette médaille effaçait-elle vraiment des années de mépris ? »}

\begin{enumerate}
\item Rédige un compte rendu de lecture de quinze lignes maximum : tu présenteras la situation, le personnage et le sens de l'extrait, sans recopier le texte.
\item Dans ton compte rendu, utilise au moins quatre substituts différents pour désigner Meka et souligne-les.
\item Relève dans l'extrait un sigle ou un emprunt éventuel ; si le passage n'en contient pas, cite un emprunt courant dans le français camerounais et propose son équivalent soutenu pour une copie d'examen.
\item Dégage en deux phrases la leçon que tu tires de cet extrait pour répondre à la question : la médaille honore-t-elle vraiment celui qui la reçoit ?
\end{enumerate}
\end{exercice}

\newpage

\coursec{Corrigés des exercices}

\begin{corrige}[Exercice 1 — QCM : L'analyse du sujet]
\begin{enumerate}
\item \textbf{Réponse b)} : définir les \cle{mots-clés} du sujet. Toute dissertation commence par l'analyse du sujet : on identifie et on définit les termes porteurs de sens avant toute recherche d'exemples ou rédaction. Rédiger l'introduction sans analyse conduit au hors-sujet.
\item \textbf{Réponse b)} : « \cle{dignité} ». C'est le terme abstrait, polysémique (respect de soi, honneur, valeur morale), dont la définition conditionne toute la réflexion. « Toujours » est un adverbe de portée important (il installe la tension), mais c'est « dignité » qui exige la définition prioritaire.
\item \textbf{Réponse b)} : d'une ou plusieurs questions. La \cle{problématique} transforme le sujet en interrogation principale (souvent introduite par « Dans quelle mesure... ? »), accompagnée de questions secondaires qui annoncent le débat.
\item \textbf{Réponse b)} : trois parties (thèse, antithèse, synthèse). Le \cle{plan dialectique}, privilégié au Baccalauréat technique, oppose deux points de vue puis les dépasse.
\end{enumerate}
\end{corrige}

\begin{corrige}[Exercice 2 — Vrai ou faux]
\begin{enumerate}
\item \textbf{Vrai.} Le \cle{paratexte} (du grec \emph{para}, « à côté de ») désigne tout ce qui entoure le texte sans en faire partie : titre, nom de l'auteur, couverture, quatrième de couverture, préface, dédicace. Son étude prépare la lecture du \cle{texte ouvroir}.
\item \textbf{Vrai.} Un \cle{plan détaillé} exige, pour chaque grande partie, au moins deux arguments, chacun illustré d'un exemple précis (littéraire, historique ou tiré de la vie courante). Un argument sans exemple reste une affirmation gratuite ; une partie à un seul argument est déséquilibrée.
\item \textbf{Faux.} Un \cle{sigle} est formé des initiales d'un groupe de mots (\emph{ONG} = organisation non gouvernementale) ; un \cle{emprunt} est un mot venu d'une langue étrangère (\emph{week-end}, venu de l'anglais). Ce sont deux procédés lexicaux distincts.
\item \textbf{Vrai.} « Meka » est le \cle{référent}, nommé une seule fois. « Le vieil homme » est une \cle{périphrase nominale} (substitut nominal) et « son » est un \cle{déterminant possessif} qui renvoie à lui : tous deux entretiennent la \cle{chaîne de référence} et évitent la répétition du nom propre.
\end{enumerate}
\end{corrige}

\begin{corrige}[Exercice 3 — Définitions]
\begin{enumerate}
\item \textbf{Le \cle{mot-clé}} : terme du sujet qui porte l'essentiel du sens et délimite le champ de la réflexion (notion abstraite, verbe d'action, adverbe de portée comme « toujours », « seul »). Il doit être défini avec précision, dans le contexte du sujet, avant toute recherche d'idées.
\item \textbf{La \cle{problématique}} : question centrale qui naît de l'analyse du sujet et met en tension au moins deux réponses possibles. Elle ne répète pas le sujet : elle le transforme en véritable problème et guide l'ensemble du développement.
\item \textbf{Le \cle{plan détaillé}} : squelette complet de la dissertation, établi au brouillon avant la rédaction. Pour chaque grande partie (I, II, III), il indique les arguments et, pour chaque argument, l'exemple précis qui l'illustre, ainsi que les transitions.
\item \textbf{Le \cle{référent} et le \cle{substitut}} : le référent est l'entité (personne, objet, idée) désignée une première fois par un nom propre ou un groupe nominal ; les substituts (pronoms, périphrases, synonymes) reprennent ce référent dans la suite du texte pour assurer la \cle{cohésion} et éviter les répétitions.
\end{enumerate}
\end{corrige}

\begin{corrige}[Exercice 4 — Sigles et emprunts]
\begin{enumerate}
\item Significations :
\begin{itemize}
\item \emph{MINESEC} : Ministère des Enseignements secondaires (Cameroun) ;
\item \emph{CMR} : Cameroun (code ISO à trois lettres du pays) ;
\item \emph{ONG} : organisation non gouvernementale ;
\item \emph{BAC} : baccalauréat.
\end{itemize}
\item Classement :
\begin{center}
\begin{tabularx}{\linewidth}{|X|X|}
\hline
\rowcolor{popblueL} \textbf{Sigles} & \textbf{Emprunts} \\
\hline
CNUCED, BEPC & week-end, football, business, tchop \\
\hline
\end{tabularx}
\end{center}
\item En copie d'examen, le correcteur attend un français clair et soutenu. Les sigles non définis à la première occurrence peuvent être incompréhensibles ; les emprunts (anglicismes, camfranglais) appartiennent souvent à la langue familière ou parlée et affaiblissent le niveau de langue. On les remplace donc par des équivalents français : \emph{week-end} $\rightarrow$ fin de semaine ; \emph{business} $\rightarrow$ affaires ; \emph{tchop} $\rightarrow$ repas, nourriture.
\end{enumerate}
\end{corrige}

\begin{corrige}[Exercice 5 — Analyser un sujet]
Sujet : \emph{« L'école est-elle le seul chemin de la réussite ? »}
\begin{enumerate}
\item Mots-clés et définitions :
\begin{itemize}
\item \emph{L'école} : institution d'enseignement formel, de l'école primaire à l'université, qui transmet des savoirs et délivre des diplômes.
\item \emph{Le seul} : adverbe de portée exclusif ; il exclut toute autre voie possible et constitue le nœud du problème.
\item \emph{Chemin} : au figuré, moyen, voie d'accès.
\item \emph{La réussite} : accomplissement social, professionnel, économique ou personnel d'un individu.
\end{itemize}
\item \textbf{Enjeu} : dans nos sociétés camerounaises, le diplôme scolaire apparaît comme le passeport obligé vers l'emploi et la considération sociale ; pourtant, de nombreux artisans, commerçants, agriculteurs ou artistes réussissent sans longs parcours scolaires. Le sujet oppose donc la valeur sociale du diplôme aux autres formes d'ascension sociale.
\item \textbf{Problématique} :
\begin{itemize}
\item Question principale : dans quelle mesure l'école constitue-t-elle la voie privilégiée, mais non exclusive, de la réussite ?
\item Question secondaire 1 : pourquoi l'école reste-t-elle, au Cameroun, la voie la plus sûre et la plus valorisée vers la réussite ?
\item Question secondaire 2 : quelles autres voies (formation professionnelle, apprentissage, entrepreneuriat, talent) peuvent y conduire, et à quelles conditions ?
\end{itemize}
\end{enumerate}
\end{corrige}

\begin{corrige}[Exercice 6 — Recherche des idées]
Sujet : \emph{« Faut-il honorer les anciens combattants ? »}
\begin{enumerate}
\item Remue-méninges (huit idées possibles) :
\begin{itemize}
\item Ils ont sacrifié leur jeunesse et risqué leur vie pour la patrie ;
\item La médaille est une reconnaissance symbolique de la nation ;
\item Meka, dans \emph{Le Vieux nègre et la médaille}, attend sa médaille avec fierté ;
\item Certaines récompenses arrivent trop tard ou sont dérisoires ;
\item La médaille de Meka n'efface pas le mépris colonial ;
\item Beaucoup d'anciens combattants vivent dans la pauvreté, sans pension ;
\item Honorer les anciens combattants transmet la mémoire aux jeunes générations ;
\item Les honneurs doivent s'accompagner d'actes concrets (pensions, soins, logement).
\end{itemize}
\item Classement :
\begin{center}
\begin{tabularx}{\linewidth}{|X|X|X|}
\hline
\rowcolor{popblueL} \textbf{Thèse (il faut les honorer)} & \textbf{Antithèse (limites des honneurs)} & \textbf{Synthèse} \\
\hline
Sacrifice pour la patrie ; transmission de la mémoire ; fierté de Meka & Récompenses tardives ou dérisoires ; pauvreté des vétérans ; la médaille n'efface pas le mépris & Honneurs sincères accompagnés d'actes concrets (pensions, reconnaissance réelle) \\
\hline
\end{tabularx}
\end{center}
\item \textbf{Problématique} : dans quelle mesure les honneurs rendus aux anciens combattants constituent-ils une juste reconnaissance de leur sacrifice, et comment éviter qu'ils ne soient qu'une cérémonie vide ?
\end{enumerate}
\end{corrige}

\begin{corrige}[Exercice 7 — Chaîne de référence]
\begin{enumerate}
\item Substituts désignant Meka : \emph{le vieux nègre}, \emph{son} (plus beau pagne), \emph{le futur médaillé}, \emph{sa} (poitrine), \emph{lui}, \emph{si longtemps méprisé}.
\item Nature de chaque substitut :
\begin{center}
\begin{tabularx}{\linewidth}{|l|X|}
\hline
\rowcolor{popblueL} \textbf{Substitut} & \textbf{Nature} \\
\hline
le vieux nègre & périphrase nominale (groupe nominal) \\
le futur médaillé & périphrase nominale (groupe nominal) \\
lui & pronom personnel tonique \\
son, sa & déterminants possessifs \\
si longtemps méprisé & participe passé en apposition (reprise qualifiante) \\
\hline
\end{tabularx}
\end{center}
\item \textbf{Effet produit} : les périphrases dressent un portrait en deux temps : « le vieux nègre » rappelle l'âge et la condition coloniale du personnage, tandis que « le futur médaillé » projette l'attente glorieuse qui anime sa journée. L'apposition « si longtemps méprisé » crée un contraste pathétique entre l'humiliation passée et l'honneur espéré : le lecteur éprouve de la sympathie pour Meka et pressent l'ironie tragique de la récompense.
\end{enumerate}
\end{corrige}

\begin{corrige}[Exercice 8 — Corriger un plan défaillant]
\begin{enumerate}
\item Trois erreurs :
\begin{itemize}
\item \textbf{Hors-sujet (pertinence)} : la partie I (« L'histoire du football au Cameroun ») ne répond pas à la question posée ; elle raconte au lieu d'argumenter.
\item \textbf{Déséquilibre} : la partie II contient une seule idée, sans exemple, et la partie III est une simple affirmation sans argument.
\item \textbf{Logique défaillante} : l'ordre est incohérent — la thèse favorable devrait ouvrir le devoir, la nuance suivre, et la synthèse conclure ; ici la réponse positive arrive en dernier sans dépassement.
\end{itemize}
\item Plan dialectique corrigé :
\begin{itemize}
\item \textbf{I. Le sport est utile à la jeunesse camerounaise (thèse).} Arg. 1 : il assure la santé physique et mentale des jeunes (ex. : les séances d'EPS et les championnats scolaires luttent contre l'obésité et le stress des examens). Arg. 2 : il forge le caractère et la discipline (ex. : l'entraînement régulier des jeunes footballeurs des académies de Douala apprend la ponctualité et l'esprit d'équipe).
\item \textbf{II. Mais le sport peut aussi détourner la jeunesse de ses priorités (antithèse).} Arg. 1 : pratiqué sans encadrement, il empiète sur le temps d'étude (ex. : des élèves qui abandonnent les classes pour les paris et les matchs de quartier échouent au BEPC). Arg. 2 : le rêve de la carrière professionnelle fait oublier que peu y accèdent (ex. : des milliers de jeunes candidats pour quelques places dans les clubs professionnels).
\item \textbf{III. Le sport est utile à condition d'être équilibré avec les études (synthèse).} Arg. 1 : l'école doit organiser une pratique encadrée et compatible avec les apprentissages (ex. : les sport-études qui concilient entraînement et scolarité). Arg. 2 : le sport peut devenir un complément de la formation et même un métier (ex. : les carrières d'éducateur sportif, d'arbitre ou de gestionnaire d'installations).
\end{itemize}
\item \textbf{Justification} : le plan corrigé répond directement à la question posée, respecte la progression dialectique (affirmation, nuance, dépassement) et garantit l'équilibre des parties : chacune contient deux arguments étayés par un exemple précis et camerounais.
\end{enumerate}
\end{corrige}

\begin{corrige}[Exercice 9 — Sujet type Baccalauréat : Le travail et la dignité]
\begin{enumerate}
\item \textbf{Analyse du sujet.}
\begin{itemize}
\item \emph{Le travail} : activité humaine, manuelle ou intellectuelle, rémunérée ou non, par laquelle l'homme transforme le monde et subvient à ses besoins.
\item \emph{Toujours} : adverbe de portée absolue ; il interdit toute exception et installe la tension du sujet.
\item \emph{Source de dignité} : ce qui confère à l'homme sa valeur, le respect de soi et la considération d'autrui.
\item \textbf{Enjeu} : on dit communément que « le travail ennoblit l'homme » ; pourtant, certains travaux (exploitants, mal payés, dégradants) semblent au contraire avilir. Le sujet oppose l'idéal du travail valorisant à la réalité du travail aliénant.
\item \textbf{Limites} : on se placera dans le cadre de la société camerounaise contemporaine, en convoquant la littérature et l'histoire.
\end{itemize}
\item \textbf{Problématique.} Question principale : dans quelle mesure le travail confère-t-il la dignité, et quelles conditions doit-il remplir pour ne pas devenir source d'humiliation ? Questions secondaires : en quoi le travail est-il traditionnellement valorisant ? Quelles formes de travail, au contraire, dégradent l'homme ?
\item \textbf{Recherche des idées} (arguments + exemples) :
\begin{itemize}
\item Le travail donne l'indépendance économique (ex. : la commerçante du marché Mokolo qui scolarise ses enfants) ;
\item Il vaut considération sociale (ex. : l'artisan menuisier respecté de son quartier) ;
\item L'oisiveté est méprisée dans la tradition (ex. : proverbes valorisant l'effort) ;
\item Certains travaux sont épuisants et mal payés (ex. : les journaliers des plantations) ;
\item Le chômage des diplômés pousse à des « métiers de débrouille » dévalorisés (ex. : les jeunes diplômés devenus moto-taximen malgré eux) ;
\item Meka a travaillé toute sa vie sans être respecté par l'administration coloniale (ex. : \emph{Le Vieux nègre et la médaille}) ;
\item Le travail librement choisi et justement rémunéré dignifie (ex. : les coopératives agricoles de l'Ouest) ;
\item La dignité dépend des conditions, non du seul fait de travailler (ex. : luttes ouvrières pour de meilleurs salaires).
\end{itemize}
\item \textbf{Plan détaillé.}
\begin{itemize}
\item \textbf{I. Le travail est source de dignité (thèse).} Arg. 1 : il procure l'autonomie et le respect de soi (ex. : la commerçante qui subvient seule aux besoins de sa famille). Arg. 2 : il confère l'estime de la communauté (ex. : l'artisan ou l'enseignant respecté ; les proverbes qui condamnent la paresse).
\item \textbf{II. Mais tout travail ne dignifie pas (antithèse).} Arg. 1 : le travail exploité et mal rémunéré avilit (ex. : les journaliers sous-payés des plantations ; les ouvriers sans contrat). Arg. 2 : le travail subi sous la domination n'apporte ni reconnaissance ni honneur (ex. : Meka, corvéable et méprisé, dont la médaille ne compense pas une vie d'humiliations dans \emph{Le Vieux nègre et la médaille}).
\item \textbf{III. Le travail dignifie à condition d'être libre, juste et reconnu (synthèse).} Arg. 1 : la dignité naît des conditions du travail (salaire juste, respect, sécurité) (ex. : les coopératives où le producteur maîtrise le fruit de son effort). Arg. 2 : la société doit reconnaître la valeur de tous les métiers (ex. : la revalorisation sociale des métiers manuels et techniques au Cameroun).
\end{itemize}
\item \textbf{Introduction (modèle).} \\
Dans nos sociétés, l'homme se définit largement par son activité : « Dis-moi ce que tu fais, je te dirai qui tu es », semble dire le sens commun. Le travail apparaît ainsi, depuis toujours, comme le fondement de la valeur de l'individu : il nourrit, il construit, il honore. Pourtant, l'expérience quotidienne montre aussi des travailleurs épuisés, mal payés, méprisés, à l'image de Meka, le héros de Ferdinand Oyono, dont une vie de labeur n'a jamais effacé l'humiliation. Dès lors, une question s'impose : le travail est-il toujours source de dignité ? Autrement dit, le fait de travailler suffit-il à rendre l'homme respectable, ou la dignité dépend-elle des conditions dans lesquelles s'exerce l'effort ? Nous verrons d'abord que le travail fonde traditionnellement la dignité de l'homme ; nous montrerons ensuite que certains travaux, au contraire, dégradent celui qui les accomplit ; nous conclurons enfin que le travail n'ennoblit qu'à condition d'être libre, juste et reconnu.
\end{enumerate}

\textbf{Barème indicatif (sur 20)} : analyse du sujet et mots-clés (4 pts) ; problématique pertinente (3 pts) ; richesse et pertinence des idées (4 pts) ; plan détaillé équilibré et dialectique (5 pts) ; introduction complète et bien rédigée (4 pts).

\textbf{Points de vigilance} : définir « toujours » et « dignité » avant tout ; ne pas confondre « travail » et « emploi salarié » ; chaque argument doit être suivi d'un exemple précis ; la synthèse ne doit pas être une simple conclusion mais un dépassement (« oui, mais à condition que... ») ; soigner les transitions entre les parties.
\end{corrige}

\begin{corrige}[Exercice 10 — Récompenses officielles et motivation des citoyens]
\begin{enumerate}
\item \textbf{Dégagement des trois moments.}
\begin{itemize}
\item \emph{Thèse} : les récompenses officielles motivent les citoyens : elles reconnaissent le mérite, stimulent l'émulation et donnent l'exemple (la fierté de Meka attendant sa médaille ; les décorations du 20 mai).
\item \emph{Antithèse} : elles déçoivent ou humilient parfois : récompenses tardives, dérisoires ou injustement distribuées (la médaille de Meka n'efface pas des années de mépris ; les distinctions parfois accordées au favoritisme).
\item \emph{Synthèse} : les récompenses ne motivent durablement que si elles sont méritées, sincères et accompagnées d'une reconnaissance concrète.
\end{itemize}
\item \textbf{Problématique} : dans quelle mesure les récompenses officielles constituent-elles un véritable moteur de l'engagement citoyen, et à quelles conditions échappent-elles à la dérision ?
\item \textbf{Plan détaillé.}
\begin{itemize}
\item \textbf{I. Les récompenses officielles motivent les citoyens (thèse).} Arg. 1 : elles consacrent publiquement le mérite et comblent le besoin de reconnaissance (ex. : Meka, dans \emph{Le Vieux nègre et la médaille}, se lève avant l'aube, revêt son plus beau pagne, tant l'honneur attendu donne un sens à sa vie). Arg. 2 : elles créent l'émulation et proposent des modèles à la jeunesse (ex. : les cérémonies de décoration de la fête nationale qui célèbrent fonctionnaires, agriculteurs et sportifs méritants).
\item \textbf{II. Mais elles peuvent décevoir, voire humilier (antithèse).} Arg. 1 : une récompense tardive ou dérisoire ne répare pas les souffrances vécues (ex. : la médaille de Meka, épinglée après une vie de corvées, laisse en lui « une étrange tristesse »). Arg. 2 : distribuées sans justice, elles perdent toute valeur et découragent (ex. : les distinctions obtenues par favoritisme ou relations, qui découragent les vrais travailleurs).
\item \textbf{III. Les récompenses ne motivent que si elles sont méritées et sincères (synthèse).} Arg. 1 : la distinction doit couronner un mérite réel, vérifié par tous (ex. : un instituteur de village décoré devant sa communauté, dont tous connaissent le dévouement). Arg. 2 : elle doit s'accompagner d'une reconnaissance concrète (pension, promotion, moyens) (ex. : les primes et avantages accordés aux lauréats des grands concours nationaux).
\end{itemize}
\item \textbf{Introduction (modèle, avec substituts mis en valeur).} \\
Chaque année, lors des fêtes nationales, le Cameroun décore des centaines de ses fils et filles. \emph{Le Vieux nègre et la médaille} de Ferdinand Oyono offre de ces cérémonies une image inoubliable : Meka, le héros du roman, attend \emph{sa} distinction comme une délivrance. \emph{Ce vieil homme}, si longtemps méprisé, croit un instant que \emph{le ruban épinglé sur sa poitrine} effacera une vie d'humiliations ; \emph{il} découvrira pourtant que l'honneur officiel ne guérit pas toutes les blessures. Dès lors, on peut se demander : les récompenses officielles motivent-elles réellement les citoyens camerounais ? Nous montrerons d'abord qu'elles stimulent l'engagement en consacrant le mérite ; nous verrons ensuite qu'elles peuvent décevoir et même humilier ; nous soutiendrons enfin qu'elles ne motivent durablement qu'à condition d'être méritées, sincères et concrètes.
\end{enumerate}

\textbf{Barème indicatif (sur 20)} : dégagement thèse/antithèse/synthèse (4 pts) ; problématique (3 pts) ; plan détaillé avec exemples dont un tiré de l'œuvre (7 pts) ; introduction avec chaîne de substituts correcte (6 pts).

\textbf{Points de vigilance} : citer précisément l'œuvre au programme (au moins un exemple tiré de Meka) ; vérifier que chaque substitut renvoie sans ambiguïté au référent ; éviter de répéter « récompense » à chaque phrase (utiliser « distinction », « honneur », « elle ») ; la synthèse doit formuler une condition, pas un compromis mou.
\end{corrige}

\begin{corrige}[Exercice 11 — Compte rendu de lecture type Probatoire]
\begin{enumerate}
\item \textbf{Modèle de compte rendu (quinze lignes maximum), substituts mis en valeur.} \\
L'extrait présente le jour tant attendu de la remise de la médaille. Meka, vieux villageois africain, s'est préparé avec soin dès l'aube : \emph{le futur médaillé} a revêtu son pagne neuf et ciré ses vieux souliers. \emph{Le vieil homme} endure ensuite une longue attente sous le soleil, au milieu des discours et des fanfares, sans que \emph{son} nom soit appelé. Lorsque l'administrateur le convie enfin, \emph{il} s'avance tremblant et reçoit la décoration sous les applaudissements. Pourtant, la joie espérée ne vient pas : \emph{ce serviteur si longtemps méprisé} éprouve une tristesse secrète, car un ruban de métal ne saurait effacer des années d'humiliation. Le passage révèle ainsi l'ironie amère de la récompense coloniale : l'honneur officiel, accordé tardivement et avec condescendance, ne rend pas à l'homme la dignité qu'on lui a refusée toute sa vie. L'extrait illustre le thème central du roman : le fossé entre les apparences de la reconnaissance et la réalité du mépris.
\item Les quatre substituts mis en valeur dans le compte rendu sont : \emph{le futur médaillé} (périphrase), \emph{le vieil homme} (périphrase), \emph{son} (déterminant possessif), \emph{il} (pronom personnel), auxquels s'ajoute \emph{ce serviteur si longtemps méprisé} (périphrase péjorative). Ils évitent la répétition du nom « Meka » et nuancent le portrait (attente, âge, souffrance).
\item L'extrait ne contient ni sigle ni emprunt. Un emprunt courant dans le français camerounais est \emph{tchop} (camfranglais, de l'anglais \emph{chop}) ; en copie d'examen, on lui substituera l'équivalent soutenu « repas » ou « nourriture » (ex. : « Meka partagea le repas de fête » et non « le tchop »).
\item \textbf{Leçon dégagée} : la médaille n'honore vraiment que si elle s'accompagne d'un respect réel et constant ; décernée après une vie de mépris, elle demeure un geste vide qui ne répare ni la dignité blessée ni les années perdues. La véritable reconnaissance se mesure aux actes quotidiens, non aux cérémonies.
\end{enumerate}

\textbf{Barème indicatif (sur 20)} : compte rendu fidèle, sans recopie, dans la limite de lignes (8 pts) ; quatre substituts variés et correctement employés (4 pts) ; question sur l'emprunt et équivalent soutenu (3 pts) ; leçon dégagée en deux phrases claires (5 pts).

\textbf{Points de vigilance} : un compte rendu reformule, il ne recopie jamais le texte ; respecter la limite des quinze lignes ; vérifier que chaque substitut se rattache sans ambiguïté à Meka ; la leçon finale doit répondre explicitement à la question posée (« la médaille honore-t-elle vraiment... ? »).
\end{corrige}

\newpage

\coursec{Fiche bilan}

\begin{popfigure}
\begin{tikzpicture}[mindmap, grow cyclic, every node/.style={concept, font=\small},
concept color=popblue,
level 1/.append style={level distance=3.4cm, sibling angle=60},
level 2/.append style={level distance=2.2cm, font=\scriptsize}]
\node[concept, text=white, font=\bfseries\small] {Dissertation : plan détaillé}
child[concept color=poporange] { node {Analyser le sujet}
  child { node {Mots-clés} }
  child { node {Type de sujet} }
}
child[concept color=popgreen] { node {Rechercher les idées}
  child { node {Brainstorming} }
  child { node {Trier, regrouper} }
}
child[concept color=poppurple] { node {Problématique}
  child { node {Question directrice} }
}
child[concept color=poppink] { node {Plan détaillé}
  child { node {Intro, 2 ou 3 parties, conclusion} }
  child { node {Arguments + exemples} }
}
child[concept color=popgold] { node {Référent et substituts}
  child { node {Pronoms, reprises} }
}
child[concept color=popdark] { node {Sigles et emprunts} };
\end{tikzpicture}
\legende{Carte mentale de la leçon : les six compétences clés pour produire le plan détaillé d'une dissertation.}
\end{popfigure}

\begin{bilan}
\textbf{Définitions clés}
\begin{itemize}
\item \cle{Dissertation} : exercice écrit de réflexion personnelle et argumentée sur un sujet donné, en introduction, développement (2 ou 3 parties) et conclusion.
\item \cle{Analyse du sujet} : dégager les \cle{mots-clés}, définir les termes, repérer les limites (temps, espace, domaine) et le type de sujet (de réflexion, de discussion, de démonstration).
\item \cle{Problématique} : question directrice, issue du sujet, à laquelle tout le devoir répond.
\item \cle{Plan détaillé} : schéma complet du devoir indiquant, pour chaque partie, l'idée directrice, les arguments et les exemples.
\item \cle{Référent} : réalité ou mot auquel renvoie un terme ; \cle{substituts} : pronoms, synonymes, périphrases qui le reprennent.
\item \cle{Sigle} : abréviation formée d'initiales (ONU, MINESEC) ; \cle{emprunt} : mot venant d'une autre langue (week-end, ndolè).
\end{itemize}

\textbf{Méthode express : du sujet au plan détaillé}
\begin{center}
\begin{tabularx}{\linewidth}{|c|X|}\hline
\rowcolor{popblueL} \textbf{Étape} & \textbf{Travail à faire} \\ \hline
1 & Lire le sujet 3 fois ; souligner les mots-clés et les définir. \\ \hline
2 & Identifier le type de sujet et formuler la problématique (une question). \\ \hline
3 & Noter toutes les idées (brainstorming), puis trier et regrouper en 2 ou 3 axes. \\ \hline
4 & Rédiger le plan détaillé : introduction (amorce, sujet, problématique, annonce du plan), parties avec arguments et exemples, conclusion (bilan + ouverture). \\ \hline
5 & Vérifier la cohérence : chaque partie répond à la problématique ; substituts variés pour éviter les répétitions. \\ \hline
\end{tabularx}
\end{center}

\textbf{Types de plans à connaître}
\begin{itemize}
\item \cle{Plan dialectique} : thèse -- antithèse -- synthèse (sujets de discussion).
\item \cle{Plan thématique} : plusieurs aspects d'une même question (sujets de réflexion).
\item \cle{Plan analytique} : constat -- causes -- conséquences/solutions.
\end{itemize}

\textbf{Texte ouvroir} : \emph{Le Vieux nègre et la médaille} de Ferdinand Oyono : repérer le titre, l'auteur, l'éditeur, la quatrième de couverture (paratextes) avant la lecture méthodique ; savoir rendre compte de la lecture (résumé, personnages, thèmes : colonialisme, illusion, humiliation).
\end{bilan}

\begin{aretenir}[Checklist avant le Bac]
\begin{itemize}
\item $\square$ Je sais repérer et définir les mots-clés d'un sujet.
\item $\square$ Je sais distinguer sujet de réflexion, de discussion et de démonstration.
\item $\square$ Je formule une problématique claire en une seule question.
\item $\square$ Je fais un brainstorming puis je trie et regroupe mes idées en 2 ou 3 axes.
\item $\square$ Je choisis le plan adapté (dialectique, thématique, analytique).
\item $\square$ Mon plan détaillé contient arguments et exemples pour chaque partie.
\item $\square$ Mon introduction comporte amorce, présentation du sujet, problématique et annonce du plan.
\item $\square$ Ma conclusion fait le bilan et propose une ouverture.
\item $\square$ J'évite les répétitions grâce aux substituts (pronoms, synonymes).
\item $\square$ J'emploie correctement sigles et emprunts, sans abus.
\item $\square$ Je connais les paratextes et les grandes lignes du \emph{Vieux nègre et la médaille}.
\item $\square$ Je relis pour vérifier la cohérence et corriger la langue.
\end{itemize}
\end{aretenir}

\findecours

\end{document}
